\documentclass[11pt]{article}
\usepackage[T1]{fontenc}
\usepackage{lmodern,amsmath,amssymb,amsthm,mathtools}
\usepackage[margin=1in]{geometry}
\usepackage{booktabs,longtable,array,enumitem,xcolor}
\usepackage[colorlinks=true,linkcolor=blue,urlcolor=blue]{hyperref}
\newtheorem{theorem}{Theorem}[section]
\newtheorem{proposition}[theorem]{Proposition}
\newtheorem{lemma}[theorem]{Lemma}
\newtheorem{corollary}[theorem]{Corollary}
\theoremstyle{definition}
\newtheorem{definition}[theorem]{Definition}
\theoremstyle{remark}
\newtheorem{remark}[theorem]{Remark}
\newcommand{\R}{\mathbb R}
\newcommand{\HH}{\mathbb H}
\newcommand{\ip}[2]{\langle #1,#2\rangle}
\newcommand{\norm}[1]{\lVert #1\rVert}
\newcommand{\src}[1]{\textnormal{\textsf{[#1]}}}
\newcommand{\ph}{\varphi}
\newcommand{\pending}[1]{\par\smallskip\noindent
  {\color{red!65!black}\textbf{Internal proof transcription outstanding.} #1}\par\smallskip}
\title{Farthest vertices and dimension bounds for compact\\
hyperbolic Coxeter polytopes\\[4pt]
\large Consolidation draft for the existing general compact paper}
\author{CoxeterPolytopes project}
\date{24 October 2026 --- draft pending mathematical review}
\begin{document}
\maketitle
\begin{center}
\fbox{\begin{minipage}{.92\linewidth}
This is a replacement draft within the existing paper family, not a seventh
paper or a new accepted edition. The six canonical TeX/PDF pairs are
unchanged. The project has accepted an internally reviewed written proof of
the bound \(17\); this draft has not been compiled or independently reviewed.
It translates the final whole-\(D_{18}\) dotted-head argument in full, but
earlier internal inventories and some component exclusions still have
precisely identified project proof homes rather than complete proofs here.
Those passages must be incorporated into the appendices before this can be
called a self-contained manuscript. An internal proof awaiting transcription
is not an external mathematical axiom.
\end{minipage}}
\end{center}
\begin{abstract}
We study a compact hyperbolic Coxeter polytope from one vertex farthest
from a fixed interior point. Its incident spherical diagram gives an
orthogonal decomposition of a single Euclidean vector of norm less than
one. Actual opposite walls, and points obtained from positive combinations
of actual walls, impose incompatible lower bounds on that vector.
The method has two branches: concave interpolation for \(A\)-chains, and
prefix interpolation followed by terminal-wall analysis for \(B\)- and
\(D\)-components. The accepted source package excludes every dimension
at least \(18\). We organize its argument around common lemmas and one
final composition, rather than successive numerical improvements.
The last obstruction, a whole \(D_{18}\) diagram, is treated by a uniform
dotted-head pressure, legally funded reference vectors, the two genuine
signed fork directions, and a local calculation at an adjacent
\(A_{18}\) vertex. Arithmeticity is nowhere assumed.
\end{abstract}

\section{The theorem and the organization of its proof}
\label{sec:organization}
The theorem targeted by this revision is the following statement.
\begin{theorem}[Accepted written theorem; consolidated proof incomplete]
\label{thm:main}
Every compact real hyperbolic Coxeter polytope has dimension at most \(17\).
\end{theorem}
Here compact means that every vertex is finite. The theorem concerns
arbitrary Coxeter reflection groups; there is no ground field, lattice,
integrality condition, or arithmetic hypothesis. It asserts an upper
bound, not sharpness or the existence of an example in dimension \(17\).
The acceptance record is
\texttt{results/compact\_dimension17\_global\_written\_acceptance\_20261024.json}.
The source ledger accompanying this draft identifies its mathematical
inputs, their exact scopes, and the remaining transcription work.

The central object is one vector \(x\), obtained from an original
farthest vertex \(p\). Energies of orthogonal components of its incident
diagram can be added because they are projections of this same vector.
Pressures arising from different auxiliary points cannot in general
be optimized separately and then added. Instead, we combine them in
one reference vector and retain its complete squared norm.

The case split occurs at the finite irreducible components of the
diagram at \(p\). Small components pay uniform energy debits.
A large \(A\)-component is handled in its zero-sum chain coordinates.
A large \(B\)- or \(D\)-component has a long common prefix; its terminal
nodes then require separate analysis. Section~\ref{sec:forests} records a
fresh complete cover at dimension \(18\), and the separate
higher-dimensional cover. A \emph{cover} always means necessary
possibilities, not realized examples. Section~\ref{sec:b18} explains the
whole-\(B_{18}\) terminal partition. Sections~\ref{sec:d18-setup}--%
\ref{sec:composition} give the whole-\(D_{18}\) closure.

The numerical theorem uses exactly two global boundaries:
every dimension at least \(19\) has already been excluded, and the
same selected farthest vertex in dimension \(18\) has only a whole
\(D_{18}\) diagram left. Neither boundary is obtained by taking a face
or deleting a node from a higher-dimensional list. We do not assume
that a face is Coxeter, or that the property under investigation is
inherited by faces. The five remaining exact-dimension-\(17\) research
families are outside this manuscript.

\section{One original vector and actual geometric points}
\label{sec:geometry}
Use the hyperboloid model with Lorentz signature \((-+\cdots+)\).
Every wall has an outward unit spacelike normal \(n\), and the polytope
is given by \(\ip{r}{n}\leq0\). Distinct actual wall normals have
nonpositive products. Fix a strict interior point \(z\), choose an
actual vertex \(p\) farthest from \(z\), and write
\begin{equation}
\label{eq:original}
 z=E(p-x),\qquad E=-\ip{z}{p}>0,\qquad
 x\in p^\perp,\qquad \norm{x}^{2}=1-E^{-2}<1.
\end{equation}
The space \(p^\perp\) is Euclidean. The incident normals at \(p\) form
a positive definite spherical diagram; compact acute simplicity
implies that there are exactly \(n\) of them in dimension \(n\).
We retain these original normals and this original choice of
\(p,z,x,E\) throughout an energy argument.

The function \(r\mapsto-\ip{z}{r}\) takes its maximum on a compact
hyperbolic polytope at a vertex. Its restriction to a unit speed
geodesic has positive second derivative equal to the function itself,
so a maximum on a segment occurs at an endpoint; successively
passing to proper faces proves the assertion. Consequently
\begin{equation}
\label{eq:allpoint}
 -\ip{z}{r}\leq E\quad\text{for every actual point }r\in P.
\end{equation}
At an adjacent vertex \(q\), write \(z=E_q(q-y)\).
Equation~\eqref{eq:allpoint} gives a comparison factor \(E/E_q\),
not a claim that \(q\) is farthest.

\subsection{Positive inverse matrices and feasible projections}
\begin{lemma}[Inverse positivity]
\label{lem:inverse}
If a positive definite symmetric matrix \(B\) has nonpositive
off-diagonal entries, then \(B^{-1}\) is entrywise nonnegative.
\end{lemma}
\begin{proof}
Choose \(s\) greater than the largest eigenvalue of \(B\) and all its
diagonal entries. Then \(M=sI-B\) is entrywise nonnegative and has
spectral radius \(s-\lambda_{\min}(B)<s\). Thus
\(B^{-1}=s^{-1}\sum_{k\geq0}(M/s)^k\) is entrywise nonnegative.
\end{proof}
\begin{lemma}[An actual positive projection]
\label{lem:projection}
Let \(n_j\), \(j\in J\), be actual wall normals with positive
definite Gram matrix \(B_J\), and let \(z\) be strict interior.
Put
\[
 q_0=z-\sum_{i,j\in J}(B_J^{-1})_{ij}\ip{z}{n_j}n_i.
\]
Then \(q_0\) is future timelike, orthogonal to all selected walls,
and has negative product with every unselected actual wall.
Its normalization is a point of the original polytope.
\end{lemma}
\begin{proof}
The added coefficients are nonnegative by Lemma~\ref{lem:inverse}.
For an unselected wall every added product is nonpositive and the
original product is negative; the selected products vanish.
Moreover
\[
 q_0^2=z^2-
   (\ip{z}{n_i})_{i\in J}^{T}B_J^{-1}
   (\ip{z}{n_i})_{i\in J}<0.
\]
Since \(\ip{z}{q_0}=q_0^2<0\), it has future orientation.
All original wall inequalities survive normalization.
If fewer than \(n\) walls are selected this is a point on a face,
not an asserted full-dimensional vertex.
\end{proof}
The same sign check applies to positive combinations of actual
normals below: selected products are calculated explicitly; every
unknown original wall has nonpositive product with every summand.
A nonzero feasible future null vector would give a point of the
ideal boundary of \(P\), contradicting compactness. A positive
affine Coxeter kernel on actual walls is an obstruction after its
future sign and all remaining wall products have been checked.

\begin{lemma}[Future-point pressure]
\label{lem:pointpressure}
Suppose \(\xi=A p-U\) is future timelike, \(\xi^2=-k<0\),
and its normalization is feasible against every actual wall.
Then the same original vector satisfies
\[
 \ip{x}{U}\geq A-\sqrt{k}.
\]
At a vertex with all-point comparison factor \(\kappa\), the
corresponding inequality is \(\ip{y}{U}\geq A-\kappa\sqrt{k}\).
\end{lemma}
\begin{proof}
Apply \eqref{eq:allpoint} to \(\xi/\sqrt{k}\):
\(-\ip{z}{\xi/\sqrt{k}}=E(A-\ip{x}{U})/\sqrt{k}\leq E\).
The modified comparison is the identical calculation.
\end{proof}
If \(U\ne0\), its support lies in actual orthogonal components
and \(\norm{U}^2=A^2-k\); Cauchy therefore also gives
\begin{equation}\label{eq:pointenergy}
 \norm{x_{\operatorname{supp}U}}^2
 \geq\frac{A-\sqrt{k}}{A+\sqrt{k}}.
\end{equation}
Normalize a future vector, not a spacelike correction.
For a cap \(c=A p-v\), clear a positive definite old block \(J\):
\[
 w=\sum_{i,j\in J}(B_J^{-1})_{ij}\beta_j n_i,\quad
 \beta_j=-\ip{c}{n_j}\geq0,\quad
 \eta=c+w,\quad \delta=\eta^2=1-\beta^TB_J^{-1}\beta.
\]
If \(\delta<0\), \(\eta\) is the future vector to normalize.
If \(\delta>0\), the actual finite projection instead uses
\[
 p+\frac{A}{\delta}\eta,\qquad
 \left(p+\frac{A}{\delta}\eta\right)^2=-1-\frac{A^2}{\delta}.
\]
When its selected products are nonpositive, its positive
coefficients check all unselected walls. The pressures are
\begin{equation}
\label{eq:corepressure}
 \ip{x}{v-w}\geq A-\sqrt{-\delta},\qquad
 \ip{x}{v-w}\geq
 \frac{A}{1+\sqrt{\delta/(A^2+\delta)}}.
\end{equation}
Old normals lie in \(p^\perp\), so \(v-w\) denotes the spatial
part of the corresponding displayed Lorentz vector.

\subsection{Opposite walls and common components}
Let \(B\) be the original incident Gram matrix, \(C=B^{-1}\), and
\[
 \omega_i=\sum_j C_{ij}n_j,\qquad
 d_i=C_{ii},\qquad u_i=\omega_i/\sqrt{d_i}.
\]
An actual opposite wall has normal
\[
 c=A p-v,\quad A>0,\quad
 v=\sum_i\beta_i\omega_i,\quad \beta_i=-\ip{c}{n_i}\geq0,\quad
 A^2=\beta^TC\beta-1.
\]
Strict interiority gives
\begin{equation}\label{eq:actualinterior}
 \ip{x}{v}<A.
\end{equation}
Along the original edge in direction \(-u_i\), of length \(L_i\),
the original farthest comparison gives
\begin{equation}
\label{eq:halfheight}
 \ip{x}{u_i}\geq\tanh(L_i/2).
\end{equation}
Whole-component projection energies sum to \(\norm{x}^2\).

\begin{proposition}[Reciprocal shadows; internal geometric interface]
\label{prop:shadows}
For an actual cap put \(T_i=(C\beta)_i/(A\sqrt{d_i})\).
If \(T_i>1\), it is the opposite wall of edge \(i\), and
\(\tanh L_i=1/T_i\). If it is not that opposite wall, then
\(\tanh L_i\geq T_i\). The value \(T_i=1\) is incompatible
with a finite endpoint. For distinct actual opposite walls,
\begin{equation}\label{eq:pair-shadow}
 \tanh L_i\,\tanh L_j\geq
 K_{ij},\qquad K_{ij}=C_{ij}/\sqrt{d_i d_j}.
\end{equation}
Opposite walls of orthogonal original directions cannot coincide.
\end{proposition}
\begin{proof}
The Schur identity says that
\(C-C_{\cdot i}C_{i\cdot}/d_i\) vanishes on row and column \(i\);
elsewhere it is the inverse of the retained old Gram matrix.
It is entrywise nonnegative. Let \(d=Dp-\sum_j\gamma_j\omega_j\)
be the actual opposite of edge \(i\), distinct from \(c\).
Acuteness and the Schur identity give
\[
 AD\geq\beta^TC\gamma
 \geq(C\beta)_i(C\gamma)_i/d_i.
\]
The actual endpoint equation is
\(\tanh L_i=D\sqrt{d_i}/(C\gamma)_i\), so
\(\tanh L_i\geq T_i\). If \(T_i>1\) this alternative is
impossible, forcing \(c=d\); its endpoint equation gives
\(\tanh L_i=1/T_i\).
If \(T_i=1\), the distinct alternative is impossible and the
same-wall equation would give an ideal endpoint. This is
excluded by compactness.
For the pair inequality, first use the contacts of the actual
opposite of \(i\) to get \(T_j\geq K_{ij}/\tanh L_i\);
the distinct alternative at edge \(j\) yields the assertion.

Finally suppose the same wall closes orthogonal directions
\(i,j\). Clearing every other old normal adds positive
coefficients and gives
\[
 c_F=A p-\frac{A}{r_i}u_i-\frac{A}{r_j}u_j,\qquad
 r_i=\tanh L_i,\quad r_j=\tanh L_j.
\]
Its product with the strict interior point remains negative,
so \(\ip{x}{u_i}/r_i+\ip{x}{u_j}/r_j<1\).
But \eqref{eq:halfheight} makes this sum at least
\(1+(\tau_i^2+\tau_j^2)/2>1\),
where \(\tau_i=\tanh(L_i/2)\). This contradiction proves
the orthogonal assertion.
\end{proof}
The displayed pair inequality concerns full tangents, not
half tangents.
For an opposite wall, its retained foreign component can have
at most one contact: two would make a cycle in its actual
finite replacement. A positive foreign cost is
\(q=\beta_j^2C_{jj}\geq1/4\). In a nontrivial component
\(C_{jj}\geq4/3\), so the only cost below \(1/3\) is a
simple \(3\)-contact on an \(A_1\), of cost \(1/4\).
Moreover
\begin{equation}\label{eq:foreigncost}
 q<A^2.
\end{equation}
Otherwise the actual reciprocal test would close that foreign
direction too, violating the orthogonal assertion, with the
equality case giving a null endpoint.
Set \(F(t)=2t/(1+t^2)\). For unit directions of correlation \(k\)
in a component of energy at most \(e\), positive heights satisfy
\begin{equation}\label{eq:pairupper}
 \tau_i+\tau_j\leq\sqrt{2e(1+k)},\qquad
 F(\tau_i)F(\tau_j)\leq\frac{4s}{(1+s)^2},\quad
 s=\frac{e(1+k)}2.
\end{equation}
The second inequality follows from concavity of \(\log F\) on
\((0,1)\). If its upper bound is smaller than \(K_{ij}\),
Proposition~\ref{prop:shadows} forces a common actual opposite.
Its finite-white extension or positive affine kernel then gives
additional actual pressures.

\section{Reference vectors, individual prices, and missing floors}
\label{sec:references}
The elementary estimate behind every certificate is
\begin{equation}
\label{eq:square}
 \norm{x}^2\geq2\ip{x}{W}-\norm{W}^2.
\end{equation}
Every term in \(\ip{x}{W}\) must be justified.
Write \(w=\sum_i b_iP_i\), \(b_i\geq0\), with actual pressures
\(\ip{x}{P_i}>a_i\), and demand \(A_w=\sum_i b_i a_i\).
For a new actual pressure \(\ip{x}{Z}\geq h\), choose
\[
 Q=Z-\sum_i t_iP_i,\qquad W=w+\mu Q.
\]
Before using pressures check every \(b_i-\mu t_i\geq0\).
Then
\begin{equation}
\label{eq:payment}
 \ip{x}{W}>A_w+\mu(h-\sum_i t_i a_i),\quad
 \norm{W}^2=\norm{w}^2+2\mu\ip{w}{Q}+\mu^2\norm{Q}^2.
\end{equation}
The available price \(b_i\) is individual to anchor \(i\).
For multiple corrections, their combined debit at each anchor
and all mutual cross terms must be retained.

Put \(P_i=e_1+\cdots+e_i\), \(S_i=\ip{x}{P_i}\).
If selected indices satisfy \(S_i>f_i\), concave interpolation
through \((0,0)\) and the selected \((i,f_i)\) gives a reference
whose coordinates are successive slopes. Its \(P_i\) coefficients
are successive slope decreases, hence nonnegative; its demand
is its squared norm. Omitting one unknown internal floor replaces
two adjacent slopes by their average, losing half their squared
difference. For \(\sqrt t\) at \(t\geq5\), the loss is less than
\(1/3200\). Nonadjacent omissions have disjoint open intervals,
so losses add. A missing floor has coefficient zero.
References are Euclidean vectors, not asserted actual points:
actual points prove pressures, then positive coefficients combine
them in \eqref{eq:square}.

\section{Common component bounds and the precise branch point}
\label{sec:components}
\subsection{Ordinary components and exceptional actual defects}
All bounds below concern the same original projection, with
arbitrary finite spectators.
\begin{center}\small
\begin{tabular}{ll@{\qquad}ll}
\toprule
Component&Energy greater than&Component&Energy greater than\\ \midrule
ordinary nonshort \(A_1\)&\(1/10\)&ordinary \(A_2\)&\(319/2000\)\\
ordinary \(I_2(m),m\geq4\)&\(1/7\)&\(A_3\)&\(1/4\)\\
\(B_3,H_3\)&\(2/7\)&\(A_4\)&\(2/7\)\\
\(B_4,D_4,F_4\)&\(1/3\)&\(H_4\)&\(11/36\)\\
\(A_5\)&\(3/8\)&\(A_6\)&\(3/7\)\\
\(B_5,B_6\)&\(3/7\)&\(D_5,D_6\)&\(5/14,\ 3/7\)\\
\(B_7,D_7\)&\(1/2\)&\(B_8,D_8\)&\(4/7\)\\
\(E_6,E_7,E_8\)&\(9/20,\ 5/9,\ 2/3\)&&\\ \bottomrule
\end{tabular}
\end{center}
An actual isolated dotted short \(A_1\) is denoted \(I\).
A second short-wall profile \(J\) uses two old \(A_1\) factors
whose joint energy exceeds \(.34\).
The two typed rank-two exceptions, each of energy \(>.09\), are
\[
 N_7=(\text{old }A_2,\ \text{new isolated native }7),\qquad
 O_7=(\text{old }I_2(7),\ \text{new isolated native }3).
\]
A bare \(I_2(7)\) is not automatically \(O_7\).
There is at most one short \(I/J\) profile or one typed exception;
they do not coexist. Acuteness of their actual isolated walls
is essential: two on orthogonal supports would have product
\(-\sqrt{h_1h_2}\) strictly between \(-1/2\) and \(0\).
The full endpoints, \(J\) case and ordinary rates are internal
inputs \src{Ord}, \src{F19-local}, not field classifications.

Useful joint bounds, counting \(I\) only once, are
\begin{equation}
\label{eq:joint}
 I+A_1>1/5,\quad I+2A_1>3/10,\quad
 I+T_2>21/100,\quad I+T_3>2/7.
\end{equation}
Here \(T_j\) is ordinary irreducible of rank \(j\).
Both typed exceptions force another \(T_2\) to cost \(>1/5\).
Another \(A_1\) costs \(>1/5\) beside \(N_7\) and \(>2/15\)
beside \(O_7\). For the latter, isolated height
\(h=\tfrac14\csc^2(\pi/7)-1<1/3\) gives squared full tangent
greater than \(3/7\). Squared half tangent exceeds \(2/15\),
since \(4(2/15)/(1+2/15)^2=120/289<3/7\).
This is a real geometric argument.
The complete contact and defect proofs, including the actual-wall
height constraints behind \eqref{eq:joint}, are given in
Appendix~\ref{app:ordinary}. The uniform actual \(B/D\) floor
is proved in Appendix~\ref{app:BDfloor}.

\subsection{The \(A\)-chain branch}
For \(A_r\), use the zero-sum subspace of \(\R^{r+1}\) with roots
\((e_i-e_{i+1})/\sqrt2\). The inverse Gram matrix is
\[
 C_{ij}=\frac{2\min(i,j)(r+1-\max(i,j))}{r+1}.
\]
A finite-white classification of actual opposite walls gives
strong dual floors except for a short list of literal helper caps.
Adjacent weak helpers are excluded using their actual mutual
product and a positive two-cap projection. Concave interpolation
uses only strong floors; helper pressures are paid at surviving
anchors. The small end variance yields stronger component bounds:
\[
\begin{array}{c|ccccc}
r&r\geq7&r\geq8&r\geq13&16&17\\ \hline
e(A_r)>&1/2&4/7&5/7&199/250&6/7
\end{array}
\qquad e(A_{18})>91/100.
\]
Later local refinements give \(e(A_{12})>13/20\) and
\(e(A_{15})>413/500\), uniformly with arbitrary spectators.
Only the first nonarithmetic component theorem of the mixed-title
\src{A13} is used, not its arithmetic packing theorem.
The shared actual feet, common-pair inventories, uniform
\(1/2,4/7,5/7\) bounds, and the \(N=17,18,19\) central-window
and paid-reference arguments are proved in Appendix~\ref{app:Achain}.
The newer \(A_{16}>.86\) proof in particular implies the older
\(.796\) bound without requiring its separate historical reference.
The fresh smaller-rank proofs are in Sections~\ref{app:A12}
and~\ref{app:A15}. Their full native-cosine inventories and
individual actual-helper payments are included. Later exact-\(17\)
research remains excluded.

\subsection{The common \(B/D\) prefix branch}
For \(B_r\), use long roots \((e_i-e_{i+1})/\sqrt2\), \(i<r\),
and short root \(e_r\). For \(D_r\), use the long chain and the
fork root \((e_{r-1}+e_r)/\sqrt2\). On a long initial segment,
\(\omega_i=\sqrt2P_i\), \(C_{ij}=2\min(i,j)\).
In type \(D\), \(x_1>\cdots>x_{r-1}>|x_r|\); the last
coordinate is signed and the forks are distinct actual duals.
\begin{proposition}[Generic prefix theorem]
\label{prop:prefix}
Let \(6\leq M\leq16\). Suppose an actual component is \(B_r\),
\(r\geq M+5\), or \(D_r\), \(r\geq M+6\). Put
\[
 f_i=\sqrt i-\frac{61}{100\sqrt2},\quad
 E_M=\frac{f_5^2}{5}+\sum_{i=5}^{M-1}(\sqrt{i+1}-\sqrt i)^2,
 \quad H_M=\left\lfloor\frac{M-5}{2}\right\rfloor.
\]
With arbitrary finite spectators,
\(\sum_{j=1}^{M+1}x_j^2>E_M-H_M/3200\).
\end{proposition}
\begin{proof}
The exact-rank inventory \src{Prefix} proves that a weak site
\(5\leq i\leq M+1\) has a pure cap
\[
 v=\gamma\omega_i+\omega_1/2,\quad
 A^2=2i\gamma^2+2\gamma-\tfrac12,\quad
 101/125<\gamma<6/5,
\]
with own label \(5\), finite label at least \(7\), or dotted own
contact, and no spectator contact. Actual weak sites are
nonadjacent. This is proved at these ranks, not transferred from
a whole-component exclusion.
If \(5\) is strong, choose a strong \(K=M\) or \(M+1\).
Interpolate through actual strong knots. Each missing internal
floor loses \(<1/3200\); at \(K=M+1\) the extra increment
\((\sqrt{M+1}-\sqrt M)^2>1/(4M+2)\geq1/66\)
pays the possible extra omission.

If \(5\) is weak, the first-six-root actual point table of
\src{Prefix} gives \(x_1>.5965\), and \(6\) is strong, with
\(S_6>2.01815\). Interpolate through \((1,.5965),(6,2.01815)\)
and a later strong endpoint. The slopes decrease:
\(.5965,.28433\), then \(1/(\sqrt K+\sqrt6)<1/5\).
Its energy exceeds
\[
 L_M=.7600299945+\frac{(\sqrt M-\sqrt6)^2}{M-6},
\]
the last term being zero at \(M=6\).
Now \(E_6<.7\). The \(k=M-6\) increments from \(6\) onward lie
in \((1/8,1/5)\), so their variance is at most
\(k(1/5-1/8)^2/4\leq9/640\).
Thus \(L_M-E_M>.76-.7-9/640>0\), directly also at \(M=6\).
This proves the weak-\(5\) branch without first assuming strong \(5\).
\end{proof}
Appendix~\ref{app:prefixinventory} supplies the complete
finite-white inventory, zero-own repair, actual two-cap
nonadjacency, all four weak-\(5\) pressures and continuous
seals, with every spectator retained. Its original rank
inequalities are part of the theorem. A short-tail continuation
at an excluded rank still requires its own new local proof.

\section{Complete forest covers and their consumers}
\label{sec:forests}
\subsection{A fresh cover in exact dimension \(18\)}
Call \(A_r\), \(r\geq7\), or \(B_r,D_r\), \(r\geq9\),
dominant. Each costs \(>1/2\), so at most one is possible.
With none the ordinary rank-over-\(14\) baseline, even allowing
the unique defect, exceeds \(1\). Exactly one is therefore present.
The fresh total-rank-\(18\) necessary cover is:
\begin{center}\small
\begin{tabular}{p{.14\linewidth}p{.75\linewidth}}
\toprule
Dominant&All retained spectator alternatives\\ \midrule
\(A_{12}\)&\(I+(A_5\text{ or }D_5)\), or \(N_7/O_7+T_4\)\\
\(A_{15}\)&ordinary \(A_3\); \(T_2+A_1\); \(I+T_2\); or \(O_7+A_1\)\\
\(A_{16}\)&ordinary \(T_2\); two ordinary \(A_1\); \(I+A_1\); or \(N_7/O_7\)\\
\(A_{17}\)&one ordinary nonshort \(A_1\), or \(I\)\\
\(A_{18}\)&none\\
\(B_{16}\)&\(N_7/O_7\)\\
\(B_{17}\)&\(I\)\\
\(B_{18}\)&none\\
\(D_{16}\)&\(N_7/O_7\)\\
\(D_{17}\)&one ordinary nonshort \(A_1\), or \(I\)\\
\(D_{18}\)&none\\ \bottomrule
\end{tabular}
\end{center}
Here ordinary \(T_2=I_2(m)\), \(m\geq3\), and
\(T_4\in\{A_4,B_4,D_4,F_4,H_4\}\).
\(H_4\) and native \(5\) are not discarded by field arguments.

For completeness, the accounting steps are as follows.
\(A_7\) leaves too many ordinary ranks;
\(A_8,\ldots,A_{11}\) are excluded by \(4/7\) plus
the rank-over-\(14\) baseline.
Beside \(A_{12}\), ordinary or \(J\) rank-six spectators cost
\(>3/7\). With \(I\), \eqref{eq:joint} excludes all rank-five
partitions except \(A_5,D_5\). With a typed exception, partition
rank \(4\): two \(T_2\)'s cost \(>.4\), \(T_3+A_1\) costs
\(>1/4+2/15\), and the other reducible partitions are stronger.
Only irreducible \(T_4\) remains. For \(A_{13},A_{14}\),
the \(5/7\) bound and \eqref{eq:joint} exclude all spectators.
At \(A_{15}\), exclude \(B_3,H_3\), three ordinary \(A_1\)'s,
\(I+2A_1\), \(J+A_1\), and \(N_7+A_1\).
At \(A_{16}\), \(.796\) excludes \(J\) but retains precisely
the listed rank-two alternatives. At \(A_{17}\) retain the
remaining \(A_1\).

For \(B_r,D_r\), the \(9/14\) bound forces spectator rank
at most \(5\), hence \(r\geq13\).
Proposition~\ref{prop:prefix}, within its rank domain, gives
\[
\begin{array}{c|rrrr}
r&13&14&15&16\\ \hline
e(B_r),e(D_r)>&.7349975&.7683275&.797425&.823735.
\end{array}
\]
At \(13,14,15\), every ordinary, \(I/J\), and typed remainder
exceeds the residual budget; \(I+T_2>.21\) closes rank \(15\).
At \(16\), two ordinary \(A_1\)'s and \(I+A_1\) are excluded;
ordinary \(T_2\) and both typed defects initially remain.
The fresh \(B_3,D_4,D_3\) layer checks of \src{F18},
Sections 5--6, extend helper coverage to \(13\), giving
energy \(>.8689525\); ordinary \(T_2>1/7\) closes that case.
The fresh \(B_{17}\) layers reach a strong \(14/15\) endpoint
and energy \(>.90715\), excluding a nonshort ordinary \(A_1\).
This reuses the interpolation proof after new local coverage,
not Proposition~\ref{prop:prefix} outside its legal ranks.

The first rows close with small shared consumers:
\[
 A_{12}+I+(A_5/D_5)>13/20+5/14=141/140,\quad
 A_{12}+N_7/O_7+T_4>13/20+.09+2/7=359/350.
\]
The improved \(413/500\) bound closes every \(A_{15}\) family:
ordinary total rank three costs \(>17/70\), \(I+T_2>.21\),
and beside \(O_7+A_1\) the total exceeds \(787/750\).
Thus the older separate ordinary-rank-three consumer is unnecessary.
The other internal proof homes are:
\begin{center}\small
\begin{tabular}{lp{.69\linewidth}}
\toprule
Row&Excluding project argument\\ \midrule
\(A_{16}\)&\src{A16-new}, ordinary and \(I+A_1\);
 \src{A16-typed}, both actual typed defects\\
\(A_{17}\)&\src{A17-ordinary}; corrected \src{A17-I}\\
whole \(A_{18}\)&\src{WholeA18}, full two-end native-gap argument\\
\(B_{16}\)&\src{B16-typed}, both actual typed defects\\
\(B_{17}\)&\src{B17-I}, dotted-head and retained-tail analysis\\
\(D_{16}\)&\src{D16-typed}, both actual typed defects\\
\(D_{17}\)&\src{D17-ordinary}; \src{D17-I1}, corrected \src{D17-I2}\\
whole \(B_{18}\)&Section~\ref{sec:b18}\\
whole \(D_{18}\)&Sections~\ref{sec:d18-setup}--\ref{sec:composition}\\
\bottomrule
\end{tabular}
\end{center}
Appendix~\ref{app:exact18layers} gives \src{F18}'s fresh local
layers, including the distinct signed \(D_4/D_3\) forks and the
actual \(A_{17}\) neighbor for \(B_{17}\).
\pending{The other exact-\(18\) row exclusions in the table and
whole-\(A_{18}\) still require their complete proof bodies.
The necessary cover and all its consumers are separate obligations.
Only identified real geometric sections of mixed-title sources
may enter; no field, lattice or integral dotted-product premise is allowed.}

\subsection{The separate complete higher-dimensional tail}
\label{subsec:higher}
The exact-\(18\) cover does not exclude higher dimensions.
The earlier internal package establishes the separate theorem:
\begin{theorem}[Internal higher-dimensional exclusion]
\label{thm:tail}
There is no compact hyperbolic Coxeter polytope of dimension at least \(19\).
\end{theorem}
\pending{The accepted project proof of Theorem~\ref{thm:tail} is not
fully transcribed here. Appendix E must contain the finite covers
and all consumers, using the published ceiling only at dimensions
at least \(30\). The section-level route below makes the obligation
finite and excludes an arithmetic cycle.}
At exact \(19\), \src{F19} has nine necessary groups:
\[
\begin{array}{c|l}
A_{16}&A_3,\ T_2+A_1,\ I+T_2,\ N_7/O_7+A_1\\
A_{17}&T_2,\ N_7/O_7\\
A_{18}&A_1\text{ ordinary or }I\\
A_{19}&\text{whole}\\ B_{18}&I\\ B_{19}&\text{whole}\\
D_{17}&N_7/O_7\\ D_{18}&I\\ D_{19}&\text{whole}.
\end{array}
\]
Source \src{G18-body}, Sections 2--4, excludes them all.
Its \(A_{19}\) branch separates zero, one and two weak ends.
Its \(B_{19}\) branch uses the corrected short-three successor.
Its \(D_{19}\) branch covers heads \(5,7,8,9\), all finite
heads at least \(10\), and the entire dotted interval.
These are actual whole-component arguments using one original
vector, not a shifted-rank cover.

The remaining numerical composition has this acyclic route:
\[
\begin{gathered}
\text{published }29\ \longrightarrow\
\text{canonical geometric exclusions }24\text{--}29\\
\longrightarrow\text{exact }23\longrightarrow\text{exact }22
\longrightarrow\text{exact }21\longrightarrow\text{exact }20
\longrightarrow\text{exact }19.
\end{gathered}
\]
Each arrow denotes a global logical composition, not geometric
dimension heredity; each cover is derived at its own rank.
The exact-\(22\) frontier has whole \(B_{22},D_{22}\) and
\(B_{21}+I,D_{21}+I\). Exact \(20\) reduces to whole
\(A_{20},B_{20},D_{20}\), each with its own actual-point proof.
Exact-\(21\) and exact-\(23\) proofs have ledger homes and
must be reproduced, not inferred from this overview.
Only general Sections 2--3 of \src{G19-mixed} and Sections
2--4 of \src{G20-mixed} occur on this route.
Their arithmetic chapters and field censuses are excluded.
The bound \(17\) does not depend on arithmetic \(15\), an
arithmetic exact-\(17\) theorem, or an exact-\(17\) field census.

\section{The whole-\(B_{18}\) terminal partition}
\label{sec:b18}
First classify the actual opposite of old node \(17\) without
assuming a strong head. The complete inventory \src{B18-T} leaves:
\begin{enumerate}[label=(\roman*)]
\item a pure far-three helper with finite own \(5\), finite own
 at least \(7\), or dotted own \(1<g<4/3\);
\item near-three with own \(6\);
\item no left contact, with short white label \(3,4,5,6,7,8\).
\end{enumerate}
All other isolated, left, bridge and affine/null roles are
excluded in that same inventory.
Sources \src{B18-F}, \src{B18-D} close the finite and full
dotted far cases. Source \src{B18-N} closes near-three/own-six
through all seven actual short-wall roles.
In the last branch, \src{B18-R} closes distinct \(17/18\)
caps and leaves seven common-cap roles:
\[
\begin{array}{c|c}
\text{own }0&\text{short }5,6,7,8\\
\text{own }3&\text{short }3,4,5.
\end{array}
\]
Source \src{B18-C} closes own \(0\)/short \(6,7,8\) and
own \(3\)/short \(4,5\). Source \src{B18-Z} closes
own \(0\)/short \(5\). The last own \(3\)/short \(3\)
role is excluded at its actual adjacent whole-\(A_{18}\)
vertex, \src{A-local}, with comparison \(11/8\).
That vertex is never assumed farthest.

This exhausts the terminal list. The complete conjunction
\src{B18-compose} and its root gate \src{V18} also bind the
earlier exact-\(18\) row closures. Thus the same selected
original vertex in dimension \(18\) has only whole \(D_{18}\)
left, without any claim about all other vertices.
\pending{The terminal appendix must transcribe the actual finite-white inventory
and each terminal exclusion, all unknown-wall checks, and the
corrected asymmetric price \(.13/10=.013\) in \src{B18-price}.
This section gives the complete composition and branch point,
but its constituent geometric proofs remain internal project inputs.}

\section{The last case: an actual whole-\(D_{18}\) vertex}
\label{sec:d18-setup}

We now give the final dotted-head argument of \src{D18-final}
in full. Its earlier internal inventory interfaces are stated
separately from the new estimates.
Keep the original \(p,z,x\) and use actual old roots
\[
 n_i=(e_i-e_{i+1})/\sqrt2\quad(1\leq i\leq17),\qquad
 n_{18}=(e_{17}+e_{18})/\sqrt2.
\]
Write
\[
 P_i=e_1+\cdots+e_i,\quad
 \Omega_i=\sqrt2P_i\ (i\leq16),\quad
 \Omega_{17}=(P_{17}-e_{18})/\sqrt2,\quad
 \Omega_{18}=(P_{17}+e_{18})/\sqrt2.
\]
In particular \(x_{18}\) is signed and
\(x_1>\cdots>x_{17}>|x_{18}|\). In this part only, set
\begin{equation}
\label{eq:d18floors}
 S_i=\ip{x}{P_i},\qquad f_i=\sqrt i-\frac{.58}{\sqrt2}.
\end{equation}
All terminating decimals in this paper denote exact rational numbers.

The head is the actual opposite wall of the long-end edge \(1\).
The complete isolation theorem \src{H} says that its normal is
\[
 H=A_Hp-\eta\Omega_1,\qquad A_H^2=2\eta^2-1>0.
\]
It also supplies the complete own-contact whitelist.
The strong terminal floor \(S_{16}>f_{16}\), its legal reference
append, and the complete fork-role list are supplied by \src{L}.
Finite heads are excluded by the complete conjunction \src{FH}.
We treat every dotted head \(\eta>1\) without a parameter cutoff.

\pending{The whole-\(D_{18}\) appendix must contain the complete head-isolation and
own-contact proof \src{H}, the full changed-threshold window
\src{D}, strong \(16\) and the endpoint-\(16\) append of
\src{L}, the finite-head conjunction \src{FH}, and the
every-head fork exclusions \src{II}, \src{IN}, \src{FF},
\src{T}. These are earlier internal geometric lemmas. In
particular strong \(5\) at the new threshold is not an inherited
unconditional fact.}

\subsection{A future point valid for every dotted head}

Put \(u=\eta^{-2}\in(0,1)\), and consider
\[
 \Xi=H+n_1/\eta.
\]
This is a positive combination of two actual normals.
Its products with \(H,n_1\) are respectively \(0\) and
\(-\eta+1/\eta<0\), and \(\Xi^2=-(1-u)<0\).
Its time is \(A_H>0\). Every other actual wall has nonpositive
product with each summand. Hence its normalization is an actual
future point of the original polytope, with no unknown wall
omitted. Lemma~\ref{lem:pointpressure} gives
\begin{equation}
\label{eq:headpoint}
 \ip{x}{Z_u}\geq h(u),\qquad
 Z_u=\frac{(2-u)e_1+u e_2}{2},\quad
 h(u)=\frac{\sqrt{2-u}-\sqrt{u(1-u)}}{\sqrt2}.
\end{equation}
Indeed the spatial part of \(\Xi\), divided by \(\sqrt2\eta\),
is \(Z_u\). We will use
\[
 \sum_j(Z_u)_j=1,\qquad R(u)=\norm{Z_u}^2=1-u+u^2/2.
\]
Equation~\eqref{eq:headpoint} is a point pressure, not an
energy lower bound, and is valid on the entire open interval.

\subsection{The four original reference cases}
\label{subsec:fourrefs}

The fresh \(.58\) inventory \src{D} has the following complete
initial alternatives: strong \(5\); a weak pure far-left-three
helper at \(5\); or an isolated dotted weak \(5\) with a chosen
strong \(J=6\) or \(7\). Thus there are four reference rows.
Here is why the initial classification has this precise scope.
Left own \(4\), including own \(4\)/white \(3\), gives a positive
actual affine-\(B\) null on the cap and old chain \(5,\ldots,18\);
the extra left contact preserves the remaining wall signs.
The newly possible near-three/own-five at \(5\) is excluded by
its actual right-\(H_4\) point and independent old strong
\(11/12\), before the new references are used. Other finite-white
and isolated finite roles are strong. A historical fifth paid
row concerned a native-five head and is outside \(\eta>1\).

We emphasize the \emph{bare} references for the three nonhelper
rows. They interpolate through the strong knots, are flat on
their first \(J\) coordinates, and have endpoint sum \(f_{16}\).
The complete contracts \src{B}, \src{L} are
\[
\begin{array}{c|c|c|c|c}
\text{initial case}&B_0&\text{first mean }b<
 &\text{unused }P_J\text{ price}>&J\\ \hline
\text{strong }5&.9565188&.3652&.151&5\\
\text{isolated dotted }5,\ \text{strong }6&.9373243&.3399&.1435&6\\
\text{isolated dotted }5,\ \text{strong }7&.9196357&.31938&.1366&7
\end{array}
\]
More explicitly, each reference \(w\) has
\[
 \ip{x}{w}>A_w,\quad 2A_w-\norm{w}^2>B_0,\quad
 \operatorname{supp}w\subset\{1,\ldots,16\},\quad
 \ip{w}{P_{16}}=f_{16},\quad \beta_{16}>.127.
\]
No head pressure has been paid. Their raw endpoint-\(15\)
lower bounds are \(.9403998,.9212053,.9035167\).
The legal endpoint-\(16\) append adds at least
\(.127^2-.00001=.016119\); if \(15\) is weak, its fictitious
knot is deleted before pressure. All late repairs preserve the
flat first block and its individual prices.

In the isolated-\(5\) rows there is also an unused actual pressure
\[
 U_5=P_5+\frac v2(e_6-e_5),\quad .4<v<1,\qquad
 \ip{x}{U_5}>1.8185.
\]
For
\[
 Q_H=Z_u-P_J/J,\qquad Q_5=U_5-(5/J)P_J
\]
the exact identities are
\begin{equation}
\label{eq:flatresiduals}
 \norm{Q_H}^2=R-1/J,\quad
 \ip{Q_H}{Q_5}=1-5/J,\quad
 \norm{Q_5}^2<
 \begin{cases}77/150,&J=6,\\194/175,&J=7.\end{cases}
\end{equation}
Both residuals have zero product with the bare flat reference.
The positive mutual cross in \eqref{eq:flatresiduals} will
not be discarded.

For the fourth row, the exact \emph{unheaded} weak-helper
reference interpolates from
\((1,a_0)\), \(a_0=.5455\), to \((6,f_6)\), then through
the later actual strong knots to \(16\), replacing each missing
knot by a chord. Its first six coordinates are
\begin{equation}
\label{eq:helperbare}
 w_1=a_0,\quad w_2=\cdots=w_6=b_0=(f_6-a_0)/5.
\end{equation}
The actual helper proves \(x_1>a_0\) and strong \(6\).
The raw contract of \src{B}, \src{E}, \src{L} is
\[
 \ip{x}{w}>A_w,\quad L_w:=2A_w-\norm{w}^2>.987759,\quad
 \ip{w}{P_{16}}=f_{16},\quad \beta_{16}>.127.
\]
Weak \(7\), later omissions, and the late \(15\) repair do not
alter \eqref{eq:helperbare}. Again no head pressure has been paid.
\pending{The four references and their exact construction are specified
here, but The whole-\(D_{18}\) appendix must also reproduce the rational increment and
omission tables proving their raw contracts, and the actual isolated-\(5\)
pressure \(U_5\), from \src{B}, \src{D}, \src{E}, \src{L}.
These are required internal reference lemmas, not script output or
unexplained lower-energy constants. Only the indicated local sections
of \src{B} and \src{E} are used, not their separate selected-branch
conclusions.}

\section{Uniform head payment and the two true fork heights}
\label{sec:uniformhead}

\subsection{The first three rows}
In the strong-\(5\) row choose \(\mu=.33\) on \(Q_H\).
Its debit \(.33/5=.066<.151\) fits the individual price.
Equations~\eqref{eq:payment}--\eqref{eq:flatresiduals} give
\[
 L_W>.9565188+.66(h-.3652)-.1089(R-1/5)
       =F(u):=.7372668+.66h(u)-.1089R(u).
\]
For isolated \(5\) with \(J=6\), use weights
\(\mu=.33,\nu=.10\) on \(Q_H,Q_5\).
The combined debit at \(J\) is \(.83/6<.1435\).
For \(J=7\) use \(\mu=.33,\nu=.11\), with debit
\(.88/7<.1366\). In each case all remaining old coefficients
are nonnegative \emph{before} the new pressures are invoked, and
\begin{align}
\label{eq:isolatedfunding}
 L_W>{}&B_0+.66(h-b)+2\nu(1.8185-5b)
       -.1089(R-1/J)\notag\\
 &-\nu^2\norm{Q_5}^2-.66\nu(1-5/J).
\end{align}
The constant terms before \(.66h-.1089R\) in the two rows are
\(22164209/30000000\) and \(51729823/70000000\), both
greater than \(.7372668\). Thus \(F\) is a common lower bound.
No optimized helper and head energies have been added.

We prove the bound uniformly on \(0<u<1\). Put
\(a=\sqrt{2-u}\), \(b=\sqrt{u(1-u)}\). Direct differentiation gives
\begin{equation}
\label{eq:hcurvature}
 h''(u)=\frac{1/(4b^3)-1/(4a^3)}{\sqrt2}
        >\frac7{4\sqrt2},
\end{equation}
because \(b\leq1/2\), \(a>1\). Hence
\(F''>.66\cdot7/(4\sqrt2)-.1089>.7\), using \(\sqrt2<1.415\).
At \(u_0=2/3\),
\[
 h(u_0)=\sqrt{2/3}-1/3,\quad R(u_0)=5/9,\quad
 F(u_0)>.9956502,\quad
 F'(u_0)=.2013-.0825\sqrt6,\quad |F'(u_0)|<.0008.
\]
The value bound uses
\(.81649^2=.6666559201<2/3\), and the derivative bound uses
\(2.4494<\sqrt6<2.4495\). Taylor's inequality and completion
of the square on the segment between \(u_0\) and \(u\) give
\begin{equation}
\label{eq:Fseal}
 F(u)>.9956502-\frac{.0008^2}{1.4}>.99564
 \qquad(0<u<1).
\end{equation}
No endpoint derivative or sampled minimum is used.

\subsection{The weak-helper row: two separate anchor prices}
\label{subsec:helperfunding}
The unheaded reference \eqref{eq:helperbare} has two available
prices
\begin{equation}
\label{eq:twoprices}
 \beta_1=a_0-b_0>.24672,\qquad
 \beta_6\geq b_0-(\sqrt7-\sqrt6)>.1023.
\end{equation}
Indeed \(2.0392<f_6<2.0394\), so
\(.29874<b_0<.29878\), while
\(\sqrt7-\sqrt6<2.6458-2.4494=.1964\).
An adjacent omitted knot decreases the next slope and increases
\(\beta_6\). The bounds follow from
\(2.4494<\sqrt6<2.4495\) and
\(.4101<.58/\sqrt2<.4102\), by positive squaring.

Define
\[
 Q=Z_u-(1-3u/5)e_1-(u/10)P_6.
\]
Its first coordinate is zero, its second is \(.4u\), and its
third through sixth coordinates are \(-.1u\). Therefore
\begin{equation}
\label{eq:helperresidual}
 \norm{Q}^2=u^2/5,\qquad \ip{w}{Q}=0.
\end{equation}
Choose \(\mu=.22\). First pay the \emph{separate} anchor debits
\[
 \mu(1-.6u)<.22<\beta_1,\qquad
 \mu u/10<.022<\beta_6
 \quad(0<u<1).
\]
The remaining original coefficients are nonnegative. Only then
add the positive new head pressure \(\mu Z_u\).
The terminal coefficient at \(16\) remains unused. The full
one-vector square gives
\begin{align}
\label{eq:Gdefinition}
 L_{w+\mu Q}
 &>.987759+.44[h-a_0+(3a_0/5-f_6/10)u]-.00968u^2\notag\\
 &>G(u):=.747739+.44h(u)+.0542784u-.00968u^2,
\end{align}
since \(3a_0/5-f_6/10>.12336\).
This argument rewrites a positive-coefficient pressure
decomposition before using it; it does not subtract old lower
floors from an unknown signed pressure.

We give the exact small-margin seal. By \eqref{eq:hcurvature},
\(G''>.44\cdot7/(4\sqrt2)-.01936>.5\).
At \(u_1=.642\), positive square tests give
\[
 \sqrt{1.358}>1.1653325,\qquad
 \sqrt{.229836}<.4794122,
\]
whose comparison squares are \(1.35799983555625\) and
\(.22983605750884\). The difference exceeds \(.6859203\).
Since
\[
 .4850188\cdot1.4142136=.68592018321568<.6859203,
 \qquad1.4142136^2>2,
\]
we have \(h(.642)>.4850188\), and hence
\[
 G(.642)>.99200425728.
\]
For the derivative use
\[
 .858<1/\sqrt{1.358}<.859,\quad
 .592<.284/\sqrt{.229836}<.593,\quad
 .353<1/(2\sqrt2)<.354.
\]
They follow by multiplying the positive radical brackets
\(1.1653<\sqrt{1.358}<1.1654\) and
\(.4794121<\sqrt{.229836}<.4794122\).
Thus
\[
 -.094518<h'(.642)<-.093545,\qquad
 .00026136<G'(.642)<.00068948<.0007,
\]
because \(G'=.44h'+.0542784-.01936u\).
Strong convexity on the entire interval now gives
\begin{equation}
\label{eq:Gseal}
 G(u)>.99200425728-.0007^2
       =.99200376728>.992.
\end{equation}

Combining \eqref{eq:Fseal} and \eqref{eq:Gseal}, every one of
the four original initial cases supplies a single legally
funded reference \(W\) with
\begin{equation}
\label{eq:uniformcontract}
 L_W>.992,\quad
 \operatorname{supp}W\subset\{1,\ldots,16\},\quad
 \ip{W}{P_{16}}=f_{16},\quad
 \text{unused }\beta_{16}>.127.
\end{equation}
All corrections have sum zero and support at most \(7\);
the late omissions and the original endpoint-\(16\) price survive.

\subsection{Both signed fork directions}
\label{subsec:forkheight}
Let \(\Delta=S_{16}-f_{16}\geq0\). The pressure decomposition
of \(W\) keeps its original terminal coefficient. Therefore
\(\ip{x}{W}>A_W+\beta_{16}\Delta\), and the same norm identity
gives
\begin{equation}
\label{eq:slack}
 \norm{x-W}^2<.008-.254\Delta,\qquad
 x_{17}^2+x_{18}^2<.008-.254\Delta.
\end{equation}
The two true unit fork duals are
\[
 \frac{\Omega_f}{3}
   =\frac{P_{17}\pm e_{18}}{\sqrt{18}},
 \qquad f=17,18.
\]
For either actual fork height,
\begin{align}
\label{eq:truefork}
 t_f:=\ip{x}{\Omega_f/3}
 &<\frac{f_{16}+\Delta+\sqrt{.016-.508\Delta}}{\sqrt{18}}\notag\\
 &<\frac{3.59+.1265}{\sqrt{18}}<.876.
\end{align}
The numerator decreases with \(\Delta\): its derivative is
\(1-.254/\sqrt{.016-.508\Delta}<0\) on its nonempty domain.
The endpoint is handled by continuity. The seals are
\(f_{16}<3.59\), \(\sqrt{.016}<.1265\),
\(\sqrt{18}>4.2426\), and
\(.876\cdot4.2426=3.7165176>3.7165\).
No absolute value has replaced the signed last coordinate.

For an actual fork edge with distance cosine \(C\),
\eqref{eq:halfheight} gives
\(t_f\geq\sqrt{(C-1)/(C+1)}\). Hence
\begin{equation}
\label{eq:forkcosh}
 C<\frac{1+.876^2}{1-.876^2}<7.6,
\end{equation}
since \(.876^2=.767376\) and
\(7.6\cdot.232624=1.7679424>1.767376\).
In the complete fork-role list \src{L}, every far-next-three
role has \(g\geq\cos(\pi/5)>.809\) and exact cosine
\(C=9g+1>8.281\). Equation~\eqref{eq:forkcosh} removes that
entire role, including all its previously classified native/null
own alternatives.

\section{A local \(A_{18}\) lemma at an adjacent actual vertex}
\label{sec:localA18}
\begin{lemma}
\label{lem:localA18}
Let \(q\) be an actual vertex of a compact Coxeter
\(18\)-polytope whose whole incident diagram is \(A_{18}\).
Write \(z=E_q(q-y)\), \(\norm{y}<1\). Suppose
\[
 -\ip{z}{r}\leq\kappa E_q
 \quad\text{for every actual point }r\text{ of the original }P,
 \qquad\kappa<11/8.
\]
Then \(\ip{y}{u_{18}}>.78\) is impossible for the positive
normalized endpoint-dual direction.
\end{lemma}
The lemma is local at \(q\), which is not assumed farthest.
We give the new endpoint proof and the precise scope of its
earlier local inventory \src{A-local}, Sections 2--5.
That source's original \(B_{18}\) Section 1 is not used.
\begin{proof}
In the zero-sum subspace of \(\R^{19}\), for \(i\leq j\),
\[
 D_i=\norm{\omega_i}^2=\frac{2i(19-i)}{19},\qquad
 C_{ij}=\frac{2i(19-j)}{19}.
\]
Call \(5\leq i\leq14\) weak if
\(\ip{y}{\omega_i}\leq\sqrt{D_i}-.87\).
The complete inventory at comparison \(11/8\), established
before any endpoint pressure, leaves only:
\begin{enumerate}[label=(\alph*)]
\item near-four, meaning own \(4\) and white \(3\) at the adjacent
 old node, in the four literal side-order-\(5/6\) orientations;
\item a pure far-end-three helper
\[
 d_i=A_iq-\gamma_i\omega_i-\omega_e/2,\quad
 e=1\text{ or }18,\quad
 \varphi/2\leq\gamma_i<17/16,
\]
 with own \(5\), finite own at least \(7\), or dotted own.
\end{enumerate}
Here \(\varphi=(1+\sqrt5)/2\).
All isolated caps are strong. The inventory includes own zero,
every finite own, parallel and unbounded dotted owns.
Every positive point checks every unknown original wall.

Some calculations locate the numerical boundary.
Every ordinary edge ratio at least \(3/2\) is strong:
with \(a=3/2,k=11/8\), the square comparison for
\(q_0=\sqrt{D_i}\in[2.7,3.1]\) is the positive polynomial
\[
 -.21q_0^3+1.1876q_0^2-1.74q_0+.7569.
\]
It is concave with endpoint values \(.583074,.519626\).
The near-four ratio of side order \(\ell\) is
\[
 \frac{1/\sqrt2+(\ell-1)/(2\ell)}
      {\sqrt{(\ell+1)/(2\ell)}}.
\]
It increases in \(\ell\) and exceeds \(3/2\) at \(\ell=7\),
since \((67+42\sqrt2)/56>9/4\). Only \(i=5,6\) with white
\(i-1\), and \(i=14,13\) with white \(i+1\), can be weak.
For pure helpers the identity
\[
 \frac{(\gamma+1/(2\ell))^2}{(\ell+1)/(2\ell)}
 -(4\gamma-2)
 =\frac{(\gamma-1-1/(2\ell))^2}{(\ell+1)/(2\ell)}\geq0
\]
gives the stated cutoff at the weakened comparison.

The local input retains genuine helper nonadjacency at this
factor. Its proof uses the following exact actual products.
For adjacent helper owns \(A,B\in[.808,1.08]\), the mutual
product \(R\) satisfies
\[
\begin{array}{c|l}
\text{same end}&R\geq1-\frac{18}{19}AB-\frac{|A-B|}{19}>-1/2\\
\text{outward}&R\geq\frac12-AB+\frac{(A-1)(B-1)}{19}>-.7\\
\text{reversed}&R\geq\frac12-AB+\frac{(A+1)(B+1)}{19}>-.7.
\end{array}
\]
The native alphabet leaves \(R=0\) in the first case and
\(R=0,-1/2\) in the others. Outward \(R=0\) contradicts
\[
 (D_i-4/5)(D_{i+1}-4/5)-C_{i,i+1}^2
 =\frac25 C_{i,i+1}-\frac{24}{25}+\frac8{95}
 \geq884/475>0.
\]
For reversed \(R=0\), set \(l=i,r=18-i\).
The exact determinant would require
\[
 4lr A^2B^2-(l+1)(r+1)(A^2+B^2)-2AB
 -(l+1)A-(r+1)B+21/4=0.
\]
Its own derivatives are positive on that rectangle.
If one own is not \(5\), it exceeds \(.9\); evaluating the
other at \(.808\), the endpoints \(l=5,13\) give
\(.9601824,1.6961824>0\), and concavity in \(l\) excludes all.
For two native-five owns the norm-two calculation in basis
\(1,\varphi\) has nonzero rational determinant coefficient
\(475\) outward and \(-361\) reversed. This uses only native
cosine algebra. Thus remaining products are \(0\) for a
common helper and \(-1/2\) for opposite helpers.

After clearing old roots, the actual two-cap Schur matrices
are \(I-t\mathbf1\mathbf1^T\), \(t=(\ell-1)/(2\ell)\), or
the matrix with diagonal
\(\delta_l,\delta_r\in(1/2,3/5]\), off-diagonal \(-1/2\).
Both are positive with positive inverse. Their positive
projections are feasible against every original wall,
including the reverse wall at \(q\).
The exact cap-time/weight estimates in \src{A-local} give
weighted mean loss less than
\[
 1/10+(11/8)(7/16)=.7015625
 \quad\text{or}\quad1/10+(11/8)(3/10)=.5125.
\]
Both are less than \(.87\), excluding adjacent weak helpers.
None of these steps uses endpoint height \(>.8\).

Now assume the new height \(\ip{y}{u_{18}}>.78\).
Each weak near-four has an actual mate, formed by releasing
its white neighbor and retaining its cap. Positive finite
Gram inversion gives an actual adjacent vertex feasible
against every wall. Its edge cosine exceeds \(11\);
for example the left mates have
\[
 C_j=\sqrt{2(i-1)/19}
 \left((20-i)/2+(19-i)/\sqrt2\right)>11,\qquad i=5,6.
\]
At comparison \(11/8\), the normalized mate height exceeds
\(7/8\). A left mate has correlation \(<.15\) with \(u_{18}\).
The sum of these heights exceeds \(1.655\), whose square
exceeds \(2.3\), while the squared norm of the sum of the
directions is \(<2.3\). This contradicts \(\norm{y}<1\).
No left near-four remains.

Some \(k=5,6\) is now strong by helper nonadjacency, with
normalized height \(>.677\).
For a right mate \(j=14,15\), the correlation formula
\[
 \rho_{ab}=\sqrt{\frac{a(19-b)}{b(19-a)}}\quad(a\leq b)
\]
gives \(\rho_{18,j}<.46,\rho_{18,k}<.161,\rho_{j,k}<.407\).
For one vector \(V=.5u_{18}+.6u_j+.4u_k\),
\[
 \ip{y}{V}>.5(.78)+.6(7/8)+.4(.677)=1.1858,\quad
 \norm{V}^2<.77+.276+.0644+.19536=1.30576.
\]
Equation~\eqref{eq:square} gives \(\norm{y}^2>1.06584\).
All near-four orientations are excluded at \(.78\).
Only nonadjacent pure weak helpers remain.

Use rational heights \(H_{19-i}=H_i\), linear on \(0\)--\(5\)
and \(14\)--\(19\), with
\[
 H_5=1.844,\ H_6=1.995,\ H_7=2.103,\ H_8=2.173,\
 H_9=H_{10}=2.207.
\]
Their positive slopes are \(.3688,.151,.108,.070,.034\),
then zero. The squares \((H_i+.87)^2\), \(5\leq i\leq9\), are
\[
 7.365796,\ 8.208225,\ 8.838729,\ 9.259849,\ 9.467929,
\]
strictly below \(140/19,156/19,168/19,176/19,180/19\).
Thus every retained positive dual is paid by a strong floor.
If \(5\) is weak use chord \(0\)--\(6\), reflect at \(14\),
and omit each internal weak knot by its neighbor chord.
After chords set \(w_i=(H_i-H_{i-1})/\sqrt2\).
The reference is zero-sum with demand \(E_0=\norm{w}^2\).

Before internal holes, the exact energies for zero, one and
two weak ends are \(.7205882,.70082285,.6810575\).
Internal curvatures are \(.043,.038,.036,.034\) and reflections;
even summing all eight losses gives
\[
 2(.043^2+.038^2+.036^2+.034^2)/4=.0028725<.002873.
\]
Nonadjacency makes the losses legitimate. End chords do not
change any eligible internal curvature. Hence
\begin{equation}
\label{eq:Alocalbase}
 E_0>
 \begin{cases}
 .6781845,&14\text{ weak: }m=6,\ K=13,\ H_K=1.995,\\
 .69794985,&14\text{ strong: }m=5,\ K=14,\ H_K=1.844.
 \end{cases}
\end{equation}
The surviving individual anchor prices are
\[
 \beta_{13}\geq(7H_6/6-H_7)/2=.11225,\qquad
 \beta_{14}\geq(6H_5/5-H_6)/2=.1089.
\]
Adjacent holes increase the anchor price. Omitted knots have
price zero; no missing strong \(7\) or \(18\) floor is paid.

Define
\[
 G_0=-e_{19}+\frac1m\sum_{\text{last }m}e_i,
\]
with \(19\) included. Then
\[
 \norm{G_0}^2=(m-1)/m,\quad \ip{w}{G_0}=0,\quad
 G_0=\omega_{18}/\sqrt2-\omega_K/(m\sqrt2).
\]
Use the new weight \(\mu=.6\).
The debit \(.6/(m\sqrt2)\) is \(<.071\) for \(m=6\),
\(<.085\) for \(m=5\), below the separate anchor prices.
The positive endpoint coefficient is paid by its actual height.
Here
\[
 -y_{19}>.78\sqrt{18/19}>.7591,
\]
since \(19(.7591)^2=10.94842339<18(.78)^2=10.9512\).
The last-block means are
\(H_K/(m\sqrt2)<.2352\) or \(<.261\).
The same local-\(y\) square, retaining \(\ip{w}{G_0}=0\), gives
\[
\begin{array}{ll}
m=6:&\norm{y}^2>
 .6781845+1.2(.7591-.2352)-.36(5/6)=1.0068645,\\
m=5:&\norm{y}^2>
 .69794985+1.2(.7591-.261)-.36(4/5)=1.00766985.
\end{array}
\]
Both contradict \(\norm{y}<1\), proving the lemma with the
local inventory established in Appendix~\ref{app:localAinputs}.
\end{proof}
Appendix~\ref{app:localAinputs} supplies the complete
local own/white inventory, isolated-cap point bounds,
positive mixed weights and native-five determinant
expansion used above. It uses only real geometry and
comparison \(11/8\), with no original-\(B_{18}\) premise.
The original-\(x\) energy is never added to this local-\(y\) budget.

\section{The near/far pair and the final composition}
\label{sec:composition}
\subsection{Every actual near/far fork edge is short}
For a near-three or far-three fork role, the complete finite-white
classification \src{T} gives an actual other endpoint \(q\)
with whole diagram \(A_{18}\). The positive reverse direction
at \(q\) is exactly its normalized new endpoint dual: \(p\)
has nonzero product with only that new incident root.
Let \(C\) be the original edge cosine,
\(\beta=\sqrt{C^2-1}/C\), and \(t=t_f<.876\).

Suppose \(C\geq5.6\). Then
\[
 E_q/E=C-\sqrt{C^2-1}\,t
       >5.6-.876(5.51)=.77324.
\]
The function increases for \(C\geq5.6\).
Here \(5.51^2=30.3601>30.36\) and
\((11/8)(.77324)>1\), so \(E/E_q<11/8\).
Also \(\beta>.9839\), because
\(31.36(1-.9839^2)=1.0016631744>1\).
The positive reverse height satisfies
\[
 \ip{y}{u_{18}}=\frac{\beta-t}{1-\beta t}>.78.
\]
Its denominator is positive and its limiting gap is
\[
 .9839-.876-.78(1-.9839\cdot.876)=.000179192>0.
\]
The original maximum supplies all-point comparison at \(q\).
Lemma~\ref{lem:localA18} excludes \(C\geq5.6\).
Every actual near/far fork edge therefore has \(C<5.6\).

\subsection{The entire near/far family}
Suppose the actual fork opposites are near and far:
\[
 c_N=Np-\gamma\Omega_f-\Omega_o/2,\qquad
 c_F=Fp-\delta\Omega_o-\Omega_1/2,
\]
with \(f,o\) the different actual forks.
Their times and exact edge cosines are
\[
 N^2=9\gamma^2+8\gamma+5/4,\quad
 F^2=9\delta^2+\delta-1/2,\quad
 C_N=\frac{18\gamma+8}{\sqrt{19}},\quad
 C_F=\frac{18\delta+1}{\sqrt{19}}.
\]
The short-edge bounds and \(\sqrt{19}<4.359\) give
\(\gamma<.912,\delta<1.301\).
Put \(U=3\gamma+4/3,V=3\delta+1/6,a^2=19/36\).
Then \(N^2=U^2-a^2,F^2=V^2-a^2\), and
\[
 (UV-a^2)^2-N^2F^2=a^2(U-V)^2\geq0,\qquad UV-a^2>0.
\]
The actual mutual magnitude consequently satisfies
\begin{align*}
 \lambda=-\ip{c_N}{c_F}
 &=NF-8\gamma\delta-\gamma/2-(9/2)\delta-1/4\\
 &\leq\delta(\gamma-1/2)-5/9\\
 &<1.301(.412)-5/9=.536012-5/9<0.
\end{align*}
This contradicts all-wall acuteness. If \(\gamma\leq1/2\)
the upper bound is already negative; the accepted own list
in fact has \(\gamma\geq\cos(\pi/4)>1/2\).
Zero, every finite, parallel and unbounded dotted mutual
are excluded at once, for all finite or dotted near/far owns.

\subsection{Exhausting the actual head and fork covers}
For a dotted head, \eqref{eq:forkcosh} removes every
far-next-three role from the complete \src{L} fork cover.
What remains is isolated, near-three, or far-three.
Their six unordered pairs are exhausted by
\begin{center}\small
\begin{tabular}{ll}
\toprule
Actual pair&Complete excluding argument\\ \midrule
isolated/isolated&\src{II}, every mutual and unequal own\\
isolated/near&\src{IN}, the entire family\\
isolated/far&\src{FF}, for every head\\
near/near&\src{T}, for every head\\
near/far&The complete argument above\\
far/far&\src{FF}, for every head\\ \bottomrule
\end{tabular}
\end{center}
Both assignments of the genuine signed forks are included.
Every pair is impossible, so no dotted head remains.

The head cover also includes own zero, finite and parallel
cases. The complete isolation proof \src{H} removes own
\(0,3,4\) by time constraints and own \(6\) by a positive
affine \([6,3]\) kernel. At own \(6\) an actual null is
\(\sqrt3H+2n_1+n_2\); at parallel own it is \(H+n_1\).
The remaining finite labels \(5,7,8,9\) and every integer
label at least \(10\) are excluded by the complete \src{FH}
conjunction. Every dotted own has just been excluded.
Thus the whole-\(D_{18}\) diagram is impossible.

Theorem~\ref{thm:tail} excludes every dimension at least \(19\).
At exact dimension \(18\), the fresh complete cover and its
row consumers, including Section~\ref{sec:b18}, leave only
whole \(D_{18}\) at the same original farthest vertex.
The head/fork argument excludes it. Hence dimension \(18\)
is impossible and Theorem~\ref{thm:main} follows at the
complete internal proof boundary, without dimension heredity.
The missing internal transcriptions identified in this draft
must be incorporated before it is a self-contained paper.

\section{Finite algorithms and conditional verification}
\label{sec:verification}
\subsection{The mathematical finite-cover calculation}
The rank covers are finite combinatorial arguments supported
by real inequalities, not field enumerations.
For fixed total rank \(n\), the forest calculation is:
\begin{enumerate}
\item List partitions of \(n\) into ranks of spherical irreducibles.
 Retain all exceptional finite types and classical types at
 their actual ranks. Keep \(I_2(m)\) as a symbolic ordinary
 rank-two family for unbounded \(m\), distinguishing its
 exceptional actual typed profiles.
\item Mark at most one allowed \(I/J/N_7/O_7\) profile by the
 actual-wall compatibility lemma. Apply joint bounds only to
 disjoint sets of components, counting each defect once.
\item Sum proved same-vector energy bounds. Discard a forest only
 when a strict lower bound is at least \(1\). If a dominant
 rank needs a fresh local layer, retain it until that
 actual-rank lemma is established.
\item Return the retained grouped necessary cover and the
 inequality discarding every other partition.
\end{enumerate}
Termination follows from finitely many rank partitions and
finite type choices at each rank. The unbounded native
parameter is represented by inequalities proved over its
entire domain. Completeness uses spherical classification,
compact acute simplicity, complete actual-defect compatibility,
and spectator-uniform component bounds. The output is a
necessary cover; no realization follows. A program testing
only finitely many native labels would not establish it.

Local own/white calculations have a different domain.
After deleting a specified original node, classify every
finite extension of the retained spherical forest by its
actual opposite wall, including own zero, all finite owns,
parallel and dotted owns. Finite white extensions are bounded
by spherical classification. Infinite own families are closed
by analytic inequalities over complete intervals.
Termination of a finite classification does not justify
discarding a continuous branch. Completeness must include
spectator attachments, affine kernels and every original
wall for every auxiliary point.

\subsection{The distinct role of scripts}
This draft introduces no scientific computation and reports
no new execution. The final continuous seals above are
written rational inequalities and calculus arguments.
A source hash, review receipt or saved maximum does not
replace them.

For a reproducibility checker attached to a finite table,
its mathematical statement must specify the exact input
domain, symbolic own parameters, exhaustive loops or recursion,
the terminating measure, exact tests and returned cover.
Its correctness proof must map every admissible geometric
case to a generated entry and each rejected entry to a
valid exclusion. The execution record then states what
that particular program returned on those inputs.
Algorithm theorem, execution claim and geometric consumer
are distinct. Historical receipts identified in the ledger
are not a new complete replay of this revision.

\subsection{The formal boundary}
A suitable Lean kernel may take as explicit parameters:
\begin{enumerate}
\item precise finite/affine classification and compact acute
 geometry statements supplied from the literature;
\item precise internal component and local-cap statements whose
 full written project proofs have been reviewed;
\item a precisely typed necessary-cover output, if exact
 enumeration is provided as a conditional premise.
\end{enumerate}
It may verify rank sums, disjoint-cover subtraction,
one-vector algebra, rational implications and the final
composition giving \(n\leq17\).
An internal geometric statement remains an internal proof
obligation when used as a parameter; it does not become
a literature theorem. Logic maps must list parameters as
well as printed axioms.

The selected historical publication kernel and its separate
actual axiom audits do not certify all new sources here.
No module was compiled or audited for this draft.
Neither the new whole-\(D_{18}\) geometry nor all its
references are claimed to be formalized. Mathematical
review, exact replay, Lean compilation, actual axiom
auditing and PDF production are separate stages.

\appendix

\section{Ordinary component proofs and actual defects}
\label{app:ordinary}
Throughout the proof appendices, \(F(t)=2t/(1+t^2)\) denotes
the full-/half-tangent conversion.
This appendix collects the contact calculations behind the small-component
rates. Each is an argument on an actual whole component with arbitrary
finite spectators, using Section~\ref{sec:geometry}. No field condition
is present. We use the spherical/affine classification and the standard
Lann\'er classification: a minimal indefinite Coxeter diagram with all
proper subdiagrams positive definite has at most five vertices.
A positive affine principal is already impossible by the actual null
argument in Section~\ref{sec:geometry}.

\begin{proposition}[Ordinary costs and the complete small defects]
\label{app:prop:ordinary}
All the ordinary component bounds in Section~\ref{sec:components}
hold with arbitrary finite spectators. For the moment use the weaker
bound \(B_3,H_3>1/4\); the improvement to \(2/7\) is proved immediately
below. The only exceptions to the rank-over-\(14\) small-component
baseline are one possible isolated short \(I\), its separately funded
\(J\) alternative, or one of the two typed rank-two defects.
\end{proposition}
Write \(m@j\) for actual contact \(\cos(\pi/m)\) at old node \(j\).
A missing contact is zero. For a wall put \(h=A^2\), own cost
\(q\) and full foreign cost \(Q\), so \(h=q+Q-1>0\).
The following calculations prove the proposition; their tables are
finite mathematical case analyses rather than empirical output.

\begin{lemma}[Common-wall forcing and isolation]\label{app:lem:ordinaryforcing}
Suppose the energy of a component is at most $e$. A pair for which
the right side of \eqref{eq:pairupper} is below $K_{ij}$ has one
actual opposite. If one wall closes every direction of the component,
then, with $v_i=(C\beta)_i/\sqrt{d_i}$ and
$R_e^2=4e/(1+e)^2$,
\begin{equation}\label{app:eq:commoncost}
 q-1\le h\le R_e^2\min_i v_i^2.
\end{equation}
If $h<1/4$, it has no foreign contact. In a path, a wall closing
every direction can have multiple contacts only at the two endpoints;
in a $D_4$ star it has at most one.
\end{lemma}
\begin{proof}
The first assertion is \eqref{eq:pairupper}. For a closed direction,
$\tanh L_i=A/v_i$ and $\tau_i^2\le e$, which gives
\eqref{app:eq:commoncost}. Every nonzero foreign cost is at least $1/4$
and is strictly below $h$ by \eqref{eq:foreigncost}; hence the isolation
assertion. If the old path between two contacts omits a direction
closed by the wall, deleting that direction leaves a cycle in an
actual finite replacement. This is impossible. The same argument
applies to every two contacts in a $D_4$ star. Finally, a single-contact
own cost $q=1$ gives an affine principal set, impossible even in the
presence of further spectators.
\end{proof}

\subsection{The short directions and their exact debits}

For a short $A_1$ opposite put $b=\beta^2$, $h=A^2$, and let $Q$
be its full foreign cost. Then
\[
 h=b+Q-1,\qquad h<b/3,\qquad Q<1-2b/3.
\]
Equation~\eqref{eq:foreigncost} rules out two foreign components.
With one, its cost must be the sole value below $1/3$, namely an
$A_1$ contact $3$ of cost $1/4$. Thus the only possibilities are
\[
 I:\ \hbox{isolated dotted }A_1,\quad 0<h<1/2;
 \qquad
 J:\ \hbox{dotted }A_1+A_{1,3},\quad 1/3<h<3/8.
\]
For completeness, with one foreign component,
$Q<h=b+Q-1$ forces $b>1$, and $Q<1/3$ then forces $Q=1/4$.
For $J$, the heavy edge contributes more than $A/2$ to
$x\cdot v$. The light shadow contributes at least
$A(1-\sqrt{1-1/(4h)})$; this is at least $A/2$ if $h\le1/3$.
Strict \eqref{eq:actualinterior} therefore gives the asserted lower bound.
Native acuteness excludes two short walls. Indeed their height
product is below $1/2$, so their product must be zero. Disjoint
supports give a negative product, heavy overlaps give a spatial
product above $1/2$, and a sole light/light overlap of two $J$ walls
gives $1/4$, below their height product. None is zero. A $J$ heavy
direction has $\tanh^2L>4/13$, hence energy $>9/100$; its light
direction has $\tanh^2L>2/3>16/25$, hence energy $>1/4$.
Thus $J$ funds both $A_1$ factors above their combined ordinary
baseline $2/10$. Every nonshort direction pays at least
$5-2\sqrt6>1/10$, since $\sqrt6<49/20$. We may conservatively
discard the one short direction, but only the isolated type $I$
can create an unfunded debit.

The rank-two exceptions are
\begin{equation}\label{app:eq:defects}
\begin{array}{c|c|c}
\hbox{old block / new isolated contact} & h=A^2 & \hbox{energy}\\ \hline
A_2/7 & (4/3)\cos^2(\pi/7)-1\in(0,1/12) & >9/100\\
I_2(7)/3 & (1/4)\csc^2(\pi/7)-1\in(0,1/3) & >9/100.
\end{array}
\end{equation}
Their spatial supports lie entirely in their own blocks.
Two defects in different blocks, or a defect and a short $A_1$,
would have actual wall product $-\sqrt{h_1h_2}\in(-1/2,0)$.
That interval is forbidden: distinct acute unit walls have product
$0$ or at most $-1/2$. Therefore
\begin{equation}\label{app:eq:debits}
\text{at most one short }A_1,\quad\hbox{or one rank-two defect,
but never both.}
\end{equation}
This argument is independent of every other component type.

\subsection{The two rank-two estimates}

For $A_2$, set $e_2=319/2000$ and $\sinh^2L_0=31/50$.
If both edges have length at least $L_0$, then their half-tangents
exceed $173/500$ and energy is greater than
$4\cdot173^2/(3\cdot500^2)>e_2$. A shorter edge cannot have two
positive own contacts: $3/3$ is affine, while $3/4$ or larger gives
$\sinh^2L\ge4\sqrt2/9>31/50$, increasing with either contact
and with $Q$. If only the other node is contacted, apply the same
test to that shorter contacted direction. For one contact of square
$b$, the shorter-edge condition is
\begin{equation}\label{app:eq:A2short}
 h=4b/3+Q-1<124b/243,\qquad Q<1-200b/243.
\end{equation}
Every positive foreign cost must be below $h$. A label $3$ gives
$h<1/4$, hence isolation and negative height. For a label $4$,
$h<62/243$ forces each foreign cost to be $1/4$; the only positive
possibility is $Q=1/2$, but then $h=1/6<1/4$. A label $5$ has
$Q<1/2$, hence at most one foreign component, whose cost exceeds
$h$ because its own cost $\ph^2/3$ is below one. A label $6$ is
affine. For a finite label $m\ge7$, $b>.81$ gives $Q<1/3$;
thus $Q=0$ or $Q=1/4$. The latter is possible in
\eqref{app:eq:A2short} only for $b<729/800$, namely $m=7,8,9,10$.
Indeed $\cos^2(\pi/11)>45/49>729/800$, using $\pi<22/7$.
A dotted contact has $Q<1/4$ and is isolated.

For an isolated finite contact the same wall closes the mate,
since its reciprocal ratio squared is $b/(4b-3)>1$.
For $m\ge8$ that actual mate has
$\tanh^2L\ge6\sqrt2-8>400/841$, so its energy exceeds
$4/25>e_2$. For an isolated dotted contact,
$b<243/200$ and its distinct mate has shadow
$\tanh^2L\ge b/(4b-3)>81/124>16/25$.
For a mixed $m=7,8,9,10$ contact, the mate's white cost is
$b+1/4>1$; the distinct shadow instead gives
$\tanh^2L\ge b/(4b-9/4)>11/17>16/25$.
Both give energy greater than $1/4$. Thus the only exception is
isolated new $7$ at old $A_2$.
The root $t=\cos(2\pi/7)$ of $8t^3+4t^2-4t-1$ lies in
$(623/1000,5/8)$: the polynomial changes sign there and is increasing
on $[1/2,1]$. Hence $b>1623/2000$ and its actual mate satisfies
$\tanh^2L=4-3/b>164/541>3600/11881$, proving its cost $>9/100$.

For an old $I_2(m)$, $m\ge4$, put $k=\cos(\pi/m)$ and
$d=(1-k^2)^{-1}$, so $C=d\left(\begin{smallmatrix}1&k\\k&1\end{smallmatrix}\right)$.
At energy at most $1/7$, \eqref{eq:pairupper} forces one common
opposite. Two contacts give an actual
$\cosh L\ge\cot(\pi/(2m))/\sqrt3>11/8$, with individual energy
$>3/19>1/7$. For one contact $\beta$,
\begin{equation}\label{app:eq:dihedralcut}
 q(1-7k^2/16)\le1,\qquad q=d\beta^2.
\end{equation}
Its minimum left side is $(9d+7)/64$. Since
$m\ge8$ gives $d\ge4+2\sqrt2>19/3$, only four old labels remain:
\[
\begin{array}{c|l}
m&\text{consequence of \eqref{app:eq:dihedralcut}}\\ \hline
4&\beta^2\le16/25<\cos^2(\pi/5):\ 3,4;\ h\le7/32<1/4,\ q\le1;\\
5&d>25/9:\ \text{only }3;\ h<7/32,\ q<1;\\
6&\text{only }3;\ q=1;\\
7&d>5:\ \text{only }3;\ Q<3/16.
\end{array}
\]
Isolation or the affine principal excludes the first three rows.
The last is isolated new $3$ at old $I_2(7)$, with $0<h<1/3$.
Its actual mate has $\cosh L=\cot(\pi/7)/\sqrt3>109/91$ by the
same cubic bracket; the squared comparison has positive numerator
$2652$. Thus its energy also exceeds $9/100$, completing
\eqref{app:eq:defects} and \eqref{app:eq:debits}.

\subsection{Rank three and four: one common-wall calculation}

For energy $1/4$, adjacent correlation at least $1/\sqrt3$ forces
a common opposite: the lowest cleared comparison is
$74/\sqrt3>39$. At energy $1/3$ the threshold $1/\sqrt2$ works,
since $51/\sqrt2>34$. For $A_4$ at energy $2/7$, the adjacent
correlations are $\sqrt{3/8},2/3,\sqrt{3/8}$; the forcing test is
$k^3+16k^2+36k-28>0$, whose lowest value is positive because
$291\sqrt6>704$. Direct inversion verifies the indicated
correlation thresholds for $A_3,B_3,H_3$ and $B_4,D_4,F_4$.
For $H_4$ use energy $11/36<1/3$. Thus a common wall closes all
directions in every case.

Number path bonds as $A_3:(3,3)$, $B_3:(4,3)$, $H_3:(5,3)$,
$A_4:(3,3,3)$, $B_4:(4,3,3)$, $F_4:(3,4,3)$, and
$H_4:(5,3,3)$. By Lemma~\ref{app:lem:ordinaryforcing}, multiple
contacts can only be at the endpoints. Two endpoint $3$ contacts
in $A_3,A_4$ are affine; the next minimum is $3/4$. In the other
rows below even $3/3$ is excluded. The displayed values are lower
bounds for $q-R_e^2v_j^2$, contradicting \eqref{app:eq:commoncost}:
\[
\begin{array}{c|c|c}
T&j&\text{lower bound}\\ \hline
A_3&2&177/200+9\sqrt2/100>1\\
A_4&2&(230+15\sqrt2)/243>1\\
B_3&2&(77+18\sqrt2)/100>1\\
B_4&3&15/16+\sqrt2/8>1\\
F_4&2&15/16+\sqrt2/4>1\\
H_3&2&(72\ph+95)/100>1.
\end{array}
\]
$D_4$ already has at most one contact. In $H_4$, two endpoint
contacts retain, after deleting old node $4$, a path at least
$(3,5,3)$ of determinant $(5-4\ph)/16<0$. Hence all contacts
are single.

For a single contact at $j$ let
$\sigma_j=\min_i(v_i^2/q)$. Formula \eqref{app:eq:commoncost} gives
$q\le(1-R_e^2\sigma_j)^{-1}$. The complete inverse calculation
is as follows; entries are in path order, and $D_4$ means leaf,center.
\begin{center}\small
\begin{tabular}{ccll}
\toprule
$T$&$d_j$&$\sigma_j$&upper $q$\\ \midrule
$A_3$&$3/2,2,3/2$&$1/9,1/3,1/9$&$225/209,75/59,225/209$\\
$A_4$&$8/5,12/5,12/5,8/5$&$1/16,1/6,1/6,1/16$&$162/155,243/215,243/215,162/155$\\
$B_3$&$3,4,2$&$1/3,1/2,1/3$&$75/59,25/17,75/59$\\
$D_4$&$2,4$&$1/4,1/2$&$16/13,8/5$\\
$B_4$&$4,6,4,2$&$1/4,1/3,1/2,1/4$&$16/13,4/3,8/5,16/13$\\
$F_4$&$4,12,12,4$&$1/2,2/3,2/3,1/2$&$8/5,2,2,8/5$\\ \bottomrule
\end{tabular}\end{center}
Here $R_e^2$ is respectively $16/25$, $56/81$ for $A_4$, and
$3/4$ for the $1/3$ rows. Dividing the upper $q$ by $d_j$
gives the finite native lists:
$A_3$ endpoints $3/4/5$, middle $3/4$;
$A_4$ endpoints $3/4$, inner $3$;
$B_3$ nodes $3;3;3/4$;
$D_4$ leaf $3/4$, center $3$;
$B_4$ nodes $3;\varnothing;3;3/4$;
$F_4$ endpoints $3$, inner $\varnothing$.
Each has $q\le1$. Equality is affine; in every other listed case
\eqref{app:eq:commoncost} gives $h<1/4$ (the largest needed bounds are
$16/75,7/155,28/405,3/16$). Isolation then gives $h=q-1<0$.
This proves the six crystallographic and $A$ rows.

For $H_3$, the minimum contact $\beta=1/2$ at node $1$ or $2$
gives tests at node $3$ equal to
$(199\ph+263)/500$ and $(61\ph+77)/125$, both above one.
At node $3$, $\beta\ge1/\sqrt2$ gives the test at node $1$
at least $(27\ph+118)/150>1$. Only $3@3$ remains, with
$q=(\ph+2)/4<1$ and $h\le4\ph^4/75<\ph/4$.
No spectator gives negative height; one has cost $Q>h$; two have
$Q\ge1/2$ and hence $h\ge q-1/2=\ph/4$. All are impossible.
For $H_4$, $R_e^2=1584/2209$; at $\beta=1/2$ the four tests
(contact node, test direction) are
\[
(1,4):\frac{4938\ph+3688}{2209},\quad
(2,4):\frac{7209\ph+5334}{2209},\quad
(3,1):\frac{2500\ph+2667}{2209},\quad
(4,1):\frac{625\ph+1219}{2209}.
\]
All exceed one, the last by $\ph>198/125$; larger contacts only
increase these positive tests. This proves $H_4>11/36$.

\subsection{The two intermediate \texorpdfstring{$A$}{A} components}

Use $C_{ij}=2\min(i,j)(N-\max(i,j))/N$, with $N=6,7$.
For $A_5$, distinct $2/3$ and $3/4$ pairs each demand a sum above
$109/100$. The vector $u_2+2u_3+u_4$ has square
$7+4\sqrt2<443/35$, and
\[
(109/50)^2-(3/8)(443/35)=209/35000>0.
\]
For $A_6$, the three pair demands are
\[
2249/2000,\qquad11547/10000,\qquad2249/2000;
\]
the vector $u_2+2u_3+2u_4+u_5$ has square
$<33781/1250$, with demand $34037/10000$ and Cauchy gap
$2181583/700000000$ at energy $3/7$.
These pair bounds follow by substitution in
$F(s/2)^2<K$. Thus a common $2/3$ occurs for $A_5$, and a common
$2/3$ or $3/4$ for $A_6$, up to reversal.

Each retained path allows one contact. The full two-contact list is
\begin{center}\small
\begin{tabular}{lll}
\toprule
Component/pair&positions&necessary finite-white labels\\ \midrule
$A_5,2/3$&$1,3$&$3/3$ or $4@1+3@3$\\
&$1,4;1,5;2,4;2,5$&both $3$\\
&$2,3$&both in $3,4,5$\\
$A_6,2/3$&$1,3$&$3/3$ or $4@1+3@3$\\
&$1,4;1,5;1,6;2,4;2,5;2,6$&both $3$\\
&$2,3$&$3,4,5$ at $2$; $3,4$ at $3$\\
$A_6,3/4$&$3,4$&both in $3,4,5$\\
&other eight left/right pairs&both $3$\\ \bottomrule
\end{tabular}\end{center}
Indeed deletion of either closed node leaves a connected finite
path or $D/E_6$ tree: a $4$ can only be terminal, and $D/E$ bonds
are all $3$. All-$3$ doubles are affine cycles. The adjacent
non-$3/3$ doubles force an actual edge of cosh at least
$(3+\sqrt2)/2>11/5$ in $A_5$; for $A_6$ the minimum formulas are
$(5+4\sqrt2)/\sqrt7$, $(4\sqrt2+6)/\sqrt{14}$, and
$(12+9\sqrt2)/\sqrt{42}$, all above $5/2$.
The $A_5$ terminal-$4$ row has $Q=0$ and cosh
$(3+\sqrt2)/2$. The $A_6$ version has white cost $9/10$, hence
$Q=0$, and is treated below. These cosh thresholds imply individual
energies above $3/8$ and $3/7$.

For single $3$ contacts the complete own costs and strict white
bounds are
\[
\begin{array}{c|c|c|c}
\text{component/pair}&j&q&Q\text{ upper}\\ \hline
A_5,2/3&1,2,3,4,5&5/12,2/3,3/4,2/3,5/12&2/3,2/3,5/8,1/2,5/8\\
A_6,2/3&1,2,3,4,5,6&3/7,5/7,6/7,6/7,5/7,3/7&2/3,2/3,3/5,2/5,2/5,3/5\\
A_6,3/4&1/6;2/5;3/4&3/7;5/7;6/7&5/8;1/2;5/8.
\end{array}
\]
End and small-$Q$ rows have $h<1/4$, or only one foreign
component with cost $Q>h$. In $A_5$ at $j=2,3$ two foreign
components would each exceed $1/3,1/4$, respectively, forcing
$Q\ge2/3$; the same gap treats $A_6,j=2$.
In the remaining $A_6,j=3,4$ rows, two foreign bonds must be $3$
since $Q<5/8$; together with the contacted old node and its two
neighbors they form affine two-fork $D_5$. This excludes all single
$3$ cases. A single internal $4$ gives the affine $(3,3,4)$ star;
an end $4$ has $Q\le1/3$ and $h<1/4$, again impossible.

A single $5$ is at a retained path end; its $H_3/H_4$ white cost
exceeds $3/4$, so $Q=0$. All cases are covered by these actual
positive replacements:
\[
\begin{array}{c|c|c|c}
\text{component/pair}&\text{contact}&\text{actual direction}&\cosh L\\ \hline
A_5,2/3&1\text{ or }5&4\text{ or }2&\ph^3/\sqrt3\\
&2&3&\ph^2\\
&3&2&\sqrt3\ph^3\\
A_6,2/3&1&4&3\ph^3/\sqrt{21}\\
&2&3&4\ph^2/\sqrt{14}\\
A_6,3/4&1\text{ or }6&4\text{ or }3&3\ph^3/\sqrt{21}\\
&3\text{ or }4&4\text{ or }3&9\ph^3/\sqrt{21}.
\end{array}
\]
The first group exceeds $11/5$, the second $5/2$, using
$\ph>8/5$, $\sqrt{21}<23/5$, $\sqrt{14}<15/4$.
No new wall is inferred: each replacement contains the already
actual $c$. Finally the $A_6$ double $4@1+3@3$ has the positive
actual future point
\[
 \xi=c+(\sqrt2-1/2)n_1+(\sqrt2-1)n_2+(\sqrt2/2)n_3,
 \qquad \xi^2=5/4-\sqrt2=-k.
\]
Its products with $c,n_1,n_3$ vanish and its product with $n_2$
is $(\sqrt2-3)/4<0$. Here $A^2=(5+4\sqrt2)/7$ and
$A^2-9k=(335-236\sqrt2)/28>0$; \eqref{eq:pointenergy} gives
energy $>(A-\sqrt k)/(A+\sqrt k)>1/2$.

\subsection{The four intermediate \texorpdfstring{$B$}{B} components}

Use $u_i=(e_1+\cdots+e_i)/\sqrt i$ and, for $i<r$,
$\omega_i=\sqrt2(e_1+\cdots+e_i)$, while
$\omega_r=e_1+\cdots+e_r$.
At energies $3/7,1/2,4/7$ the relevant forcing signs are
$127\sqrt{2/3}-100$, $(69\sqrt{6/7}-52)/7$, and
$(199\sqrt{6/7}-176)/7$, all positive.
Thus $B_5,B_6$ have a common $2/3$, and $B_7,B_8$ a common $6/7$.

For $B_5,B_6$, deletion $2$ retains $A_1+B_3/B_4$.
The right contact is $3@3$, with the sole $B_5$ exception $3@5$
giving $F_4$; deletion $3$ allows one left $3/4/5$.
For $3@3$, own $0$ is affine, own $3$ is a triangle, and own
$\ge4$ gives $\cosh L_3\ge2+3/\sqrt2>3$.
With no right contact the complete choices are $3/4/5$ at $1$ or
$2$: $3@2$ and either $4$ are affine; $3@1$ has
$h=Q-1/2<1/6$ and closes a forbidden foreign direction.
The two $5$ cases have $Q=0$: $5@1$ also closes $4$ with
$\cosh L_4=\ph^3>3$, and $5@2$ has
$\cosh^2L_3=2\ph^2/(2-\ph)>9$.
For the $B_5$ short-end $3@5$, finite $F_4$ forbids left $1$
and all spectators; own $0$ is affine and the other replacement
forces own $3$, giving $\cosh L_2=2+\sqrt2$.
Every branch exceeds the individual threshold $11/5$.

For the $6/7$ pair, a wall without a right contact has a single
left $3$ or endpoint $4$. Apart from far $3@1$ these give an
affine fork, double-$4$, or star. A minimal nonpositive principal
for the far branch must contain $c,7$, since deletion $7$ is
finite; their connecting path $c,1,\ldots,7$ has eight vertices.
Spectators attach through $c$ and do not shorten it. This violates
the Lann\'er order bound, excluding the branch with all spectators.
In $B_7$, let $\beta$ be the contact to the retained right $A_1$:
\[
\begin{array}{c|l|l}
\beta/\text{own }6&\text{left contacts}&\text{pressure or rejection}\\ \hline
\beta\ge\cos(\pi/5),\ \text{any}&\text{none}&\cosh L_6\ge\sqrt6\cot(\pi/5)>3\\
\beta_4,0&\text{any allowed}&\text{affine }(4,4)\\
\beta_4,>0&\text{none}&\cosh L_6\ge\sqrt{24}\\
\beta_3,>0&\text{none}&\cosh L_6\ge2+\sqrt2\\
\beta_3,0&\text{none or }1&\text{affine }F_4\\
\beta_3,0&\text{one }j=2,\ldots,5&\cosh L_6\ge1+3/\sqrt2>3.
\end{array}
\]
Here $\beta_m=\cos(\pi/m)$. The last row retains the full foreign
cost: its white residual is
$\Delta\le3/4-j(6-j)/12\le1/3$, and
$v_6\ge\sqrt{12}(j/12+1/(2\sqrt2))$.
Thus no allowed $E_7$ extension is silently removed.

In $B_8$ the retained right $B_2$ has one simple $3$ at either
end. At long end $7$, own $0/3$ is affine and own $4$ gives
$\cosh L_7\ge6+7/\sqrt2>7$.
At short end $8$, own $>0$ gives
$\cosh L_6\ge\sqrt6+\sqrt3>11/3$.
At own $0$, the left-white budget permits only $3@1$ or $3@5$:
none or $1$ gives affine $F_4$, but $3@5$ is an allowed bridge,
with $\Delta\le1/12$ and
$v_6=\sqrt{12}(5/12+1/(2\sqrt2))$.
It has $\cosh L_6\ge5+3\sqrt2>9$; it is estimated, not removed
by an incorrect $F_4$ sealing claim. These complete lists prove
the four asserted $B$ costs.

\subsection{Extended stars, \texorpdfstring{$D_7$}{D7}, and the \texorpdfstring{$D_8$}{D8} closure}

Let $S$ be a central $D_4$ star with an extended arm. If one wall
closes its center and three neighbors, every minimal nonpositive
principal contains $c$ and all four nodes of $S$, since deletion
of any of them is an actual finite replacement. Compactness and
the Lann\'er order bound make that principal exactly $S+c$.
Two contacts give a nonpositive proper triangle or cycle.
A sole center $3$ is affine, and a larger center contact has a
nonpositive proper star. A sole leaf contact is $3/4/5$ from its
finite proper path: $3$ gives finite $D_5$, $4$ affine, leaving
only $5$. If that leaf extends, it retains a nonpositive $(5,3,3)$
star; if another arm extends, it retains a nonpositive $(5,3,3,3)$
path. Each omits a supposedly closed $S$ direction, a contradiction.
Further negative contacts cannot repair the positive nonpositive
quadratic witness. Thus the four opposites cannot all coincide.

For $D_5,D_6$ the minimum central correlations are
$\sqrt{3/5},\sqrt{2/3}$. At energies $5/14,3/7$ the forcing
signs are $544\sqrt{3/5}-362$ and $127\sqrt{2/3}-100$, positive.
For $E_6,E_7,E_8$ the arm lengths are $(1,2,L)$ with $L=2,3,4$.
In raw norm-two coordinates the central Schur residual is
$\delta=2-[1/2+2/3+L/(L+1)]$; on an arm of length $m$, squared
central correlation is $t/(\delta+t)$, $t=m/(m+1)$.
The minima are $\sqrt3/2,\sqrt{6/7},\sqrt{15}/4$.
At energies $9/20,5/9,2/3$ even $4e/(1+e)^2$ lies below these
correlations. All four opposites would coincide, proving the five
extended-star costs.

For $D_7$ write $o=r-2$, $a=r-3$, with fork leaves $b,f$.
At energy at most $1/2$, correlation $\sqrt{4/5}$ forces a
common $o,a$. Deletion $a$ retains $A_3=b-o-f$, so its right
contact is none, middle $3$, or endpoint $3/4/5$; deletion $o$
allows one left contact. The complete surviving profiles are
\begin{equation}\label{app:eq:D7roles}
 4@a+3@o,\qquad\text{isolated }5@b,\qquad3@a+4@b,
\end{equation}
up to exchanging the fork leaves. To check the list: middle $3$
forces $4@a$ (own $0/3$ is affine); endpoint $5$ seals all white
extensions by $H_4$, and a contact at the removed $a$ leaves a
nonpositive $(3,3,3,5)$ path. Endpoint $4$ forces $3@a$, otherwise
an affine terminal-$4$ fork remains. Endpoint $3$ excludes every
left contact by cycles or retained $F_4$; a foreign $3$ is affine
$E_7$, a foreign $4$ is an affine terminal-$4$ fork, and larger
foreign bonds are nonfinite. Its isolated height is negative for
$r=7$ and it contains affine $E_8$ for $r\ge8$.
With no right contact, internal-left $3$ and endpoint $4$ are
affine; far-left $3$ needs a Lann\'er path of at least six vertices.
This accounts for all spectators as well as the removed own contact.
The first and third roles have
$\cosh L_o=a+(r-2)/\sqrt2>7$.
The isolated $5$ has
$\tanh^2L_a=1-(9-4\sqrt5)/(r-3)$; at $r=7$ its half-tangent
square is $4/\sqrt5-1>3/4$. Each role therefore contradicts
energy at most $1/2$, proving the $D_7$ cost.

For $D_8$, common $4/5$ has exactly $4@4+3@5$.
Indeed deletion $4$ permits simple $3$ at a right $D_4$ leaf,
and deletion $5$ allows one left $3$ or endpoint $4$.
Long-leaf $3$ forces $4@4$ and forbids foreign contacts by the
retained $(3,4,3,3)$ path. A fork $3$ with a left $3$ gives an
affine cycle; without a left contact it is affine $E_8$; a left
$4$ is nonpositive. With no right contact only far $3@1$ survives
the affine tests, but its path to $5$ has six vertices.
Common $5/6$ has the preceding center-neighbor list. Both common
pairs thus force an actual half-tangent square above $3/4$.
At energy at most $4/7$, their opposites are distinct.

Here is the remaining three-direction calculation. Put
$q=\sqrt{4/5}$, $k=\sqrt{5/6}$, $\lambda=(1+k)/2$,
$\alpha=k\sqrt{1-k}/2$, $e=4/7$, and
$P(t)=t/(1+\sqrt{1-t})^2$.
Write $h_i=x\cdot u_i$ and decompose
$x=aw+bv+zt_0+y$, where $w,v$ are the normalized sum and
difference of $u_5,u_6$, $u_4=qu_5+t_0/\sqrt5$, and
$y\perp w,v,t_0$. The full-tangent laws imply
\[
 h_4h_5\ge P(q),\qquad
 T:=(1+\alpha)b^2+\lambda z^2
 \le\delta:=\lambda e-P(k)<1/400.
\]
For the second inequality take the tangent line to
$4p/k-(1+p)^2$ at its first zero $p=P(k)$, with
$p=h_5h_6=\lambda a^2-(1-\lambda)b^2$.
Keeping the negative quadratic terms in $h_4h_5$ gives
\[
 h_4h_5\le q\lambda e+A_0|b|+B_0|z|
              -qkb^2-q\lambda z^2,\qquad
 A_0^2=8/105,\quad B_0^2=4/35.
\]
Weighted Cauchy bounds the linear part by $\sqrt{C_0T}$, where
$C_0=(8/105)/(1+\alpha)+(4/35)/\lambda<47/250$.
Also $qk/(1+\alpha)>.71$, $q>.71$, and $q\lambda e<.489$.
Consequently
\[
 h_4h_5<.489+\sqrt{(47/250)T}-.71T
 <101781/200000<.509<P(q),
\]
a contradiction. The rational seals are
$.8944<q<.8945$, $.9128<k<.9129$,
$.29<\sqrt{1-k}<.2953$, $.13<\alpha<.14$,
and $.95<\lambda<.95645$.
The tight $\delta$ comparison reduces to
$91280000\cdot70000-167780209\cdot38083=26300653>0$;
the last square-root comparison uses
$\sqrt{47/100000}<271/12500$.
This proves $D_8>4/7$ and completes Proposition~\ref{app:prop:ordinary}.

\subsection{The strengthened rank-three rates}
\label{app:strongrankthree}
We prove the improvement \(B_3,H_3>2/7\) used in the forest
accounting. Suppose the component energy is at most \(2/7\).
The adjacent correlations are at least \(1/\sqrt2\), so the
forcing argument above gives one actual opposite for the whole
component, with
\[
 h\leq R^2\min_i v_i^2,\qquad R^2=56/81.
\]
All foreign costs remain included in \(h=q+Q-1\).
Multiple contacts can occur only at the two path endpoints.

For \(B_3\), endpoint contacts \(u,w\) give
\[
 q-R^2v_2^2=(131u^2+106w^2+50\sqrt2uw)/81.
\]
If either label is at least \(4\), its minimum is \(13/9\)
or \(443/324\), both greater than \(1\). Larger contacts
increase the expression. At \(u=w=1/2,Q=0\), the actual
half-tangent squares in directions \(2,3\) are
\((2\sqrt2-1)/7\); increasing \(Q\) increases both lengths.
Their correlation is \(1/\sqrt2\), so the positive two-direction
cover costs at least \((10\sqrt2-12)/7>2/7\).
For one contact the diagonals are \(3,4,2\), with minimum
ratios \(v_i^2/q=1/3,1/2,1/3\). Thus the maximum own costs
are \(243/187,81/53,243/187\).
Nodes \(1,2\) admit only native \(3\), and node \(3\) only
\(3,4\), since \(243/374<\varphi^2/4\).
Every \(q=1\) case is affine; every remaining \(q<1\) case
has \(h<1/4\), forcing \(Q=0\) and \(h=q-1<0\).

For \(H_3\), the endpoint quadratic has positive coefficients;
its minimum at \(u=w=1/2\) is
\((200\varphi+287)/324>1\).
At node \(1\) or \(2\), even contact \(1/2\), tested in
direction \(3\), gives respectively
\[
 (599\varphi+823)/1620>1,\qquad
 (181\varphi+237)/405>1.
\]
At node \(3\), contact at least \(1/\sqrt2\), tested in
direction \(1\), gives \((75\varphi+374)/486>1\).
Only contact \(1/2\) remains, with
\[
 q=(\varphi+2)/4<1,\qquad h\leq14\varphi^4/243.
\]
No foreign component gives \(h<0\); exactly one is impossible
by \(Q<h\). Two or more give
\[
 h\geq q-1/2=\varphi/4>14\varphi^4/243,
\]
the final difference being \((75\varphi-112)/972>0\).
This covers all native/dotted contacts and all spectators.

\subsection{A rank-free joint payment for an isolated \(I\)}
\label{app:Ijoint}
This rank-free argument transcribes \src{F19-local}, \src{A16-local},
and the real scalar barrier isolated in the source ledger.
Let the actual isolated short wall be
\[
 c_0=\sqrt h\,p-\sqrt{1+h}\,e,\qquad0<h<1/2,
\]
with \(e\) the old unit \(I\) normal.
For an actual opposite outside \(I\), write its time as \(A\)
and its contact with \(e\) as \(\gamma\).
The finite white endpoint retains \(e\), so \(\gamma=0\)
or \(1/2\leq\gamma<1\). These opposites are distinct.
For the chosen unit dual \(u\perp e\), its spatial derivative
has square at most \(1+A^2-\gamma^2\).
If \(\gamma=0\), the native gap gives \(A^2\geq1/(4h)\);
if \(\gamma>0\), acuteness gives
\(A^2\geq(1+h)\gamma^2/h\).
Since \(\gamma^2<1\), the derivative comparison yields in
both cases
\begin{equation}\label{app:eq:Ifulltangent}
 \tanh^2L\geq\frac1{1+4h}.
\end{equation}
For the nonzero case, substitute the lower bound on \(A^2\);
the difference from \(1/(1+4h)\), after positive denominators
are cleared, has numerator \(h((5+4h)\gamma^2-1)>0\).

Let \(q\) be the actual \(I\)-edge half-tangent square.
Then \(h=4q/(1-q)^2\), \(0<q<1/5\).
If \(\theta\) is the half-tangent square of the other edge,
the increasing function \(4s/(1+s)^2\) and
\eqref{app:eq:Ifulltangent} show \(\theta>1/5-q\).
The exact cross numerator for the strict comparison is
\[
 (1-q)^2(6-5q)^2-20(1-5q)(1+14q+q^2)
 =(39q-4)^2+q^2(40-10q+25q^2)>0.
\]
Thus the same original vector satisfies
\(e_I+\ip{x}{u}^2>1/5\).
For another whole \(A_1\) this proves \(I+A_1>1/5\);
adding a second ordinary \(A_1>.1\) proves \(I+2A_1>.3\).

For the sharper \(I+T_2\) bound set \(s=\sqrt q\in(0,1/3)\)
and \(t=\sqrt{1+4h}-2\sqrt h\).
Equation~\eqref{app:eq:Ifulltangent} gives normalized height
at least \(t\) in every other actual edge direction, and
\[
 e_I\geq s^2,\qquad(1-s^2)(1-t^2)=8st.
\]
We need \(t+\tfrac32s>.767\).
At \(t_0=.767-1.5s>0\), the left-minus-right of that relation is
\[
 P(s)=.411711-3.835s+9.338289s^2-2.301s^3+2.25s^4.
\]
Putting \(v=s-.2\in(-.2,2/15)\), its coefficients are
\(.00343456,-.3038044,8.497689,-.501,2.25\).
Its cubic term is at least \(-.0668v^2\), so
\[
 P(s)>.0034-.304|v|+8.43v^2>0,
\]
as the last quadratic has negative discriminant.
The original relation decreases strictly in positive \(t\),
proving the barrier. This is a scalar real argument; no
arithmetic chapter of its source is used.

The ordinary \(T_2\) energy is greater than \(1/7\).
If \(s\geq519/2000\), \(s^2+1/7>.21\).
Below that split, \(s^2+(.767-1.5s)^2\) decreases to
\(.2100353125>.21\). Thus \(I+T_2>.21\), counting \(I\) once.

For \(I+T_3>2/7\), \(B_3,H_3\) already have the intrinsic
stronger bound. For \(A_3\), if \(h>1/5\), the actual
\(I\)-edge gives \(e_I>3/70\), and
\(3/70+1/4>2/7\).
If \(h\leq1/5\), all three other half-tangent squares are
at least \(1/5\). The normalized \(A_3\) dual Gram \(K\)
has \(K^{-1}\mathbf1>0\) and
\(\mathbf1^TK^{-1}\mathbf1=5-2\sqrt3\);
positive Cauchy gives energy at least
\((5-2\sqrt3)/5>3/10>2/7\).
The same positive-cover principle will be used below:
if \(K^{-1}\mathbf1>0\) and every height is at least
\(\sqrt\theta\), energy is at least
\(\theta\mathbf1^TK^{-1}\mathbf1\).

\subsection{Propagation from each actual typed rank-two defect}
\label{app:typedpropagation}
The complete propagation proof is the real geometric argument of \src{Old7}
and its old-\(A_2\) companion.
For \(N_7\), put \(b=\cos^2(\pi/7)<5/6\) and
\(h_N=4b/3-1<1/12\).
An opposite in another component retains the whole old
\(A_2\), so has no contact there or one contact of white
cost \(q\in[1/3,1)\).
In the contact case its spatial product with the isolated
wall has square at least \(bq/3\).
Acuteness and the bound on the chosen derivative give
\[
 A^2\geq\frac{bq}{3h_N},\qquad
 \tanh^2L\geq\frac{b}{b+6h_N}
             =\frac{b}{9b-6}>\frac59.
\]
The first ratio is minimized at \(q=1/3\); all denominators
are positive. Without contact the native gap instead gives
\(\tanh^2L\geq1/(1+4h_N)>5/9\).
Every other half-tangent square therefore exceeds \(1/5\).
In particular another \(A_1\) or irreducible \(T_2\) has
energy greater than \(1/5\).

For \(O_7\), put \(d_7=\csc^2(\pi/7)\) and
\(h_O=d_7/4-1\in(0,1/3)\).
No opposite in another component can contact the retained
old \(I_2(7)\): its minimum white cost \(d_7/4>1\) violates
positive definiteness. Its actual product with the isolated
typed wall is thus \(-\sqrt{h_O}A<0\), and the native gap gives
\begin{equation}\label{app:eq:O7height}
 A^2\geq1/(4h_O)>3/4
\end{equation}
for every opposite in the other component.
This immediately gives another \(A_1\) cost \(>2/15\)
by the full-/half-tangent comparison in the main text.

To prove another irreducible \(T_2=I_2(m)\) costs \(>1/5\),
suppose the contrary. With \(k=\cos(\pi/m)\) and
\(d=(1-k^2)^{-1}\), distinct opposites are impossible:
for \(m\geq4\), their tangent product is at most
\(F(\sqrt{1/5})^2=5/9<k\); for \(m=3\), use the pair sum
bound to get \(F(\sqrt{3/20})^2=240/529<1/2\).
There is one common actual opposite, and both actual
edge cosines are at most \(3/2\).
For own contacts \(b,g\geq0\) and entire foreign cost \(Q\),
the finite endpoint equations are
\[
\begin{gathered}
 Q+b^2<1,\quad Q+g^2<1,\quad
 h=d(b^2+2kbg+g^2)+Q-1>3/4,\\
 C_b=\frac{\sqrt d(b+kg)}{\sqrt{1-Q-g^2}},\qquad
 C_g=\frac{\sqrt d(g+kb)}{\sqrt{1-Q-b^2}}.
\end{gathered}
\]
Every nonzero own is consequently a finite native cosine.
There is at least one, and \(Q=0\) or \(Q\geq1/4\).

If only \(b=\beta>0\), the mate has
\(C_g^2=1+h/(1-Q-\beta^2)\).
The bounds \(h>3/4,C_g\leq3/2\) force
\(Q+\beta^2<2/5\), hence \(Q=0,\beta=1/2\).
Then \(h=d/4-1>3/4\) implies \(d>7>d_8\), so \(m\geq9\).
But \(d>81/\pi^2>81/10>31/4\), giving
\(C_g^2=(d-1)/3>9/4\), a contradiction.

If both contacts are positive and \(Q>0\), either their
maximum is at least \(1/\sqrt2\), giving
\(C_b^2\geq2+4\sqrt2/3>9/4\), or both are \(1/2\).
For old \(m=3\), the latter gives \(h=Q<3/4\);
for \(m\geq4\), it gives
\(C_b^2\geq(1+k)/(2(1-k))>5/2\).
If both contacts are positive and \(Q=0\), a contact of
label at least \(5\) gives \(C_b^2>3\); two \(4\)'s give
\(C_b^2=(1+k)/(1-k)\geq3\).
For \(3/4\), old \(m\geq4\) gives \(C_b\geq2\),
while old \(m=3\) gives \(h=\sqrt2/3<3/4\).
For \(3/3\), old \(m\geq5\) gives
\(C_b^2=(1+k)/(3(1-k))>3\);
old \(m=3,4\) give \(h=0,1/\sqrt2<3/4\).
This exhausts all finite labels and all foreign costs,
contradicting the common-wall low-energy branch.
Thus \(O_7+T_2\) has its other component energy \(>1/5\).


\section{The uniform actual \(B/D\) component floor}
\label{app:BDfloor}
The next proof supplies the \(9/14\) input used before the long-prefix
argument. Its small boundary cases refer only to the complete ordinary
contact calculations in Appendix~\ref{app:ordinary}. All displacements
remain in the indicated component, and all positive actual-root
constructions retain every original spectator wall.

\begin{proposition}[Actual $B/D$ floor]\label{app:prop:BDnine}
Every whole $B_r$ or $D_r$ component with $r\ge9$ has actual
projection energy greater than $9/14$, with arbitrary finite
spectators.
\end{proposition}

\begin{proof}
Assume energy at most $e=9/14$. In $B$ it suffices to bound the first
seven coordinates; all displacements below lie there. The normalized
directions in either type are $u_i=(e_1+\cdots+e_i)/\sqrt i$.
We first establish distinctness of $4/5,6/7$, then use those laws to
exclude all common $1/2,2/3$ profiles. This order is noncircular.

For an internal common $i/(i+1)$ with right retained rank at least
four, finite classification permits a right long-leaf $3$, or none.
The first forces own $4$: own $0/3$ is affine fork/triangle, and
any foreign contact retains the nonpositive $(3,4,3,3)$ path.
Without a right contact, internal-left $3$ and endpoint $4$ are
affine. Only far $3@1$ remains, but its path to $i+1$ has at least
six vertices, contradicting the Lann\'er bound. Right $D$-fork
contacts give affine $E_8$ or a cycle with the left contact.
Thus the sole common profile is $4@i+3@(i+1)$, $Q=0$, with
$\cosh L_{i+1}=i+(i+1)/\sqrt2>7$.
The $D_8$ right-$D_4$ and center-neighbor boundary calculations are
in Section~\ref{app:ordinary}; the center-neighbor calculation also
handles $D_9$ at $6/7$.
In $B_9$, the extra right short-end $3@9$ gives $F_4$, excluding
left and foreign white contacts. The removed own $6$ remains to check:
own $0$ is affine $F_4$, own $4$ the affine $(4,3,3)$ star, and own
$3$ gives $\cosh L_6=(1+1/\sqrt2)\sqrt{12}>11/2$ and
$\tau_6^2>9/13>e$. Consequently $4/5$ and $6/7$ are distinct.

Write $h_i=x\cdot u_i$. Distinct $6/7$ gives
$h_6h_7\ge k/(1+\sqrt{1-k})^2>9/16$, where
$k=\sqrt{6/7}>576/625$ with squared margin $21318$.
Hence $x\cdot w>3/2$ for $w=u_6+u_7$.
We record four local tests with their exact margins.

\emph{Edge $2$ with $\cosh L_2\ge1+\sqrt2$.}
Then $h_2^2\ge\sqrt2-1$. Projection of $w$ about $u_2$ has
$\rho=1/\sqrt3+\sqrt{2/7}<1.114$ and residual square
$S=29/21+4\sqrt{2/21}<2.62$.
Its Cauchy envelope $\rho t+\sqrt{S(e-t^2)}$ decreases for
$t\ge q=\sqrt{\sqrt2-1}$: its derivative condition follows from
$\rho^2e<13/16<6/5<\|w\|^2q^2$.
Since $q<.644$, $e-q^2<.23$, and $\sqrt{2.62\cdot.23}<.777$,
$x\cdot w<1.114(.644)+.777=1.494416<3/2$.
These follow by squaring, using $1.414<\sqrt2<1.4143$,
$1/\sqrt3<.578$, $\sqrt{2/7}<.535$, $\sqrt{2/21}<.309$.

\emph{Isolated $5@1$.}
Its height square is $(\ph-1)/2$ and its actual direction-$4$
shadow has $\tau_4^2=\ph-1$. For the same $w$, projection about
$u_4$ has $\rho=2(1/\sqrt6+1/\sqrt7)<787/500$ and
$S=16/21+4/\sqrt{42}<69/50$. Use
$1/\sqrt6<.409$, $1/\sqrt7<.378$, $\sqrt{42}>162/25$;
the last residual comparison has cleared margin $23$.
Since $\sqrt e<401/500$, $e-(\ph-1)<1/40$ and
$(93/500)^2>69/2000$, Cauchy gives
$x\cdot w<362087/250000<3/2$.

\emph{Isolated $5@2$.}
Its height square is $\ph$. The direction-$1$ shadow gives
$x_1>31/50$, since $F(31/50)^2<161/200<\ph/2$ with cleared
margin $6541881$. Its actual edge $3$ has
$\cosh L_3=\sqrt2\ph^2>37/10$; hence
$S_3=x_1+x_2+x_3>131/100$, using
$3\tau_3^2>81/47>(131/100)^2$, of margin $3433$.
Distinct $4/5$ gives $x\cdot w_0>57/40$, $w_0=u_4+u_5$.
Put $v=e_1+e_2+e_3$, $k_0=1/2+1/\sqrt5<237/250$, and
$x_0=(31/50,69/200,69/200,0,\ldots)$.
Then $x_0=(11/40)e_1+(69/200)v$, $x\cdot x_0>x_0^2>31/50$.
The residual $w_0-k_0v$ has square $13/20+1/\sqrt5<11/10$
and excess demand $>9/50$.
All coefficients of $x_0+(w_0-k_0v)/6$ remain positive, and its
Cauchy margin is at least
\[
 (31/50+(9/50)/6)^2-e(31/50+(11/10)/36)=3/700>0.
\]

\emph{Contacts $4@2+3@3$.}
The white $A_2$ cost is $2/3$, so $Q=0$ or $1/4$; the latter has
$\cosh L_3=4+3\sqrt2>7$. At $Q=0$, put $s=\sqrt2$.
The actual floors satisfy
$\tau_2^2=(1+2s)/7$, $\tau_3^2=(2s-1)/3$,
$\tau_2\tau_3=1/\sqrt3$.
Their positive equality vector has $E_0=(6s-4)/7>.6407$.
The positive old-$A_2$ correction of $c$ is an actual future point
of square $-s/3$, giving pure displacement
\[
 U=\frac{2+s}{2}e_1+\frac{1+s}{3}(e_2+e_3+e_4),\qquad
 x\cdot U\ge A-\sqrt{s/3},\quad A^2=(5+4s)/2.
\]
Its projection coefficients on $u_2,u_3$ are
$\rho_2=(2s+1)/6$, $\rho_3=(1+s)/\sqrt3$;
its residual square is $(7+4s)/12<1.055$, with excess demand
$\Delta>.06$.
Indeed $1.4142<s<1.4143$ gives $A>2.3083$,
$\sqrt{s/3}<.6867$, $\tau_2\in(.7395,.7396)$,
$\tau_3\in(.7806,.7808)$, $\rho_2<.6381$, $\rho_3<1.394$,
and excess $>.06122604$.
The old positive coefficients exceed $.3,.52$; subtracting
$(\rho_2,\rho_3)/16$ removes less than $.04,.088$.
The corrected dual therefore has Cauchy margin at least
\[
 (.6407+.06/16)^2-e(.6407+1.055/256)
       =8825113/11200000000>0.
\]
This is an actual-point constraint, not an edge-only assertion.

Now exhaust common $1/2$. Deletion $1$ retains right rank at least
eight, so the right contact is none or long-leaf $3@2$.
The latter forces own $\gamma\ge1/\sqrt2$ (own $0/3$ is affine),
and $\cosh L_2=(1+\gamma)/\sqrt{1-\gamma^2-Q}\ge1+\sqrt2$,
already excluded. Without a right contact,
$A^2=2\gamma^2+Q-1>0$ and $\gamma^2+Q<1$.
Own $0$ has zero derivative; own $3$ has height below $1/4$ and
forces a foreign closure; own $4/6$ is affine.
Own $5$ with $Q>0$ has $\cosh L_2\ge\ph^2>1+\sqrt2$;
at $Q=0$ use the isolated-$5$ test. Finite $m\ge7$ forces $Q=0$;
$m\ge8$ uses the edge-$2$ test, leaving only isolated $7$.

For common $2/3$, deletion $2$ retains right $B/D$ rank at least
seven; deletion $3$ allows one left $3/4/5$ at $1$ or $2$.
A right long $3@3$ forces own $3/4/5$: own $0$ is affine, own $3$
a triangle, own $4$ is the tested $4+3$ profile, and own $5$ has
$Q=0$, $\cosh L_3=(2+5\ph)/\sqrt2>7$.
Without a right contact the full left list is
$3@1,4@1,5@1,3@2,4@2,5@2$.
The first has $A^2=Q-1/2<1/6$ and forces foreign closure;
$4@1,3@2,4@2$ are respectively affine terminal-$4$, two-fork,
and $(4,3,3)$ star. The $5$ cases have $Q=0$ by white cost
$\ph^2/3>3/4$ and were tested above.
For $D_9$ one extra right-fork $3$ gives retained $D_7+c=E_8$.
It excludes old-$1$ and foreign contacts, but the removed own $2$
must still be checked: own $0/3/4$ gives affine $E_8$/cycle/star,
and own $5$ retains nonpositive $(3,5,3)$. At higher rank the fork
attachment already contains affine $E_8$. No finite branch is omitted.

For the final isolated $7@1$, put $t=\cos(2\pi/7)$,
$g=\cos(\pi/7)$, and
\[
 \xi=c+g^{-1}n_1+(1-t)g^{-1}n_2
     =\sqrt t\,p-\frac V{\sqrt{1+t}},\qquad
 V=t(e_1+e_2)+(1-t)e_3.
\]
It has $\xi\cdot c=\xi\cdot n_1=0$,
$\xi\cdot n_2=(1/2-t)/g<0$ and
$\xi^2=-(1-t)(2t-1)/(1+t)<0$.
All coefficients and its time coefficient are positive, proving
actual feasibility and future orientation. Thus
$x\cdot V\ge\sqrt{t(1+t)}-\sqrt{(1-t)(2t-1)}$.
The increasing cubic $8t^3+4t^2-4t-1$ gives $.6234<t<.6235$,
with endpoint values $-.000926392768,.000105223$.
Hence the first radical is $>1.0059$, the second $<.305$,
and $.95<\|V\|<.9588<.959$, so $x\cdot\widehat V>.73$.
The projection of $w$ onto $\widehat V$ is
$\rho=(1/\sqrt6+1/\sqrt7)(1+t)/\|V\|\in(1.3,1.35)$,
and residual square is $<11/5$.
The Cauchy envelope decreases above $.73$, since
$\rho^2e<6561/5600<3(.73)^2<\|w\|^2(.73)^2$.
It gives $x\cdot w<1.35(.73)+.492=591/400<3/2$;
the final squared radical margin is $277/1750000$.
Thus $1/2,2/3,6/7$ all have distinct actual opposites.

Their exact pair sums are
\[
 x\cdot(u_1+u_2)>1.091,\quad
 x\cdot(u_2+u_3)>1.265,\quad
 x\cdot(u_6+u_7)>1.512.
\]
Set $W=.1(u_1+u_2)+.055(u_2+u_3)+.305(u_6+u_7)$. Direct positive
Gram estimates give
\[
\begin{gathered}
 x\cdot W>.639835,\qquad W^2<.636435825,\\
 .639835^2-\frac9{14}(.636435825)
 =\frac{70423123}{280000000000}>0.
\end{gathered}
\]
To verify the pair sums, use $F(a)F(b)\le F((a+b)/2)^2$.
The correlation lower bounds
\[
7071/10000,\qquad81649/100000,\qquad9258/10000
\]
have squared margins $959,\allowbreak322397,\allowbreak26052$ below the exact values.
At sums $m/1000$, $m=1091,\allowbreak1265,\allowbreak1512$, the cleared differences
\[
p(4000000+m^2)^2-16dm^2\,1000000
\]
are respectively
\[
40828209713231,\qquad358399613480625,\qquad52443946917888.
\]
The six upper Gram entries of the three pair vectors are
$3.415,3.633,3.852,3.102,1.899,2.474$, obtained by positive
squaring of their prefix correlations. Their expansion is the stated
bound on $W^2$. This contradicts Cauchy and proves the proposition.
\end{proof}

\section{Coupled actual pressures in an \(A\)-chain}\label{app:Achain}

This shared proof consolidates \src{Ahalf}, \src{Afour}, and
\src{A13}; the later central-window proofs consolidate the geometric
window of \src{A16}, \src{A16-new}, \src{A17}, and \src{A18}.

The crucial feature of this section is that a foot at one end is
tested against \emph{actual} opposites at the other end. A collection
of separate edge estimates would lose that coupling.

\begin{theorem}[Uniform actual \(A\) pressure]\label{app:thm:Ahalf}
With arbitrary finite spectators,
\[
 \|x_{A_r}\|^2>\frac12\qquad(r\ge7).
\]
\end{theorem}

We first prove the uniform step for rank at least eight, then give the
rank-seven base in Subsection~\ref{app:A7base}. Put \(N=r+1\), and realize
\[
 n_i=(e_i-e_{i+1})/\sqrt2,\quad
 C_{ij}=\frac{2\min(i,j)(N-\max(i,j))}{N},\quad
 h_i=x_A\cdot u_i .
\]
The root space is the sum-zero subspace of \(\R^N\). All its
inequalities below use the actual projection \(x_A\).

\subsection{A forcing window and the complete common-pair inventory}

Suppose first that \(N\ge9\) and \(\|x_A\|^2\le1/2\). Consider the
four pairs \(1/2,3/4,(r-3)/(r-2),(r-1)/r\).
If their respective opposites are distinct, the outer pair sums
exceed \(1\), and the inner sums exceed \(19/16\). These follow
from \eqref{eq:pair-shadow}, concavity of \(\log F\), and the minimum
correlations \(\sqrt{3/7}\), \(\sqrt{3/5}\) at \(N=8\).
For the inner comparison,
\[
 F(19/32)^2=\frac{1478656}{1918225}
 <\frac{387}{500}<\sqrt{3/5}.
\]
For the outer one, \(F(1/2)^2=16/25<\sqrt{3/7}\).
Set
\[
 L=2u_1+2u_2+u_3+u_4,\qquad
 R=2u_r+2u_{r-1}+u_{r-2}+u_{r-3}.
\]
The correlations on the left increase to \(\sqrt{i/j}\), whence
\[
 L^2,R^2<
 12+6\sqrt2+4/\sqrt3+4\sqrt{2/3}+\sqrt3<28.
\]
The last bound follows, for example, by the positive rational
upper bound \(2351/84\). For \(a=(2,2,1,1)\), direct coordinates give
\[
 L\cdot R=\left(\sum_{i=1}^4a_i\sqrt{\frac{i}{N-i}}\right)^2.
\]
This decreases in \(N\); at \(N=9\) its square root is
\(\sqrt2+2\sqrt{2/7}+2/\sqrt5<17/5\).
Thus \((L+R)^2<80\), whereas distinctness of all four pairs gives
\(x_A\cdot(L+R)>51/8\). The Cauchy gap is \(41/64>0\).
After reversing the numbering, a common pair is therefore \(1/2\)
or \(3/4\). This is a relabeling, not a symmetry assumption on \(P\).

Here is the finite inventory used for the target in
Theorem~\ref{app:thm:Ahalf}. It is useful to state it before making any
energy-dependent discard. Write \(Q\) for the total spectator cost.
Every retained finite diagram is a forest, so there is at most one
contact on each retained own path. The foreign costs satisfy
\eqref{eq:foreigncost}, are at least \(1/4\), and have the gap
\((1/4,1/3)\).

For a common \(3/4\) opposite, its own contact positions are in
\(\{1,2,3\}\) and \(\{4,\ldots,r\}\). A double profile away from
positions \(3,4\) survives one deletion as a connected rank-\(r\)
finite diagram. Its two new bonds must be \(3\): a weighted bond
would be internal in a \(B\) path, and \(D,E_8\) are simply laced.
The two \(3\)'s instead give an affine cycle in the full diagram.
This also covers the rank-eight \(E_8\) boundary.
At positions \(3,4\), the labels are
\[
 p\in\{3,4,5\},\qquad q\in\{3,4\}.
\]
The pair \(3,3\) is affine. The two smallest remaining assignments
have actual endpoint cosines
\begin{equation}\label{app:eq:Acentral}
 \sqrt{\frac3{2N}}\bigl(N-3+\sqrt2(N-4)\bigr),\qquad
 (3\sqrt2+4)\sqrt{\frac{N-4}{2N}},
\end{equation}
both greater than \(3\) at \(N=9\) and increasing thereafter.
Increasing a contact or \(Q\) strengthens the appropriate endpoint.

A single \(3\) at positions \(3,\ldots,r-2\) contains an affine
\(E_8\) tree with arms \(1,2,5\), or an affine \(E_7\) tree with
arms \(1,3,3\), obtained by truncation if needed. The other
single-\(3\) possibilities satisfy
\[
\begin{array}{c|c|c}
 \text{position}&Q\text{ bound}&A^2\\ \hline
 1&Q<5/8&Q-(N+1)/(2N)<1/8\\
 2&Q<1/2&Q-2/N\\
 r-1&Q<2/(N-3)&Q-2/N\\
 r&Q<1/2+1/(2(N-3))&
          Q-(N+1)/(2N)<3/(2N(N-3)).
\end{array}
\]
The end heights are below \(1/4\); the middle rows allow at most
one foreign component and its cost exceeds \(A^2\). All contradict
\eqref{eq:foreigncost}. A single \(4\) is an affine \(3,3,4\) star
internally; at either end its white condition gives \(Q<1/4\),
so \(Q=0\) and its height is negative. A single \(5\) can occur
only at \(1\) or \(3\), has \(Q=0\) from its \(H_4\) cost
\(3\ph^2/8>3/4\), and at \(3\) closes direction \(4\) with
\begin{equation}\label{app:eq:Aiso5central}
 \cosh L_4=3\ph^3\sqrt{(N-4)/N}.
\end{equation}
At \(1\) it gives foot F below. A zero own contact cannot close
an own edge. This completes the common \(3/4\) list.

For a common \(1/2\) opposite there is possibly a contact at \(1\)
and at most one on \(2,\ldots,r\). For double positions \(1,j\),
\(j\ge3\), deletion \(2\) gives a full connected finite path or
branch. A branch has only \(3\)'s; a path can have a \(4\) only
at its terminal bond \(1-c\), with \(j=3\) or \(r\). All \(3\)'s
give an affine cycle. The profile \(4@1+3@r\) contains the affine
\(F_4\) path \(2,1,c,r,r-1\). The remaining \(4@1+3@3\) has
\(Q<1/(2(r-1))\), hence \(Q=0\), and gives foot H.
Adjacent contacts \(1,2\) have an arbitrary finite native first
label and second label \(3\) or \(4\); \(3,3\) is affine.

For single \(3\), the same internal affine exclusions apply.
At \(1,r,r-1\), respectively, the bounds are
\[
 Q<3/4,\quad Q<1/2+1/(2r),\quad Q<2/r.
\]
The first two give height below \(1/4\), and the last leaves no
foreign cost and negative height. At \(2\),
\(Q<1/2+1/(2r)\) allows at most two foreign components.
One exceeds \(A^2=Q-2/N\). With two, both contacts are \(3\);
\(c,n_2,n_1,n_3\) and their two foreign neighbors give the exact
affine two-fork diagram. A single \(4\) is affine internally;
at \(1\), \(Q<1/2\) allows at most one cost greater than
\(A^2=Q-1/N\); at \(r\), \(Q<1/r\) leaves negative height.
A single label at least \(5\) is possible only at \(1\).
Label \(6\) gives the affine \(3,6\) path. Labels at least \(7\)
have \(Q<1-\beta_1^2<1/4\), hence \(Q=0\).
Label \(5\) with \(Q=0\) gives F; with \(Q>0\) it remains in L.
Thus the remaining end profiles are
\[
 \mathrm L:\ \text{adjacent, not }(3,3);
 \ \text{isolated native }\ge7;
 \ \text{isolated }5\text{ with }Q>0,
 \qquad \mathrm H,\mathrm F .
\]
No foreign branch has been suppressed by assuming a full star finite.

\subsection{Three feasible feet prove the half bound}

For H and F take the following positive actual combinations:
\[
\begin{array}{c|c|c}
 &\xi&k=-\xi^2\\ \hline
 \mathrm H&
 c+(\sqrt2-\tfrac12)n_1+(\sqrt2-1)n_2+
                         \tfrac{\sqrt2}2n_3&
 \sqrt2-\tfrac54\\
 \mathrm F&
 c+\displaystyle\sum_{j=1}^4((j+1)\ph-2j)n_j&
 34\ph-55 .
\end{array}
\]
For H the products with \(c,n_1,n_3\) vanish and the \(n_2\)
product is \((\sqrt2-3)/4<0\). For F the products with
\(c,n_1,n_2,n_3\) vanish and the \(n_4\) product is \(3\ph-5<0\).
Every other actual inequality follows from positivity. Both have
\(Q=0\), future coefficient \(A>0\), and pure \(A\)-spatial part,
with
\[
 A_{\mathrm H}^2=\frac32+\sqrt2-\frac{11/2+3\sqrt2}{N},
 \qquad
 A_{\mathrm F}^2=\frac{\ph^2(N-1)}{2N}-1.
\]
At \(N=9\),
\[
 A_{\mathrm H}^2-9k_{\mathrm H}
   =\frac{437-300\sqrt2}{36}>0,\qquad
 A_{\mathrm F}^2-9k_{\mathrm F}
   =\frac{50(89-55\ph)}9>0.
\]
The time norms increase with \(N\). Equation~\eqref{eq:pointenergy}
therefore excludes these profiles under energy at most \(1/2\).
The central doubles and \eqref{app:eq:Aiso5central} already have
\(\cosh L>3\), so are excluded by \eqref{eq:halfheight}.

For L, first clear all foreign components positively to obtain
\(c_F=A p-v_A\), \(c_F^2=1-Q\). Put
\[
 b=\beta_1,\ d=\beta_2,\ B=b+2d,\quad
 h=\frac43(b^2+bd+d^2)+Q-1,\quad z_0=h/B^2,
\]
\[
 \lambda_1=(4b+2d)/3-h/B,\qquad
 \lambda_2=2B/3-2h/B,\qquad
 \xi=c_F+\lambda_1n_1+\lambda_2n_2=A p-U .
\]
The endpoint condition gives
\[
 B\lambda_1=2bd+1-Q>0,\quad
 B\lambda_2=2(1-b^2-Q)>0,
\]
and direct multiplication gives
\[
 \xi\cdot c=\xi\cdot n_1=\xi\cdot S=0,\quad
 \xi\cdot n_2=-3h/(2B)<0,\quad
 k=-\xi^2=B^2(z_0-3z_0^2).
\]
Here \(S\) denotes all cleared spectator roots. The complete L
range is
\begin{equation}\label{app:eq:Ldomain}
 \frac8{81}<z_0<\frac13.
\end{equation}
Indeed adjacent contacts with one label at least \(4\) give
\(z_0>1/10\); isolated labels at least \(7\) have \(b>9/10\),
so \(z_0=4/3-b^{-2}>8/81\); isolated \(5\) with \(Q\ge1/4\)
gives \(z_0\ge3\ph-14/3>2/15\).
The upper bound is precisely \(1-b^2-Q>0\).
Thus \(\xi\) is actual, future timelike and has no foreign spatial part.

For \(a_N=2/3-2/N\), its full vector data are
\begin{equation}\label{app:eq:Lvector}
 A^2=B^2(a_N+z_0),\quad
 U^2=B^2(a_N+3z_0^2),\quad
 U_{\rm raw}=\frac B{\sqrt2}(2/3+z_0,2/3+z_0,2/3-2z_0,0,\ldots,0).
\end{equation}
Subtracting \(\sqrt2 B/N\) from every coordinate centers the raw
vector. Its normalized actual demand
\(g=(A-\sqrt k)/\sqrt{U^2}\) satisfies \(g^2>9/20\).
At \(a_N=4/9\) the equivalent quadratic is
\(2523z_0^2-720z_0+484/9>0\), of negative discriminant
\(-72976/3\); increasing \(N\) strengthens it.

Let \(W=u_{r-1}+u_r\). Direct coordinates give
\begin{equation}\label{app:eq:rho}
 \rho=\frac{U\cdot W}{\sqrt{U^2}}
 =\frac{2/\sqrt{N(N-2)}+\sqrt{2/(N(N-1))}}
        {\sqrt{a_N+3z_0^2}} .
\end{equation}
The numerator decreases and denominator increases with \(N\).
At \(N=9,z_0>8/81\), its squared upper seal is
\(48114/129500<3721/10000\), so \(\rho<.61\).
If the far opposites are distinct, \(x_A\cdot W>1\).
Perpendicular Cauchy instead gives
\[
 x_A\cdot W<\frac{.61}{\sqrt2}
              +\sqrt{(1/2-9/20)4}<.89<1.
\]
If they are common, their inventory supplies another L foot:
H and F have already been excluded. Normal reversal changes
the raw sum from \(\sqrt2 B\) to \(-\sqrt2 B'\).
The disjoint raw supports therefore give the \emph{positive}
centered product
\[
 U\cdot U'=2BB'/N,\qquad
 \operatorname{corr}(U,U')<\frac{2/N}{a_N}
                 =\frac3{N-3}\le\frac12.
\]
The two actual demands force energy greater than
\(4(9/20)/3=3/5\). This proves the half bound.


\subsection{The rank-seven base}\label{app:A7base}

Set $N=8$ and suppose $\|x_A\|^2\le1/2$. The correlations of
pairs $1/2,6/7$ are $\sqrt{3/7}$ and those of $3/4,4/5$ are
$\sqrt{3/5}$. The pair-sum calibrations already proved therefore
apply. If all four pairs have distinct opposites, put
\[
 w=2u_1+2u_2+u_3+2u_4+u_5+2u_6+2u_7.
\]
Then $x_A\cdot w>51/8$, whereas direct Gram summation gives
\begin{align*}
 w^2={}&\frac{64}7+\frac{32}3+\frac{16}5+4
 \\
 &+\frac{64}{\sqrt{21}}+\frac{64}{\sqrt{105}}
 +\frac{16}{\sqrt7}+\frac{64}{3\sqrt5}\\
 &+\frac{16}{\sqrt3}+8\sqrt{3/5}<\frac{1208}{15}.
\end{align*}
The first four terms sum to less than 28, and the last six are
respectively less than $14$, $32/5$, $32/5$, $10$, $28/3$, $32/5$; each comparison
follows by squaring positive rational bounds. Cauchy is impossible
since
\[
(51/8)^2-(1208/15)/2=359/960>0.
\]
Reversing the numbering
if necessary, a common pair is $1/2$ or $3/4$.

\paragraph{The central pair.}
For a common $3/4$ opposite there is at most one contact on
each of the paths $1,2,3$ and $4,5,6,7$. Native bonds have labels
$3,4,5$: a larger label with a retained old bond is not finite.
For two contacts outside positions $3,4$, the following finite
white tests exhaust the position pairs:
\[
\begin{array}{c|l}
 (1,4),(1,7),(2,4),(2,7),(3,5),(3,7)&
       \text{a rank-seven path after deleting }3\text{ or }4\\
 (1,5),(1,6),(2,5),(2,6)&D_7\text{ or }E_7\text{ after deleting }3\\
 (3,6)&D_7\text{ after deleting }4 .
\end{array}
\]
Both new bonds must be 3 in these rows; their cycle in the full
diagram is affine. At positions $3,4$, the $3,3$ pair is affine;
every other assignment gives an endpoint cosine greater than 3,
already for the minimal $3,4$ and $4,3$ assignments in
\eqref{app:eq:Acentral} at $N=8$. It therefore forces energy greater
than $1/2$.

For a single 3-contact, the complete own costs and sufficient
white bounds are
\[
\begin{array}{c|rrrrrrr}
 j&1&2&3&4&5&6&7\\\hline
 \beta^t C\beta&7/16&3/4&15/16&1&15/16&3/4&7/16\\
 Q<&5/8&1/2&5/8&3/5&2/5&2/5&3/5 .
\end{array}
\]
Here $A^2=\beta^t C\beta+Q-1$. One foreign component is
impossible by $q<A^2$ in~\eqref{eq:foreigncost}. Every row except
$j=3,4$ allows at most one, or has $A^2<1/4$. In rows $3,4$,
two foreign components must both start with a 3-contact; the
new wall, the contacted old node, its two old neighbors and the
two foreign neighbors give the exact affine two-fork $D_5$
diagram. The case $Q=0$ has nonpositive height.

A single 4 is an affine $3,3,4$ star internally. At an end
$Q<1/4$, so $Q=0$ and the height is negative. A single 5 is
possible only at positions 1 or 3, and its retained $H_4$ cost
$3\ph^2/8>3/4$ forces $Q=0$. Position 3 gives direction-four
cosine greater than 3. Position 1 gives foot F above. At $N=8$
its time norm satisfies
\[
 A_F^2=(7\ph-9)/16,\qquad
 A_F^2-9k_F=(7911-4889\ph)/16>0.
\]
For example $\ph<16181/10000$, proved by positive squaring,
gives $7911-4889\ph>1091/10000$. Thus the actual-point
inequality~\eqref{eq:pointenergy} excludes this last row.
The central common-pair case is complete.

\paragraph{The end pair.}
For common $1/2$, a double profile has positions $1,j$.
For $j=2$ its second label is 3 or 4 and its first label is any
finite native label. At $j=3,7$ the labels are $3$ at $j$ and
$3$ or $4$ at 1. At $j=4,6$, deletion 2 gives $E_7$ and
both labels are 3. At $j=5$, deletion 2 already contains affine
$E_6$. Two 3-labels always give an affine cycle. The
$4@1+3@7$ row contains the affine $F_4$ path
$2,1,c,7,6$; the $4@1+3@3$ row has $Q=0$ and is foot H.
The displayed H formula gives
$A_H^2-9k_H=(193-134\sqrt2)/16>0$ at $N=8$.

For a single 3, the own costs are the same seven entries above,
while the sufficient $Q$ bounds at positions $1,\ldots,7$ are
\[
 3/4,\quad4/7,\quad2/7,\quad1/7,\quad1/7,\quad2/7,\quad4/7.
\]
All except position 2 have $A^2<1/4$. At position 2, one foreign
component contradicts $q<A^2$; with two, each exceeds $1/4$,
so each is at least $1/3$, contrary to $Q<4/7$. A single 4 is
affine internally. At 1, $Q<1/2$ permits at most one foreign
component, whose cost exceeds $A^2=Q-1/8$; at 7, $Q<1/7$
leaves none and gives negative height.

A single label at least 5 is possible only at 1. Label 6
contains the affine $3,6$ path. A label at least 7 forces
$Q<1-\beta_1^2<1/4$, hence $Q=0$. Label 5 with $Q=0$ was
excluded by F. Thus the only remaining profiles are exactly L:
adjacent contacts other than $3,3$, isolated label at least 7
with $Q=0$, or isolated 5 with $Q>0$.

\paragraph{The far-end comparison.}
The positive spectator clearing and actual L-foot formulas
\eqref{app:eq:Ldomain}--\eqref{app:eq:Lvector} remain valid at $N=8$,
now with $a_N=5/12$. The lower bound $z_0>8/81$ follows as
before: adjacent contacts give $z_0>1/10$, an isolated label
at least 7 has $\beta_1>9/10$, and isolated 5 with $Q>0$
has $z_0\ge3\ph-14/3>2/15$. For the label-seven bound,
$8t^3-4t^2-4t+1$ is $-1/125$ at $9/10$ and strictly
increasing from $\sqrt3/2$ to 1.

The actual demand now satisfies $g^2>11/25$: equivalently
$A^2/k>324/49$, whose cleared difference is
\[
 972z_0^2-275z_0+245/12>0.
\]
Its discriminant is $-3755$. With $W=u_6+u_7$, direct
coordinates give
\[
 0<\rho^2=
 \frac{5/42+1/(2\sqrt{21})}{5/12+3z_0^2}
 <\frac{603612}{1170300}<\frac{324}{625},
 \qquad W^2=2+2\sqrt{3/7}<10/3.
\]
Here the numerator is less than $23/100$ and
$z_0^2>64/6561$ gives the stated rational seal. If the far
opposites are distinct, their pair law gives $x_A\cdot W>1$.
Perpendicular Cauchy instead gives
\[
 x_A\cdot W<\frac{18}{25\sqrt2}
       +\sqrt{(1/2-11/25)10/3}
 =\frac{18}{25\sqrt2}+\frac1{\sqrt5}<\frac{24}{25}.
\]
If they are common, the reversed end inventory supplies another
L foot. Normal reversal is $e_j\mapsto-e_{9-j}$, so its raw
sum is $-\sqrt2B'$ on the disjoint final three coordinates.
Consequently $U\cdot U'=BB'/4>0$, and their normalized
correlation is less than $3/5$. The sum of these two normalized
directions has square less than $16/5$ and actual scalar product
greater than $2\sqrt{11/25}$. Cauchy forces
\[
 \|x_A\|^2>\frac{44/25}{16/5}=\frac{11}{20}>\frac12.
\]
Both far-end alternatives are impossible. This proves the
rank-seven base and completes Theorem~\ref{app:thm:Ahalf}.


\subsection{The paired-end Gram and the uniform \(4/7\) bound}
\label{app:Afour}
\begin{proposition}
Every actual \(A_r\) component, \(r\geq8\), has energy
greater than \(4/7\), with arbitrary finite spectators.
\end{proposition}
\begin{proof}
Put \(N=r+1\geq9\) and suppose energy is at most \(4/7\).
For distinct outer \(1/2\) pairs and inner \(3/4\) pairs,
the height sums exceed \(41/40\) and \(1219/1000\),
and likewise at the opposite end. At \(N=9\) the
correlations are \(\sqrt7/4,\sqrt{5/8}\); the exact tests are
\[
 F(41/80)^2=\frac{43033600}{65302561}
 <33/50<\sqrt7/4,\qquad
 F(1219/2000)^2=\frac{23775376000000}{30095768093521}
 <79/100<\sqrt{5/8}.
\]
The cross-product margins are \(3304513,28079388159\).
The correlations increase with \(N\).
For \(M_i=u_i+u_{N-i}\), \(1\leq i\leq j\leq4\),
\[
 \ip{M_i}{M_j}=2N\sqrt{\frac{i}{j(N-i)(N-j)}}.
\]
Its logarithmic derivative is
\(1/N-1/(2(N-i))-1/(2(N-j))<0\).
Thus every paired-end Gram entry decreases with \(N\).
For \(W=2M_1+2M_2+M_3+M_4\),
\[
 W_9^2=\frac{1116}{35}+\frac{60}{\sqrt7}
       +\frac{36}{\sqrt{10}}+\frac{72}{\sqrt{70}}
 <\frac{747}{10}.
\]
The last three terms are \(<227/10,57/5,87/10\);
these imply the sharper \(2614/35<747/10\).
If all four selected pairs were distinct, then
\[
 \ip{x_A}{W}>3269/500,\qquad
 (3269/500)^2-(4/7)(747/10)=104527/1750000>0,
\]
contradicting Cauchy. Reversing the numbering, a common
\(1/2\) or \(3/4\) therefore occurs.

Use the complete common-pair inventory above before its
half-energy discards. At \(4/7\), the individual edge
cutoff is \(C=11/3\). The two minimal central-double
cosines in \eqref{app:eq:Acentral} exceed \(11/3\)
already at \(N=9\), and increase thereafter.
The single central \(5\) cosine in
\eqref{app:eq:Aiso5central} also exceeds \(11/3\).
The remaining left profiles are precisely L, H, F.
The same finite-white, affine and full-spectator checks
apply; no earlier small-target exclusion of H/F is used.

For \(\xi=A p-U\), write
\(g^2=(A-\sqrt k)/(A+\sqrt k)\) and \(d=U^2/B^2\),
where \(\sqrt2B>0\) is its raw coordinate sum.
At \(N=9\) the exact shared feet satisfy
\[
\begin{array}{c|c|c|c}
\text{type}&g^2>&d>&
 \rho=\ip{U/\norm U}{u_{r-1}+u_r}<\\ \hline
\mathrm L&9/20&4/9&.61\\
\mathrm H&53/100&1/3&3/4\\
\mathrm F&11/20&2/9&9/10.
\end{array}
\]
For L, its formulas and negative-discriminant seal above
prove the first two entries. The correlation numerator
has square
\[
 (2/\sqrt{63}+1/6)^2
 =23/252+2/(3\sqrt{63})<22/125,
\]
giving \(\rho^2<48114/129500<3721/10000\).
For H, \(A^2=(8+6\sqrt2)/9>9/5\),
\(k=\sqrt2-5/4<1/6\), so \(A^2/k>54/5>23409/2209\).
Also
\[
 d=\frac{77-12\sqrt2}{9(11+6\sqrt2)}>1/3
\]
by \(22>15\sqrt2\). For F,
\(A^2=(4\varphi-5)/9>184/1125\),
\(k=34\varphi-55<167/12500\); their ratio exceeds
\(961/81\), with integer margin \(5752125\).
Its density \(d>2/9\) is equivalent to
\(979>605\varphi\). These seals follow from
\(809/500<\varphi<40451/25000<16181/10000\),
proved by positive squaring.
The correlation bounds then follow from
\((22/125)/(1/3)<9/16\) and
\((22/125)/(2/9)<81/100\).
At fixed native/foreign parameters, \(A^2,d,g^2\)
increase with \(N\), while \(k,B\) stay fixed.

Let \(W_f=u_{r-1}+u_r\). Its square is
\(<2+\sqrt2<7/2\). For every foot,
\[
 \rho=\frac{2/\sqrt{N(N-2)}+
                 \sqrt{2/(N(N-1))}}{\sqrt d}.
\]
For \(N\geq10\), L has the improved \(\rho<8/15\):
the numerator square is at most \(5/36\) at \(10\),
and \(d>7/15+64/2187>49/100\);
thus \(\rho^2<125/441<64/225\).
If the far pair is distinct, its demand is \(>41/40\).
Perpendicular Cauchy, with \(v=\ip{x_A}{U/\norm U}\),
gives
\[
 \ip{x_A}{W_f}\leq
 \rho v+\sqrt{(4/7-v^2)(W_f^2-\rho^2)}.
\]
For L this decreases with \(v\) above its demand floor
and increases with \(\rho\) up to its bound; the sign
condition is \(v^2W_f^2>(4/7)\rho^2\).
At \(N\geq10\) it is
\[
 <(8/15)\sqrt{9/20}+\sqrt{17/40}
 <9/25+33/50=51/50<41/40.
\]
At \(N=9\), use \(\rho<.61,W_f^2<10/3\) instead:
the two terms are \(<.41,.6\), the latter square
being \(1510229/4200000<9/25\).
For H and F, without derivative estimates, use
\(v<\sqrt{4/7}<19/25\) and their pressure floors.
For \(N\geq10\) the bounds are
\(.57+\sqrt{29/200}<.97\) and
\(.7+\sqrt{3/40}<.98\).
At \(N=9\) they improve to \(.95,.97\).
Every far-distinct case is impossible.

If the far pair is common, it supplies a right L/H/F
foot. Normal reversal gives negative raw sum
\(-\sqrt2B'\). Except F/F at \(N=9\), their raw
supports are disjoint and
\[
 \ip{U}{U'}=2BB'/N,\qquad
 \operatorname{corr}(U,U')=\frac{2}{N\sqrt{dd'}}.
\]
The following complete positive-sum Cauchy seals hold
at \(N=9\), and improve thereafter:
\[
\begin{array}{c|c|c|c}
\text{pair}&\rho<&\text{sum demand}>&\text{energy}>\\ \hline
L/L&1/2&2\sqrt{9/20}&3/5\\
L/H&3/5&139/100&19321/32000\\
L/F&5/7&141/100&139167/240000\\
H/H&2/3&36/25&3888/6250\\
H/F&5/6&73/50&15987/27500.
\end{array}
\]
Every last entry exceeds \(4/7\).
For F/F, at any \(N\geq9\) the two actual isolated
end-five walls have exact product
\[
 1-\varphi^2/2+\varphi^2/N\in(-1/2,0).
\]
The lower inequality uses \(\varphi^2<3\);
the upper follows at \(N=9\) from \(\varphi>8/5\)
and strengthens with \(N\).
This violates the native gap without an unsupported
overlap formula. All cases are excluded.
\end{proof}

\subsection{A fourth actual foot and the rank-\(13\) threshold}
\label{app:Afivesevenths}
\begin{proposition}
Every actual \(A_r\) component, \(r\geq13\), has energy
greater than \(5/7\), with arbitrary finite spectators.
\end{proposition}
\begin{proof}
Put \(N=r+1\geq14\) and suppose the energy is at most \(5/7\).
The outer and inner selected correlations increase with \(N\).
At \(N=14\) the exact pair calibrations are
\[
 F(263/500)^2=(263000/319169)^2
  <6791/10000<\sqrt{6/13},
\]
\[
 F(641/1000)^2=(1282000/1410881)^2
  <8257/10000<\sqrt{15/22}.
\]
Thus distinct outer and inner pairs have sum demands
\(>263/250\) and \(>641/500\).
For \(M_i=u_i+u_{N-i}\), the paired Gram identity above
decreases entrywise with \(N\).
At \(14\), \(W=M_1+M_2+2M_3+2M_4\) has diagonal
contribution \(55489/2145<25.869\). Its six off-diagonal
contributions, in order \(12,13,14,23,24,34\), are
\[
 \frac{28}{\sqrt{78}},\
 \frac{112}{\sqrt{429}},\
 \frac{56}{\sqrt{130}},\
 \frac{112}{\sqrt{198}},\
 \frac{28}{\sqrt{15}},\
 112\sqrt{3/110},
\]
respectively less than
\(3.171,5.408,4.912,7.960,7.230,18.497\).
Hence \(W^2<73.047<1461/20\) at every \(N\geq14\).
Four distinct pairs would give
\[
 \ip{x_A}{W}>904/125,\qquad
 (904/125)^2-(5/7)(1461/20)=53923/437500>0.
\]
A common left \(1/2\) or \(3/4\) follows after relabeling.

Recheck the full common-pair inventory at the new
individual cutoff \(C=6\). A central double with second
label at least \(4\) has the first cosine in
\eqref{app:eq:Acentral}, greater than
\(.3(11+14)>6\) at \(N=14\).
The \(5/3\) central double and isolated central \(5\)
have cosines
\[
 (3\varphi+4)\varphi^2\sqrt{(N-4)/N},
 \qquad3\varphi^3\sqrt{(N-4)/N},
\]
both \(>6\) at \(14\), using
\(\sqrt{5/7}>4/5,\varphi>8/5,\varphi^2>13/5,\varphi^3>4\).
They increase with \(N\).
The central \(4/3\) case must now be retained.
Its foreign cost is zero and its actual positive foot is
\[
 \mathrm C:\quad\xi=c+n_3+(2-\sqrt2)n_4,\qquad k=3\sqrt2-4.
\]
Products on \(c,n_3\) vanish; the \(n_4\) product is
\(1-\sqrt2<0\), and all unknown wall products follow
from positivity. The time is positive, with
\[
 A_C^2=4+3\sqrt2-\frac{17+12\sqrt2}{N},\quad
 U_C^2=8-\frac{17+12\sqrt2}{N},\quad B_C=3\sqrt2/2+2.
\]
The complete surviving profiles are L,H,F,C.
All L spectators are cleared in their entirety.

For L, at \(N=14\), \(a_N=11/21\) and
\(8/81<z_0<1/3\). Its bound \(g_L^2>12/25\)
is equivalent to
\[
 4107z_0^2-1200z_0+1859/21>0,
\]
whose discriminant is \(-99884/7\).
Also \(d_L>11/21+64/2187>11/20\), with
last positive comparison \(26880>24057\).
The other exact data at \(14\), with \(s=\sqrt2\), are
\[
\begin{array}{c|c|c|c|c}
 &A_{14}^2&k&B^2&d_{14}\\ \hline
H&(31+22s)/28&s-5/4&(11+6s)/4&
 (66-6s)/(7(11+6s))\\
F&(13\varphi-15)/28&34\varphi-55&(\varphi+1)/4&
 (1525-939\varphi)/(7(\varphi+1))\\
C&(39+30s)/14&3s-4&(17+12s)/2&
 (95-12s)/(7(17+12s)).
\end{array}
\]
For H, \(A^2>11/5,k<33/200\), so \(A^2/k>40/3\).
For F, \(A^2>21/100,k<167/12500\), so
\(A^2/k>2625/167\). For C, \(A^2>5,k<1/4\).
All exceed \(961/81\), proving \(g^2>11/20\).
The density bounds \(d>.3\) reduce to
\(143>62s,\ 15229>9411\varphi,\ 593>372s\).
Use \(s<3/2\) and
\(\varphi<40451/25000\); the cancellation-sensitive
middle comparison has integer margin \(40639\).
All these demands and densities increase with \(N\).

If either far selected pair is common, its reflected
inventory supplies a right L/H/F/C foot.
The raw supports, each of length at most \(5\), are
disjoint. Normal reversal gives positive correlation
\(\rho=2/(N\sqrt{dd'})\).
The complete sum-direction tests are
\[
\begin{array}{c|c|c}
\text{pair type}&\rho<&\text{energy}>\\ \hline
L/L&13/50&16/21\\
L/\text{non-L}&9/25&20449/27200\\
\text{non-L}/\text{non-L}&10/21&231/310.
\end{array}
\]
For the middle row,
\(\rho^2<200/1617<(9/25)^2\) and the demand sum
is \(>\sqrt{12/25}+\sqrt{11/20}>1.43\).
Every energy exceeds \(5/7\).

Otherwise both far selected pairs are distinct.
Their weighted demand is
\[
 \ip{x_A}{R}>452/125,\quad
 R=u_{N-1}+u_{N-2}+2u_{N-3}+2u_{N-4}.
\]
Direct within-end Gram bounds give \(10<J=R^2<29\).
For a left foot,
\[
 \rho=\ip{U/\norm U}{R}=\frac{K_N}{\sqrt d},
\]
\[
 K_N=2\sqrt{\frac8{N(N-4)}}+
 2\sqrt{\frac6{N(N-3)}}+
 \sqrt{\frac4{N(N-2)}}+\sqrt{\frac2{N(N-1)}}.
\]
Its terms decrease with \(N\); at \(14\) they are
\(<.479,.395,.155,.105\), so \(K_N<1.134<8/7\).
Thus \(\rho<31/20\) for L and \(<21/10\) otherwise.
The square seals are
\(1280/539<(31/20)^2,\ 640/147<(21/10)^2\).
Perpendicular Cauchy gives
\[
 \ip{x_A}{R}\leq
 \rho v+\sqrt{(5/7-v^2)(J-\rho^2)},\qquad
 v=\ip{x_A}{U/\norm U}.
\]
On the admissible rectangle it decreases with \(v\)
above its pressure floor and increases with \(\rho\)
below its bound, since \(v^2J>(5/7)\rho^2\).
This holds already with \(J>10\).
Increasing \(J\) to \(29\) is safe.
For L the resulting bound is
\[
 (31/20)\sqrt{12/25}+\sqrt{436199/70000}
 <217/200+5/2=717/200.
\]
For every other foot it is
\[
 (21/10)\sqrt{11/20}+\sqrt{56557/14000}
 <63/40+201/100=717/200.
\]
Both are \(<452/125\), contradicting the far demand.
This exhausts all actual profiles and proves the
nonarithmetic component theorem.
\end{proof}

\subsection{One central-window lemma for \(N=17,18,19\)}
\label{app:Acentralwindow}
The next lemma separates the geometric inventory from the
reference chosen for a particular rank.
\begin{lemma}
Let \(N\in\{17,18,19\}\), and let an actual \(A_{N-1}\)
component have the original farthest-vector projection \(x_A\).
On \(5\leq i\leq N-5\), call \(i\) weak when
\(\ip{x_A}{\omega_i}\leq\sqrt{D_i}-.61\).
Before any contrary energy bound, every weak opposite is:
\[
\begin{array}{ll}
\text{near-four:}&v=\omega_i/\sqrt2+\omega_{i-1}/2,\
 i=5,6,\text{ or a reflected role};\\
\text{pure helper:}&v=\gamma\omega_i+\omega_e/2,\
 e=1,N-1,\quad Q=0,\\
 &\gamma\geq\varphi/2,\quad\gamma<17/16,\
 \text{own }5,\ \text{finite }\geq7,\text{ or dotted}.
\end{array}
\]
Literal weak helpers are nonadjacent, with all spectators
retained. If the energy is at most \(.86\), \(6/7\), or
\(.91\), respectively for \(N=17,18,19\), no near-four
role remains either.
\end{lemma}
\begin{proof}
Both retained old paths have rank at least \(4\).
The complete finite-white list is: none; a simple \(3\)
at an end or an admissible \(D/E\) position; an endpoint
\(4\); or two simple end \(3\)'s, one on each path.
An \(H/F\) completion has insufficient rank, and a joined
completion has rank \(N-1>8\). Finite classification forces
the ordinary \(A\)-path bridge. An additional branch there
cannot be finite. This classification is applied to the
actual endpoint, with the entire spectator forest.

For zero own, internal \(3/4\) and double-white roles
give affine stars or cycles. A far \(3\) has time square
\(<1/(2m)\), \(m\geq5\); a next-to-far \(3\) has
\(Q<2/m\) and time square less than its positive foreign
cost. Either forces an orthogonal closure. End \(4\)
has \(Q<1/m\leq1/5\), hence no spectator and negative
time square. Thus own contact is positive.

For an isolated own, clear all spectators as in the
prefix proof, using the sharper \(A\)-finite pressure
\[
 h_i>\sqrt{D_i-1}-\sqrt{u-1}\quad(u>1),\qquad
 h_i>\sqrt{D_i-u}-\sqrt{u(1-u)}\quad(0<u<1).
\]
At \(u=1\) there is an actual feasible null.
Here \(D_i\geq120/17\), and
\[
 a_i=\frac1{\sqrt{D_i}+\sqrt{D_i-1}}<.1955.
\]
Indeed the denominator exceeds \(2.656+2.461=5.117\)
and \(5.117\cdot391>2000\).
For \(u<1\) the loss is
\(<(.1955+1.019)/2=.60725<.61\).
For foreign \(u>1\), \(u-1\leq(2-\varphi)^2\),
so the loss is \(<.1955+2-\varphi<.5775\).
For pure finite owns at least \(8\), singleton loss is
\(<.1955+\tan(\pi/8)<.1955+.4145=.61\),
since \((2829/2000)^2>2\).
Own \(5\) has its actual \(H_4\) loss
\(<25/76+5/24<.61\); own \(7\) has its actual
\(A_2\)-core loss \(<2000/8181+9/25<.61\).
The latter uses
\(\sqrt{D_i}>53/20,\sqrt{D_i-100/81}>12/5\).
Own \(3,4,6\) and parallel retain their affine/null
principals. Thus every isolated cap is strong.

The strength ratio \(T(q)\) increases here, and
\[
 q\leq\sqrt{19/2}<25/8,\qquad
 T(25/8)=3/2-1441/430050<3/2.
\]
The residual table of the preceding finite-white
classification gives ordinary ratios \(>761/500\):
near simple-\(3\)/own at least \(5\) gives at least
\((101/125+2/5)/\sqrt{3/5}\), \(D\)-position gives
\((101/125+1/5)/\sqrt{2/5}\), and \(E\), end-\(4\)
and two-white rows are stronger. Smaller owns have
the same affine principals. Foreign far helpers have
ratio \(>227/(25\sqrt{35})>761/500\);
foreign near-four has ratio square \(>121/35\).
Consequently every weak cap has \(Q=0\).
The near-four ratio increases with side order \(m\),
and at \(m=7\) has square
\((67+42\sqrt2)/56>9/4\).
Thus only side orders \(5,6\) remain.
For a pure helper the residual identity
\[
 \frac{(\gamma+1/(2m))^2}{(m+1)/(2m)}-(4\gamma-2)
 =\frac{(\gamma-1-1/(2m))^2}{(m+1)/(2m)}
\]
and \(T<3/2\) give \(\gamma<17/16\).

For adjacent literal helpers, the general-\(N\)
AM--GM product bounds in the common-pair calculation
leave mutual \(0\) for the same helper end, and
\(0,-1/2\) for opposite ends. Outward mutual \(0\)
is excluded by the positive projected determinant
\[
 (D_i-4/5)(D_{i+1}-4/5)-C_{i,i+1}^2
 =\tfrac25C_{i,i+1}-24/25+8/(5N)>0.
\]
For reversed mutual \(0\), put \(l=i,s=N-i-1\).
The required determinant is
\[
 P=4lsA^2B^2-(l+1)(s+1)(A^2+B^2)-2AB
 -(l+1)A-(s+1)B+(N+2)/4=0.
\]
Its own derivatives are positive on \([.8,1.08]\),
and it is concave in \(l\), \(5\leq l\leq N-6\).
At \(N=17\), both owns at least \(7\) give
\(P(.9,.9)\geq14.632>0\). One \(5\) and one own at
least \(8\) give, at \((.809,.923)\) or its reverse,
minimum \(1.53076164678>0\).
These endpoint values increase for \(N=18,19\):
at \(l=5\) the coefficient of \(N\) is
\(20A^2B^2-6(A^2+B^2)-B+1/4>0\)
at both comparison pairs and their reversals.
If owns are \(5,7\), the determinant is a polynomial
of degree at most \(2\) in \(\cos(\pi/7)\) over
\(\mathbb Q(\sqrt5)\), with nonzero linear coefficient.
But \(2\cos(\pi/7)\) has irreducible polynomial
\(t^3-t^2-2t+1\), and degree \(3\) over
\(\mathbb Q(\sqrt5)\). Thus it cannot vanish.

For two native-five helpers, use raw norm-two duals
\(C_{ij}=\min(i,j)-ij/N\) and
\(v_i=\varphi\omega_i+\omega_e\).
Same-end adjacent difference square is
\(\varphi^2(N-1)/N<4\), forcing positive wall product.
For opposite ends, with orientation \(\epsilon=\pm1\),
\[
 \norm{v_i-v_{i+1}}^2=
 \varphi^2(N-1)/N+2(N-2)/N-4\epsilon\varphi/N<6.
\]
The raw native gap forces mutual \(0\).
For \(H_i=v_i^2-2\), \(S=v_i\cdot v_{i+1}\),
the rational coefficient of
 \(N^2(H_iH_{i+1}-S^2)\) in basis \(1,\varphi\)
is \(N(N+6)\) outward and \(-N^2\) reversed.
It cannot vanish. Hence opposite helpers have
mutual exactly \(-1/2\) whenever they survive.

The retained two-cap matrices are the common-helper
\(I-t\mathbf1\mathbf1^T\), or the opposite-helper
matrix with diagonal in \((1/2,3/5]\), cross \(-1/2\).
 They have positive inverses. Positive actual projection
 retains the entire old spectator block.
 Put \(\lambda=H^{-1}(A_i,A_j)^T>0\),
 \(\mathcal D=A_i\lambda_i+A_j\lambda_j\) and
 \(Z=W_i\omega_i+W_j\omega_j\).
 For a common left helper,
 \(W_i=\lambda_i\gamma_i+(\lambda_i+\lambda_j)/(2i)\),
 \(W_j=\lambda_j\gamma_j\); reflect for a common right helper.
 Outward helpers give
 \(W_i=\lambda_i(\gamma_i+1/(2i))\),
 \(W_j=\lambda_j(\gamma_j+1/(2(N-j)))\).
 Reversed helpers give
 \(W_i=\lambda_i\gamma_i+\lambda_j/(2i)\),
 \(W_j=\lambda_j\gamma_j+\lambda_i/(2(N-j))\).
 Thus \(W=W_i+W_j>.8L\), \(L=\lambda_i+\lambda_j\).
 All cap times exceed \(2\): the smallest time square is
 \(>(16/25)(120/17)+(4/5)(10/17)+8/17-1
 =379/85>4\).
 For a common helper the same explicit inverse as in
 the prefix proof gives \(\sqrt{1+\mathcal D}/W<7/16\).
 For opposite helpers,
 \(5/6<\lambda_i/\lambda_j<6/5\), and
 \[
 \mathcal D/L^2\leq3/55,\quad
 A_i+A_j\leq L/10,\quad L>40.
 \]
 The first inequality follows by maximizing
 \((.6(t^2+1)-t)/(t+1)^2\) on \([5/6,6/5]\).
 Consequently
 \(\sqrt{1+\mathcal D}/W
 <\sqrt{15/176+1/1024}<3/10\).
The spatial rationalization loss is \(<W/10\):
use \(\sqrt{120/17}>2.65\),
\(\sqrt{110/17}>2.54\).
The weighted mean losses are thus \(<43/80\)
or \(<2/5\), both \(<.61\). Adjacent weak helpers
are impossible independently of an energy target.

A near-four has an actual positive finite mate
with cosine \(>7\), hence normalized height \(>.86\).
At the opposite end there is another mate, or a
strong anchor of normalized height \(>.77\).
The correlations are bounded respectively by
\(<.5,<.44,<.41\) at \(N=17,18,19\).
They follow by substituting mate index at most
\(5\) and anchor index at least \(N-6\) in
\(\sqrt{j(N-k)/(k(N-j))}\).
At the three contrary energies, Cauchy would give
\[
 1.63^2<.86(3),\quad
 1.63^2<(6/7)(2.88),\quad
 1.63^2<.91(2.82),
\]
each false. The reflected roles satisfy the same
bounds. This removes all near-fours.
\end{proof}

\subsection{Shared actual endpoint cores and three reference tables}
\label{app:Aendpointcores}
For a weak end \(i=5\) with helper \(e=1,N-1\), let \(j\)
be the old neighbor of the helper end. All three old
selected roles are present in each \(N=17,18,19\).
Set \(v=\gamma\omega_5+\omega_e/2\). The actual
time squares are
\begin{equation}\label{app:eq:Aendtimes}
 A_{\mathrm{far}}^2=
 \frac{10(N-5)\gamma^2+10\gamma-(N+1)/2}{N},\qquad
 A_{\mathrm{near}}^2=A_{\mathrm{far}}^2+
                        \frac{2(N-10)}N\gamma.
\end{equation}
These follow directly from the actual dual Gram.
Reflected weak ends have the reflected formula.

There are three positive actual constructions:
\[
\begin{array}{c|c|c}
\text{own}&\text{old correction }b_c&\text{residual}\\ \hline
5&\gamma n_5+(2/3)n_e+(1/3)n_j&
 \delta=2/3-\gamma^2>0\\
\text{finite }\geq7&\gamma n_5+2(1-\gamma^2)n_e&
 -k,\ k=(1-\gamma^2)(4\gamma^2-3)>0\\
\text{dotted}&n_5/\gamma&-k,\ k=1-\gamma^{-2}>0.
\end{array}
\]
For own \(5\), \(\eta=c+b_c\) is spacelike and the
actual future finite point is
\(p+(A/\delta)\eta\), of square \(-1-A^2/\delta\).
It is orthogonal to all four selected roots, while
all other wall products are nonpositive.
For finite owns at least \(7\), put \(d=1-\gamma^2\).
The core \(\xi=c+b_c\) has products zero with \(c,n_5\)
and \(-1/2+2d<0\) with \(n_e\).
The dotted core has products zero with \(c\) and
\(-\gamma+\gamma^{-1}<0\) with \(n_5\).
All time coefficients are \(A>0\). Positivity checks
every other actual wall, including the other helper,
every spectator and all unknown walls.
Thus, for \(U=v-b_c\), the actual point pressures give
\begin{equation}\label{app:eq:Acorepressure}
 \ip{x_A}{U}>A-\sqrt{\delta}\quad(\text{own }5),
 \qquad
 \ip{x_A}{U}\geq A-\sqrt{k}\quad(\text{other rows}).
\end{equation}
Dropping the positive old-normal products yields
the same lower bounds for \(\ip{x_A}{v}\).

The rational concave references needed below have
linear end segments, reflection \(H_{N-i}=H_i\),
and the following central heights:
\[
\begin{array}{c|rrrrrr}
N&H_5&H_6&H_7&H_8&H_9&H_{10}\\ \hline
17&2.046&2.176&2.259&2.300&2.300&\\
18&2.077&2.218&2.314&2.371&2.389&\\
19&2.104&2.255&2.363&2.433&2.467&2.467.
\end{array}
\]
Only the independent entries before reflection are used.
Positive square tests \((H_i+.61)^2<D_i\) pay every
retained strong floor. If \(5\) is weak use its chord
to \(6\); reflect at the other end; omit every
internal weak knot by its neighbor chord.
Nonadjacency proves that all losses add.
The complete zero/one/two-weak-end tables are
\[
\begin{array}{c|c|c|c}
N&\text{number of weak ends}&\text{initial energy}&
 \text{maximum internal loss}\\ \hline
17&0&.8626932&3(.047^2/4)=.00165675\\
17&1&.8302129333\ldots&.00165675\\
17&2&.7977326666\ldots&2(.042^2/4)=.000882\\
18&0&.8954558&.002025\\
18&1&.8640827333\ldots&.00151875\\
18&2&H_6^2/6+.012789&3(.039^2/4)\\
19&0&.9258842&4(.043^2/4)=.001849\\
19&1&.8955541833\ldots&.001849\\
19&2&.8652241666\ldots&3(.038^2/4)=.001083.
\end{array}
\]
The repeating decimals abbreviate exact rational
sums of squared rational slopes.
These tables are complete: after zero/one/two end
chords, the eligible internal weak intervals and
nonadjacent cardinalities give the indicated
maximal counts. Each loss is one quarter of the
squared second difference. End chords preserve
the neighboring strong anchor; an adjacent omission
can only increase its available individual price.
Every omitted knot still has price zero.

\begin{proposition}[The \(N=17\) paid reference]
\label{app:prop:A16}
An actual \(A_{16}\) component has energy \(>43/50\)
with arbitrary finite spectators.
\end{proposition}
\begin{proof}
Assume contrary \(e\leq.86\). The central-window
lemma removes near-fours. Use the \(N=17\) table.
Zero weak ends already give \(e\geq.86103645\).
For two weak ends the reference is flat on both
six-coordinate end blocks, has energy
\(E\geq1195276/1500000\), and separate anchor capacities
\[
 \beta_6,\beta_{11}\geq
 (7H_6/6-H_7)/2=839/6000.
\]
For one weak end the block lengths are \(6,5\),
\(E>.8285\), and
\(\beta_6\geq839/6000,\beta_{12}\geq349/2500\).

For a left weak helper, project its core displacement
\(U\) onto vectors constant on each end block.
In the six/six case,
\[
 \Pi U=(5\gamma/6)\omega_6+\omega_{11}/12
       \quad(\text{far}),\qquad
 \Pi U=(5\gamma/6+1/12)\omega_6
       \quad(\text{near}).
\]
All old-root corrections have mean zero in their
own block. Thus \(Y=U-\Pi U\) sums to zero on each
block, \(\ip{w}{Y}=0\), and its square is
\[
\begin{array}{c|c}
\text{own}&\norm{Y}^2\text{ in the far orientation}\\ \hline
5&(2/3)\gamma^2+1/12\\
\text{finite }\geq7&4\gamma^4-(16/3)\gamma^2+29/12\\
\text{dotted}&(5/3)\gamma^2-19/12+\gamma^{-2}.
\end{array}
\]
The near orientation adds \(\gamma/3\).
These identities follow by expanding the actual
\(N=17\) Gram and
\[
 b_c^2-2\ip{v}{b_c}
 =k_c^2+l_c^2+m_c^2-l_cm_c-2\gamma k_c-l_c,
\]
where \(b_c=k_cn_5+l_cn_e+m_cn_j\).
The finite polynomial increases for \(\gamma^2>.81\)
and is \(<13/12\). The dotted one increases for
\(1<\gamma^2<289/256\) and is
\(<229/768+256/289<148/125\), with cleared margin \(271\).
Consequently every orientation has
\(\norm{Y}^2<148/125+17/48<8/5\).

The projection demand is \(H_6(5\gamma/6+1/12)\).
Let \(\Delta\) be the actual core demand minus this
quantity. For native \(5\), the fresh time bound
\(A_{\mathrm{far}}>2.136\) and radical \(<.111\)
give
\[
 \Delta>2.136-.111-H_6(91/120)>.327.
\]
For finite owns at least \(7\), use
\(A_{\mathrm{far}}(.9)>2.391,\sqrt k\leq1/4\);
for dotted, use \(A_{\mathrm{far}}(1)>2.667\)
and radical \(<.338\). They give, respectively,
\[
 \Delta>2.391-.25-5H_6/6=983/3000>.327,
 \quad
 \Delta>2.667-.338-11H_6/12>.327.
\]
The deficits increase with \(\gamma\):
\[
 (A_{\mathrm{far}}')^2-120/17
       =1105/(289A_{\mathrm{far}}^2)>0,\qquad
 \sqrt{120/17}>5H_6/6.
\]
The near time is larger with the same radical and
reference pairing. These inequalities therefore
cover all finite/dotted intervals and both orientations.

For one weak end, the far projection instead uses
\((5\gamma/6)\omega_6+\omega_{12}/10\).
Its residual square decreases by \(1/60\), and its
deficit decreases by exactly
\(H_5/10-H_6/12=349/15000\).
Hence \(\norm Y^2<8/5,\Delta>.303\).

Use weight \(\mu=.1\), paying each original anchor
before adding its actual core pressure. One near
foot spends at most \(31/320\) at \(6\); a far foot
spends less there and \(1/100\) at \(12\).
For two ends the combined debit at each \(6/11\)
is at most
\[
 .1(5(17/16)/6+1/6)=101/960<839/6000.
\]
All surviving original coefficients are nonnegative.
With \(W=w+\mu\sum Y\), one original-vector square gives
\[
 e\geq E+2\mu\sum\Delta-\mu^2\norm{\sum Y}^2.
\]
For one weak end the gain exceeds
\(.2(.303)-.01(8/5)=.0446\), giving \(e>.8731\).
For two ends retain the full cross by
\(\norm{Y_L+Y_R}^2\leq2\norm{Y_L}^2+2\norm{Y_R}^2<32/5\).
The gain exceeds
\(.2(2\cdot.327)-.01(32/5)=167/2500\),
and
\[
 e>1195276/1500000+167/2500
   =1295476/1500000>.86.
\]
All three cases contradict the contrary bound.
\end{proof}

\subsection{Raw helper-end pressures at \(N=18,19\)}
\label{app:Arawends}
For weak \(5\), put \(z=x_1\) for helper \(1\) and
\(z=-x_N\) for helper \(N-1\). Then
\(\ip{x_A}{v}=\gamma h_5+z/\sqrt2\).
The shared actual cores give deficits
\(>.313\) at \(N=18\) and \(>.306\) at \(N=19\).
Here are the complete seals.
The weak floors are \(h_5<2.078\) and \(h_5<2.105\).
For native \(5\), radical is \(<.111\), own \(<.81\),
and times are \(>2.155,>2.17\); the deficits
exceed \(.36082,.35395\).
For finite owns at least \(7\), the core radical
has square \(k=(1-\gamma^2)(4\gamma^2-3)\leq1/16\):
\[
\begin{array}{c|c|c|c|c|c}
\text{own}&\gamma\text{ bracket}&\sqrt k<&
 A_{18,\rm far}>&A_{19,\rm far}>&
 \text{deficit at }18;\ 19\\ \hline
7&(.9009,.901)&.2158&2.4154&2.434&
 .327322;\ .321595\\
8&(.9238,.9239)&.2464&2.4795&2.499&
 .3132358;\ .3077905\\
\geq9&\gamma>.93&1/4&2.497&2.517&
 .31446;\ .30935.
\end{array}
\]
The native-seven polynomial
\(8\gamma^3-4\gamma^2-4\gamma+1\) increases
above \(\sqrt3/2\); its signs at \(.9009,.901\)
have numerators \(-569738168,257608\) and
denominators \(10^{12},10^9\).
Native eight is \(\sqrt{2+\sqrt2}/2\).
The residual polynomial increases on both brackets;
substitution of upper endpoints gives the radical
bounds. The tight eight comparison is
\(14640879\cdot41436484<607129600000000\).
Times come from \eqref{app:eq:Aendtimes} at the
lower endpoints. Every finite label at least \(9\)
has \(\gamma>.93\), by monotonicity and
\(8(.93)^3-6(.93)-1<0\).
The functions \(A_{\rm far}-2.078\gamma\) or
\(A_{\rm far}-2.105\gamma\) increase, since
\(A_{\rm far}'>\sqrt{D_5}\).
Thus the last row covers every finite label.

For dotted \(1<\gamma<17/16\), retain the exact
floor \(h_5\leq\sqrt{D_5}-.61\).
At \(18\),
\[
 A_{\rm far}-\sqrt{D_5}\gamma=
 \frac{(5/9)\gamma-19/36}
      {A_{\rm far}+\sqrt{D_5}\gamma}>1/216.
\]
Its numerator is \(>1/36\), denominator \(<6\).
At \(19\) the same difference is positive.
The function \(.61\gamma-\sqrt{1-\gamma^{-2}}\)
decreases throughout the interval: its radical
derivative exceeds \(1\).
Using \(\sqrt{33}/17<.338\), its endpoint lower
bound is \(>.310125\), and adding \(1/216\)
at \(18\) gives \(>.313\).
Consequently
\[
 z>.442\quad(N=18),\qquad z>.432\quad(N=19).
\]
A near helper at \(19\) has the same radical and
time larger by
\[
 A_{\rm near}-A_{\rm far}>
 \frac{(18/19)(.8)}7>.108,
\]
using \(A_{\rm far}<3,A_{\rm near}<4\).
Its deficit is \(>.414\), so \(z>.58\).


\begin{proposition}[The last two ordinary \(A\)-rates]
\label{app:prop:A17A18}
With arbitrary finite spectators,
\(e(A_{17})>6/7\) and \(e(A_{18})>91/100\).
\end{proposition}
\begin{proof}
At \(N=18\), assume \(e\leq6/7\).
The central-window lemma removes near-fours.
The reference table excludes zero and one weak end:
the retained energies are \(.8934308\) and
\(>.8625639833\ldots\).
With two weak ends, the reference is flat on both
six-coordinate end blocks, of signed mean \(b\)
with \(1/4<b<.262\), and
\[
 E\geq H_6^2/6+.012789-3(.039^2/4)
      =9978827/12000000.
\]
One actual weak helper supplies endpoint \(z>.442\)
at either end. Change that signed endpoint of the
reference to \(.442\) and its other five block
coordinates to \((6b-.442)/5\).
The block sum is preserved; its last coordinate
is \(>.2116>.096\), above the next coordinate.
All new positive duals are paid by the old
strong floors and this actual endpoint pressure.
The single reference's energy gain is
\[
 (6/5)(.442-b)^2>(6/5)(.442-.262)^2=.03888.
\]
Hence
\[
 e>9978827/12000000+.03888
   =10445387/12000000>6/7.
\]
Only one helper is used; either orientation works.

At \(N=19\), assume \(e\leq.91\); no near-four
remains. The reference table gives
\[
 E_0\geq.9240352,\quad E_1>.8937,\quad E_2>.86414.
\]
A strong end is a length-five block with
\(.29<b_0<.298\); a weak end is a length-six block
with \(1/4<b_1<.266\).
Changing a signed endpoint to \(z_0>b\)
raises energy exactly \(m(z_0-b)^2/(m-1)\),
preserving the block sum.
For \(m=5,z_0=.432\), the other coordinates
exceed \(.2545\), above \(.151/\sqrt2\);
for \(m=6,z_0\leq.58\), they exceed \(.184\),
above \(.108/\sqrt2\).
Neighbor omissions only decrease the next
coordinate. Opposite end blocks are disjoint.

Zero weak ends already exceed \(.91\).
With one weak end and a near helper,
\[
 e>.8937+(6/5)(.58-.266)^2=1.0120152.
\]
Otherwise its helper is the opposite strong end:
\[
 e>.8937+(5/4)(.432-.298)^2=.916145.
\]
With two weak ends, any near helper gives
\(.86414+(6/5)(.58-.266)^2=.9824552\).
If both helpers are far, they use distinct
opposite endpoints. Paying both in the disjoint
length-six blocks gives
\[
 e>.86414+2(6/5)(.432-.266)^2=.9302744.
\]
These exhaust the end patterns and orientations,
always on one modified reference for the same \(x_A\).
\end{proof}


\section{The local \(A_{18}\) inventory at comparison \(11/8\)}
\label{app:localAinputs}
This appendix supplies the earlier geometric body used in
Lemma~\ref{lem:localA18}; it does not use that lemma's
endpoint pressure or its original \(B_{18}\) application.
We work at the actual whole-\(A_{18}\) vertex \(q\)
with the same strict interior point and comparison
\(-\ip{z}{r}\leq\kappa E_q\), \(\kappa<11/8\).

\subsection{All zero-own and isolated boundaries}
On \(5\leq i\leq14\), deletion leaves path orders
\(\ell=i,r=19-i\), both at least \(5\).
The complete one-node finite-white list is none;
one simple \(3\) at an end, a \(D\) position or an
admissible \(E\) position; one terminal \(4\);
or two simple end \(3\)'s joining the two paths.
Two contacts in one path make a nonfinite cycle.
A joined completion has rank \(18\), hence is
classical with only new simple end bonds.
There are no foreign components at this vertex.

At zero own, a global endpoint \(3\), next-to-end
\(3\), or endpoint \(4\) has time square
\(9/19-1,17/19-1,18/19-1<0\), respectively.
Other simple-\(3\) placements retain affine
\(E_8\) arms \(1,2,5\) or \(E_7\) arms \(1,3,3\).
Internal \(4\) gives the \(4/3/3\) star; two
simple white contacts close an all-three cycle.
No contact gives zero derivative and time square \(-1\).
The affine kernels are positive: for \(E_7\)
take center \(4\) and arms \(2;(3,2,1);(3,2,1)\);
for \(E_8\), center \(6\), arms
\(3;(4,2);(5,4,3,2,1)\).
For a \(4/3/3\) star, take center \(\sqrt2\),
doubled leaf \(1\), simple leaves \(1/\sqrt2\);
all-three cycles have coefficients \(1\).
Positive cap time makes every such null nonzero
future, and all unknown original wall products
are nonpositive. Compactness excludes them.

Isolated own \(3\) retains the affine \(E_7\)
star with the cap as short arm; own \(4\) has
the \(4/3/3\) star; own \(6\) the \(6/3\) chain,
with kernel \((\gamma,1,1/2)\); parallel own
has null \(c+n_i\).
White \(3\)/own \(3\) closes an all-three cycle.
Far white \(4\)/own \(3\) retains the \(E_7\)
star and far white \(4\)/own \(4\) the doubled
star. For pure helpers own \(3,4,6\) or parallel
retains respectively the path cycle, doubled
star, \(6/3\) chain or \(c+n_i\).
Additional unselected contacts preserve all
wall signs in these actual null arguments.

For isolated positive own put \(\tau=\gamma^{-2}\).
For \(0<\tau<1\), the actual positive core
\(c+\gamma\tau n_i\) gives
\[
 h_i>\sqrt{D_i-\tau}
           -\kappa\sqrt{\tau(1-\tau)}.
\]
At \(\tau=1\) it is a feasible null.
For \(\tau>1\), positive finite projection on
\(c,n_i\) gives the exact pressure
\[
 h_i\geq
 \sqrt{\frac{D_i-1}{D_i-\tau}}\,
       \bigl(\sqrt{D_i-1}-\kappa\sqrt{\tau-1}\bigr).
\]
The future point, not its spacelike correction,
is normalized. Every unknown wall is checked
by the positive selected-root coefficients.
Here \(D_i\geq140/19\),
\(a_i=(\sqrt{D_i}+\sqrt{D_i-1})^{-1}<5/26<.193\).
For \(0<\tau<1\), the loss is
\[
 <a_i\tau+\kappa\sqrt{\tau(1-\tau)}
 <(.193+1.389)/2=.791<.87.
\]
The radical seal is \(.193^2+(11/8)^2<1.389^2\).
For finite owns at least \(8\), the positive-bracket
finite formula has loss
\(<5/26+(11/8)(.415)<.764\).

Native \(5\) uses its positive \(H_4\) projection,
old coefficients \(\gamma(3/2,1,1/2)\),
residual \(\delta_5=(7-3\sqrt5)/16<1/36\),
and cap time square \(>3.7\).
The finite-point radical factor
\(\sqrt{1+\delta_5/A^2}<1.01\) gives loss
\[
 <25/76+(11/8)\frac{1.01}{6(.8)}<.619.
\]
Native \(7\) uses the actual negative core
\(c+(4\gamma/3)n_i+(2\gamma/3)n_{\rm neighbor}\),
whose loss is
\[
 <10000/41877+(11/8)(9/25)<.735.
\]
Use \(.9<\cos(\pi/7)<10/11\), verified
by its increasing cubic, and
\(\sqrt{D_i-100/81}>2.47\).
Thus every isolated finite or dotted cap is
strong at loss \(.87\) and comparison \(11/8\).

\subsection{Every nonisolated ratio and the surviving helpers}
The ordinary ratio threshold \(3/2\) at this
comparison is proved in Lemma~\ref{lem:localA18}
by its explicit concave polynomial.
For one white \(3\) at side position \(t\)
in a path of order \(\ell\),
\[
 R=\frac{\gamma+t/(2\ell)}
          {\sqrt{1-t(\ell-t)/(2\ell)}}.
\]
Near \(3\), after the own-\(3\) cycle, has ratio
\(>(.808+2/5)/\sqrt{3/5}>3/2\),
except the own-\(4\) near-four role.
At \(D\) positions \(t=2,\ell-2\), after the
own-\(4\) star, ratio is
\[
 \geq\frac{\gamma+1/\ell}{\sqrt{2/\ell}}
 >\frac{.808+1/5}{\sqrt{2/5}}>3/2.
\]
The lower expression increases for \(\ell\geq5\).
At the allowed \(E\) positions \(\ell=6,7,8\),
denominator is at most \(1/2\) and own \(>.8\).
Near endpoint \(4\) exceeds \(3/2\) even at
own \(3\). Far endpoint \(4\) has own \(3,4\)
nulls, and otherwise ratio
\(\geq\gamma\sqrt\ell>3/2\).
Two-white-\(3\) has ratio at least
\(\gamma\sqrt{D_i}+1/\sqrt{D_i}>3/2\)
already at \(\gamma=1/2\).
All larger finite and dotted owns strengthen
these tests. This exhausts the finite-white table.

Only near-four with side order \(5,6\), or
pure far helpers, remain weak. The near-four
and pure-helper formulas, the order-\(7\) seal,
and the entire helper cutoff \(\gamma<17/16\)
are derived in the main proof.
This completes its own/white inventory without
a farthest hypothesis at \(q\).

\subsection{Native-five determinant and positive mixed weights}
For two adjacent native-five helpers use raw
norm-two duals \(C_{ij}=\min(i,j)-ij/N\).
Put \(H_i=v_i^2-2,S=v_i\cdot v_j,U=H_i-S,V=H_j-S\)
and \(P=i(N-i-1)+1\), with \(j=i+1\).
For outward helpers, the exact expansions are
\[
 NS=P+(P+N-2)\varphi,\quad
 NU=(i-N-2)+(N-i+1)\varphi,\quad
 NV=-i-3+(i+2)\varphi.
\]
For reversed helpers they are
\[
 NS=P+(P+N)\varphi,\quad
 NU=(i-N-2)+(3i-N-1)\varphi,\quad
 NV=-i-3+(2N-3i-4)\varphi.
\]
Since \(H_iH_j-S^2=S(U+V)+UV\), substituting
\(\varphi^2=\varphi+1\) gives rational coefficients
\(N(N+6)\) and \(-N^2\), respectively.
At \(N=19\) these are \(475,-361\), both nonzero.
The raw difference-square argument in the
central-window lemma forces mutual zero first,
so neither determinant can satisfy that zero
mutual. No arithmetic criterion is used.

For the surviving mixed pairs, use the explicit
positive Schur matrices, coefficients
\(\lambda=H^{-1}(A_i,A_j)^T\), and all four
orientation formulas for \(W_i,W_j\) in the
central-window lemma. They are purely geometric,
so remain valid at comparison \(\kappa\).
Cap times still exceed \(2\); hence
\(\sqrt{1+\mathcal D}/W<7/16\) for a common
helper and \(<3/10\) for opposite helpers.
With \(N=19\), \(\sqrt{D_i}>2.7\) and
\(\sqrt{C_{i,i+1}}>2.6\).
Rationalizing \(E_c-\norm Z\), as in the
prefix proof, bounds the spatial loss by \(W/10\).
The actual mixed-point pressure now is
\[
 \ip{y}{Z}\geq
  1+\mathcal D-\kappa\sqrt{1+\mathcal D}.
\]
Thus its weighted mean loss is \(<.7015625\)
or \(<.5125\), respectively.
Both are less than \(.87\), proving the
required nonadjacency against every original wall.
Together with the new endpoint/mate/reference
argument already written, this completes all
earlier local inputs of Lemma~\ref{lem:localA18}.


\subsection{The fresh \(N=16\) inventory and \(A_{15}>.826\)}
\label{app:A15}
\begin{proposition}
With arbitrary finite spectators, an actual \(A_{15}\)
component has energy greater than \(413/500\).
\end{proposition}
\begin{proof}
This proof consolidates \src{A15-new} with only the fresh
\(N=16\) inventory/reference portions of \src{A15-two}; its older
numerical consumers are not needed.
Use the exact interval \(5\leq i\leq11\), loss \(.61\),
and \(55/8\leq D_i\leq8\).
The finite-white list is the two-path list already
proved: each retained path has rank at least \(4\),
and a joined completion has rank \(15>8\).
The zero-own affine and foreign-shadow repairs need
only side order at least \(5\), so give precisely
the same exhaustive cases at this rank.

For isolated caps, put
\(a=(\sqrt D+\sqrt{D-1})^{-1}<1/5\).
The cleared \(0<u<1\) loss is
\(<(.2+\sqrt{1.04})/2<.61\).
For foreign \(u>1\) it is
\(<.2+2-\varphi<.582\).
For pure finite labels at least \(8\), use the
exact singleton pressure
\[
 T=r\sqrt{\frac{r-s}{r+s}},\quad
 r=\sqrt{q^2-1},\quad
 s=\tan(\pi/m)\leq\sqrt2-1<.415.
\]
Its loss \(q-T\) decreases with \(q\) and increases
with \(s\). For the first monotonicity, the
positive derivative comparison reduces to
\[
 (1-s^2)r^4+2sr^3+(2s^4-s^2)r^2-2s^3r+s^4
 \geq r^4/4>0
\]
on \(r\geq1,0\leq s\leq1/2\).
At \(q_0=\sqrt{55/8}<2.623\),
\(2.423<r_0=\sqrt{47/8}<2.425\), and
\[
 T^2>(47/8)(2.008/2.84)=11797/2840>(2.03)^2.
\]
The loss is \(<.593<.61\).
Own \(5\) has its positive \(H_4\) loss
\(<25/76+5/24\); own \(7\) has
\(<10000/40257+9/25<.61\).
Own \(3,4,6\), parallel and the \(u=1\)
case have their actual affine/null exclusions.
Thus all isolated caps are strong.
The ordinary white ratios exceed \(3/2>T(q)\)
using only side order at least \(5\).
The complete weak list is near-four at \(5,6\)
or its reflection, and pure helpers
\(\gamma\omega_i+\omega_e/2\),
\(e=1,15,\ Q=0,\ \varphi/2\leq\gamma<17/16\).
All positive-foreign roles are strong.

For adjacent helpers, the same actual-product
bounds leave mutual \(0\) for a common helper
and \(0,-1/2\) otherwise.
Outward zero contradicts
\[
 (D_i-4/5)(D_{i+1}-4/5)-C_{i,i+1}^2
 =(2/5)C_{i,i+1}-24/25+1/10\geq41/25>0.
\]
For reversed zero, the polynomial of the
central-window proof has \(l+s=15\) and
constant \(9/2\). It increases in each own
on \([.8,1.08]\) and is concave in \(l\).
At its endpoint \(l=5\) or \(10\), both
owns at least \(7\) give \(P(.9,.9)=11.88>0\).
One \(5\) and an own at least \(8\) give
the exact minimum
\(P(.809,.923)=.0907947698>0\);
the reflected endpoint is larger by
\(5(.923-.809)\).
Two native-five helpers are impossible by
the raw determinant proof, valid for \(N\geq15\).
Native \(5/7\) is impossible by the same
nonzero quadratic over \(\mathbb Q(\sqrt5)\)
versus cubic degree argument. No field
restriction is imposed on any other wall.
The surviving opposite-helper mutual is
therefore \(-1/2\).
The positive mixed Schur point appends all
spectators; its times satisfy
\(A_i^2>699/160>4\).
The explicit inverse weights give the same
mean losses \(<43/80\) or \(<2/5\).
The spatial term is \(<W/10\), since
\(\sqrt{55/8}>2.6,\sqrt{25/4}=2.5\).
Helpers are consequently nonadjacent before
any energy target is used.

Assume now \(e\leq413/500\).
A near-four's actual mate has demand \(>.86\).
At the other end, a mate or a strong \(10/11\)
has demand \(>.76\); their correlation is at most
\(\sqrt{3/11}<11/20\).
Cauchy would give
\(1.62^2<(413/500)(31/10)=2.5606\),
whereas \(1.62^2=2.6244\).
Thus all weak sites are pure helpers.

Use reflected rational heights
\[
 H_5=H_{11}=a=2.012,\quad
 H_6=H_{10}=C=2.128,\quad
 H_7=H_9=2.196,\quad H_8=2.218
\]
and linear ends.
Flatten an outer weak end to its strong
neighbor; omit all internal weak knots
by chords. The exact complete bounds are
\[
\begin{array}{c|c|c}
\text{weak ends}&\text{maximum internal loss}&E>\\ \hline
0&3(.048^2/4)=.001728&.8264648\\
1&2(.048^2/4)=.001152&.7928637333\ldots\\
2&2(.046^2/4)=.001058&1138171/1500000.
\end{array}
\]
Every positive dual uses an actual strong
floor. For one weak end the block lengths
are \(6,5\), with separate prices
\(\beta_6\geq43/300,\beta_{11}\geq179/1250\).
For two weak ends they are \(6,6\), and
\(\beta_6,\beta_{10}\geq43/300\).
Adjacent omissions increase these capacities.

Use the same three actual endpoint-core
constructions of
Section~\ref{app:Aendpointcores}, now with
the freshly computed times
\[
 A_{\rm far}^2=(220\gamma^2+20\gamma-17)/32,\quad
 A_{\rm near}^2=A_{\rm far}^2+3\gamma/4.
\]
In the six/six reference, the core's block
projection is
\((5\gamma/6)\omega_6+\omega_{10}/12\)
far, and
\((5\gamma/6+1/12)\omega_6\) near.
Fresh substitution of the \(N=16\) dual Gram
gives exactly the three residual polynomials
in the \(N=17\) payment proof, with the
same near addition \(\gamma/3\).
This is the fixed-block variance identity,
not an imported different-rank energy bound.
Thus \(Y=U-\Pi U\) is orthogonal to the
reference and \(\norm Y^2<8/5\).

Let \(\Delta\) be the core demand minus
\(C(5\gamma/6+1/12)\).
All orientations and own intervals have
\(\Delta>1/3\). Indeed
\((A_{\rm far}')^2-55/8=15/(4A_{\rm far}^2)>0\).
At native \(5\), use time \(>2.11\),
radical \(<.111\), projection \(<C(91/120)\);
the deficit is \(>5779/15000\).
For finite owns at least \(7\), use
\(A_{\rm far}(.9)=\sqrt{28/5}>2.366\)
and radical \(\leq1/4\), giving \(>257/750\).
For dotted owns, use
\(A_{\rm far}(1)=\sqrt{223/32}>2.639\)
and radical \(<.338\), giving \(>1051/3000\).
Monotonicity of \(A_{\rm far}\) minus its
linear projection makes these full-interval
bounds; near times are larger.
In the six/five case, a far projection
uses \(\omega_{11}/10\) instead.
Its variance decreases \(1/60\), and
its deficit decreases
\(a/10-C/12=179/7500\).
Hence \(\Delta>2321/7500>.309\).

Pay \(\mu=.1\) before using each core pressure.
One foot costs at most \(31/320<43/300\)
at the weak-side strong anchor and, if far,
\(1/100<179/1250\) at the opposite anchor.
For two feet, the combined cost at each
\(6/10\) is at most \(101/960<43/300\).
Thus \(W=(w-\mu\sum\Pi U)+\mu\sum U\)
has a nonnegative old strong-floor part and
the actual core pressures, giving
\[
 e\geq E+2\mu\sum\Delta-\mu^2\norm{\sum Y}^2.
\]
One weak end gains \(>.0458\), hence
\(e>.7928+.0458=.8386\).
For two ends retain the full cross through
\(\norm{Y_L+Y_R}^2<32/5\).
The gain exceeds \(26/375\), so
\[
 e>1138171/1500000+26/375
   =1242171/1500000=.828114>413/500.
\]
Zero weak ends already exceed the target.
This proves the component theorem.

Every ordinary rank-three forest costs
\(>17/70\), so this new component rate
also excludes its ordinary spectators directly.
The other retained spectators \(I+T_2\)
and \(O_7+A_1\) cost \(>.21\) and
\(>.09+2/15\); each sum likewise exceeds \(1\).
A separate historical ordinary-rank-three
consumer is therefore unnecessary.
\end{proof}


\subsection{The fresh \(N=13\) component and its native-cosine split}
\label{app:A12}
\begin{proposition}
With arbitrary finite spectators, an actual \(A_{12}\)
component has energy greater than \(13/20\).
\end{proposition}
\begin{proof}
The following complete native-cosine and paid-helper argument transcribes
\src{A12-new}, at its general real geometric scope.
Use \(N=13\), the exact window \(5\leq i\leq8\),
loss \(.65\), and \(80/13\leq D_i\leq84/13\).
The same complete two-path finite-white classification
applies at these actual ranks: each path has rank
at least \(4\), a joined extension has rank \(12>8\),
and its only connected finite completion is the
ordinary path. The zero-own repairs are unchanged:
internal \(3/4\) and two-white roles have affine
principals; endpoint \(3\) or next-to-end \(3\)
has time square below its positive foreign cost,
forcing an orthogonal closure; endpoint \(4\)
has \(Q<1/5\), hence negative time square.
Every own contact is positive.

The strength ratio now satisfies, for the entire
window,
\[
 T(q)=20/13-\frac{q-.65}{q(2q-.65)}
 <T(3)=5809/4173<7/5.
\]
Near simple \(3\), \(D/E\) positions and endpoint
\(4\) have the same larger ordinary residual ratios
as before. Even the near-four exception has minimum
\((1/\sqrt2+2/5)/\sqrt{3/5}>7/5\), so it is strong.
Far white \(4\)/own \(3\) retains affine \(F_4\).
Two-white-\(3\) has ratio at least
\(q/2+1/q>7/5\) at positive own; zero own is affine.
Positive foreign cost strengthens these tests;
the exceptional foreign far/helper ratios remain
\(>761/500\) or have square \(>121/35\).

For isolated owns put
\(a_0=(\sqrt D+\sqrt{D-1})^{-1}<.212\).
The actual cleared \(0<u<1\) pressure has loss
\(<(.212+1.023)/2=.6175\);
foreign \(u>1\) has loss \(<.212+.383=.595\).
Pure finite owns at least \(8\) have loss
\(<.212+\tan(\pi/8)<.627\).
The native-five \(H_4\) point has loss
\(<125/352+5/24=595/1056<.57\).
Native seven's \(A_2\) core has loss
\(<.264+\sqrt{13/120}<.601\),
using \(.9<\gamma<.902\) and
\(\gamma^{-2}>49/40\).
Own \(3,4,6\), parallel and \(u=1\) have their
actual affine/null exclusions.
All spectators are cleared positively and every
unknown wall is checked as in the shared proofs.
Thus the complete weak list is pure
\[
 v=\gamma\omega_i+\omega_e/2,\quad e=1,12,\quad Q=0.
\]
The identity \(\sigma^2/\delta\geq4\gamma-2\)
and \(T<7/5\) force
\(\varphi/2\leq\gamma<.99<1\).
Only own \(5\) or finite own at least \(7\)
remains; dotted owns have been excluded by the
actual edge inequality.

We prove fresh nonadjacency, without importing a
larger-rank determinant. For adjacent owns \(A,B\),
the same-helper product satisfies
\[
 R\geq1-AB+(AB-|A-B|)/13>1-.99^2>0,
\]
contradicting acuteness. Both opposite-helper
AM--GM bounds are \(>-1/2\), so native gap forces
mutual \(R=0\).
For outward helpers, the retained parameters obey
\(\eta<1/(2\gamma^2)<25/32<4/5\), whereas
\[
 (D_i-.8)(D_{i+1}-.8)-C_{i,i+1}^2
 =(2/5)C_{i,i+1}-24/25+8/65
 \geq428/325>0.
\]
This makes zero mutual impossible by the actual
time/spatial determinant identity.

For reversed helpers put \(l=i,s=12-i\);
\(l+s=12,\ ls=35\) or \(36\).
The exact time and spatial formulas are
\[
 A_c^2=\frac{2l(s+1)A^2+2lA-7}{13},\quad
 B_c^2=\frac{2(l+1)sB^2+2sB-7}{13},
\]
\[
 S=\frac{2lsAB+(s+1)A+(l+1)B+1/2}{13},\qquad
 A_c^2B_c^2-S^2=P/13,
\]
\[
 P=4lsA^2B^2-(l+1)(s+1)(A^2+B^2)-2AB
 -(l+1)A-(s+1)B+15/4.
\]
Zero mutual requires \(P=0\).
The polynomial increases in either variable on
\([.8,17/16]\): its derivative bracket is at
least \(1.56ls-13\geq41.6\).
If neither own is \(5\), \(A,B>.9\), and the
minimum \(P(.9,.9)\) at \(ls=35\) is \(3.624>0\).
If both are \(5\), direct reduction gives
\(P=((ls-56)\varphi-13)/4<0\), also nonzero.

It remains \(A=\varphi/2\), \(B=\cos(\pi/m)\),
\(m\geq7\), including its reversed placement.
As a polynomial in \(B\), \(P\) has leading
coefficient \(a=ls\varphi-13\ne0\) and linear
coefficient \(b=-\varphi-s-1\ne0\).
If it vanished, \(B\) would have degree at most
\(2\) over \(\mathbb Q(\varphi)\), hence degree
at most \(4\) over \(\mathbb Q\).
The real cyclotomic degree formula gives
\(\phi(2m)\leq8\).
Its complete even list with \(2m\geq14\) is
\(14,16,18,20,24,30\).
Indeed prime divisors can only be \(2,3,5,7\);
the totient product bounds leave exactly these
even integers. Thus
\(m=7,8,9,10,12,15\).
For \(7,9\), the native cubic degree remains
\(3\) over the quadratic \(\mathbb Q(\varphi)\).
For \(8\), the quartic real cyclotomic field
has unique quadratic subfield \(\mathbb Q(\sqrt2)\):
its generator \(\cos(\pi/8)\) has a cyclic
order-four conjugation, whose square fixes
\(\sqrt2\).
For \(12\), the field is
\(\mathbb Q(\sqrt2,\sqrt3)\), with quadratic
subfields \(\sqrt2,\sqrt3,\sqrt6\).
Neither contains \(\sqrt5\); both native
degrees remain \(4\) after adjoining \(\varphi\).
For \(10\), \(B^2=(\varphi+2)/4\) has
nonsquare rational norm \(5/16\); its minimal
quadratic has zero linear coefficient, unlike \(P\).
For \(15\), the minimal quadratic is
\[
 4B^2-2(\varphi-1)B-(\varphi+1),
\]
whose discriminant squareclass \(3(2+\varphi)\)
has nonsquare rational norm \(45\).
Proportionality to \(P\) would require
\[
 2b+a(\varphi-1)=ls+11-2s-15\varphi=0,
\]
impossible. All native cases are excluded.
This is algebra of literal finite cosines,
not an arithmeticity restriction on the polytope.
Weak sites are consequently nonadjacent.

Use heights
\(H_5=H_8=a=1.830,\ H_6=H_7=b=1.891\),
linear ends and reflection.
Their strong-floor tests are
\(2.480^2<80/13,\ 2.541^2<84/13\).
Flatten a weak outer end and omit internal weak
knots by neighbor chords. The complete reference
bounds are
\[
\begin{array}{c|c|c}
\text{weak outer sites }5,8&
 \text{internal omissions}&E\geq\\ \hline
0&\text{at most one, loss } .061^2/4&
2690283/4000000\\
1&\text{at most one, same loss}&1901431/3000000\\
2&6,7\text{ both strong}&3575881/6000000.
\end{array}
\]
Zero outer weak sites already gives \(e>.65\).
For one outer weak site, orient it to \(5\).
The block lengths are \(6,5\), with prices
\(\beta_6\geq b/12,\ \beta_8\geq.1525\).
For two outer weak sites they are \(6,6\),
with \(\beta_6=\beta_7=b/12\).
Neighbor omissions only increase capacities;
no weak floor is charged.

Choose one actual weak outer helper.
Use the finite-five or finite-at-least-seven
actual core in Section~\ref{app:Aendpointcores};
their selected roots are still present at
\(N=13\), with all-wall signs and future
orientations unchanged.
The freshly computed times are
\[
 A_{\rm far}^2=(80\gamma^2+10\gamma-7)/13,\qquad
 A_{\rm near}^2=A_{\rm far}^2+6\gamma/13.
\]
Project its displacement onto the constant
end blocks. In the six/six reference,
\[
 \Pi U=(5\gamma/6)\omega_6+\omega_7/12
 \quad\text{far},\qquad
 \Pi U=(5\gamma/6+1/12)\omega_6
 \quad\text{near}.
\]
For six/five, replace the far \(\omega_7/12\)
by \(\omega_8/10\).
The residual \(Y=U-\Pi U\) has zero mean in
each block and \(\ip{w}{Y}=0\).
Fresh block variance gives
\[
 \norm{Y}^2=(2/3)\gamma^2+1/12\quad(5),\qquad
 4\gamma^4-(16/3)\gamma^2+29/12
                 \quad(\text{finite }\geq7),
\]
with near addition \(\gamma/3\).
A six/five far residual decreases \(1/60\).
Thus every case has \(\norm Y^2<17/12<8/5\).
Both reference pairings obey
\[
 \ip{w}{\Pi U}\leq b(5\gamma/6)+a/10.
\]
For native five, \(A_{\rm far}>2.027\),
radical \(<.111\), pairing \(<1.459425<1.46\),
so its deficit exceeds \(.456\).
For finite owns at least seven,
\(A_{\rm far}(.9)^2=334/65>2.266^2\);
the pairing is \(1.60125\), radical at most \(1/4\),
giving deficit \(>.41475\).
Its deficit increases because
\((A_{\rm far}')^2-80/13=585/(169A_{\rm far}^2)>0\).
Near times are larger and pairings smaller.
Uniformly \(\Delta>.4\).
In particular the far-native-five pairing is
bounded by \(1.46\), not the false value \(1.438\).

Pay \(\mu=.1\) for this one helper.
A far foot spends \(<.0825\) at \(\omega_6\)
and at most \(.01\) at the opposite strong
anchor; a near foot spends \(<.090834\)
at \(\omega_6\).
They fit \(b/12>.15758\) and \(.1525\)
separately before pressure is used.
The single vector \(W=w+.1Y\) gives
\[
 e>E+2(.1)(.4)-(.1)^2(8/5)=E+.064.
\]
Even the worst reference yields
\[
 e>3575881/6000000+8/125
  =3959881/6000000
  =13/20+59881/6000000.
\]
All weak sets, orientations and native labels
are exhausted, with one original-vector budget.
\end{proof}


\section{The actual long-prefix inventory and its point pressures}
\label{app:prefixinventory}
We supply the geometric input of Proposition~\ref{prop:prefix}.
Its exact domain is \(6\leq M\leq16\), \(B_r\) with
\(r\geq M+5\), or \(D_r\) with \(r\geq M+6\).
For \(5\leq i\leq M+1\), deletion of old \(i\) leaves a
left \(A_{i-1}\) of rank at least \(4\), and a right
\(B_s\), \(s\geq4\), or \(D_s\), \(s\geq5\).
All roots used at \(5\) and its neighbors are actual long
roots. No short-tail exclusion is imported.
Write \(D_i=2i,q_i=\sqrt{2i}\) and
\(h_i=\ip{x}{\omega_i}=\sqrt2S_i\).

\subsection{The strong-ratio test}
An edge is strong if
\[
 \frac{\cosh L_i}{q_i}>T(q_i),\qquad
 T(q)=\frac{1+(1-.61/q)^2}{q(1-(1-.61/q)^2)}
 =\frac{100}{61}-\frac{q-.61}{q(2q-.61)}.
\]
Its half-tangent then exceeds \(1-.61/q_i\).
In particular \(T(q)<100/61\).
For the small left layers, clearing positive denominators
gives \(T(q)<38/25\) on \([5/2,10/3]\), and
\(T(q)<8/5\) on \([3,4]\). The resulting quadratics
are concave and positive at both endpoints.

\subsection{Clearing all spectators from an isolated own}
For an isolated own \(\gamma\), retain all contacted spectator
blocks, with spatial sum \(v_S\) and full cost \(Q\). Put
\[
 c=A p-\gamma\omega_i-v_S,\quad
 u=(1-Q)/\gamma^2>0,\quad A^2=\gamma^2(D_i-u).
\]
If \(u<1\), the positive actual-normal combination
\[
 \xi=c+v_S+\frac{1-Q}{\gamma}n_i
\]
has products zero on \(c,S\), product
\(\gamma(u-1)<0\) on \(n_i\), square
\(-\gamma^2u(1-u)\), and time \(A>0\).
Inverse positivity expands \(v_S\) in actual spectator
normals with nonnegative coefficients. Every unknown
wall product is nonpositive. Its actual point gives
\[
 h_i>\sqrt{D_i-u}-\sqrt{u(1-u)}.
\]
At \(u=1\) it is a nonzero feasible future null.
For \(u>1\), the positive finite projection on
\(S+n_i+c\), with residual \(\gamma^2(u-1)\), gives
\[
 h_i>\sqrt{D_i-u}-\sqrt{u-1}.
\]
Own \(3,4,6\) retain an affine fork, the \(4/3/3\)
star, or the \(6/3\) path; a parallel own has null
\(c+n_i\). Thus \(\gamma\geq\varphi/2\).
If \(Q>0\), \(u\leq3/\varphi^2=6-3\varphi<7/6\).
For \(u>1\) the loss is \(<1/5+2/5<.61\).
For \(u<1\) it is at most
\[
 u/6+\sqrt{u(1-u)}
 \leq(1+\sqrt{37})/12<3/5.
\]
Every foreign isolated cap is strong.

At \(Q=0\), isolated own \(5\) has an actual
\(H_4\) correction on three consecutive old nodes,
with \(\delta_5=(7-3\sqrt5)/16<1/36\).
Concavity bounds its corrected spatial product by
\(\gamma h_i\); its loss is \(<25/76+5/24<.61\).
Own \(7\) has the negative \(A_2\) core, of loss
\(<400/1701+9/25<.61\), using \(.9<\gamma<10/11\).
Finite own at least \(8\) has singleton pressure
\(h_i>\sqrt{D_i-1}-\tan(\pi/m)\), of loss
\(<4/21+5/12<.61\).
For dotted own use \(c+n_i/\gamma\) and \(u=\gamma^{-2}\).
Each correction is positive in actual roots.
No projected mirror is declared native.

\subsection{The complete finite-white lists}
A left \(A_{i-1}\) permits one contact: endpoint \(3\),
endpoint \(4\), next-to-end \(3\), or an admissible
\(E_6,E_7,E_8\) position \(3\).
Their Schur residuals are
\[
 (i+1)/(2i),\quad1/i,\quad2/i,\quad1/4,\ 1/7,\ 1/16.
\]
For a simple contact at position \(j\) from the far end,
\(\sigma=\gamma+j/(2i)\); for endpoint \(4\),
\(\sigma=\gamma+j/(\sqrt2i)\).
The actual edge ratio is \(\sigma/\sqrt{\delta-Q}\).
The right list is
\[
\begin{array}{c|l}
B_s,\ s\geq4&\text{long-leaf }3\\
D_s,\ s\geq5&\text{long-leaf }3;
 \text{ fork }3\text{ only for }s=5,6,7.
\end{array}
\]
Two contacts to a retained component create a nonfinite cycle.
A white bridge joining left and right has rank \(r>8\);
classification forces finite type \(B\) or \(D\), with simple
endpoint attachments. A further branch at its internal new
node is not finite. Exceptional fork pieces are retained
with their full allowed foreign costs.

Zero own needs a separate test. Internal simple/doubled left
contacts, right long contacts and full bridges retain affine
forks, stars or cycles. Right \(D\)-fork contact retains
affine \(E_8\). For the separate far-\(3\) case,
\[
 A^2=Q-\tfrac12<1/(2i)\leq1/10.
\]
Any foreign contact has derivative square at least \(1/4>A^2\),
forcing an orthogonal closure; at \(Q=0\) time square is
negative. Zero own has therefore not been omitted.

For positive own, the complete direct-strength check uses
the displayed residuals. Near simple-\(3\) has
\(\delta\leq3/5,\sigma\geq\gamma+2/5\);
next-to-end has \(\delta=2/i,\sigma\geq\gamma+1/i\);
\(E\)-rows have \(\delta\leq1/4\).
For \(\gamma>101/125\), near simple-\(3\) exceeds
\(100/61\) from \(i=8\), and next-to-end from \(i=7\).
At \(i=5,6,7\), the smaller interval tests for \(T\)
give the same conclusion. Own \(3,4,6\) are affine
except near-\(4\)/white-\(3\), whose \(B\) cosine
\((i-1)+i/\sqrt2\) is already strong for \(i\geq5\).
In \(D\), own \(4\) without right contact is affine on
the actual right fork. White \(4\) is stronger, except
far white \(4\)/own \(3\), with affine \(3,4,3,3\) path.

Right long \(3\) has
\(\delta=1/2-Q,\sigma=\gamma+1/2\).
Own \(0,3\) is affine; the next contact has ratio
at least \(1+1/\sqrt2>100/61\).
For right fork \(3\) in \(D_s\), \(s=5,6,7\),
\[
 \delta=1-s/8-Q,\qquad \sigma=\gamma+1/4.
\]
Own \(0,3,4,6\) retains affine \(E_8\), a cycle,
the \(3,3,4,3\) path, or the \(6,3\) path.
The next own gives ratio
\((\cos(\pi/5)+1/4)/\sqrt{3/8}>1.72>100/61\).
All finite foreign extensions of the \(E_6,E_7,E_8\)
piece remain included through \(Q\).
A full bridge has residual at most \(1/(2i)\),
numerator at least \(\gamma+1/2+1/(2i)\);
positive own is strong, and zero own gives an affine cycle.

A far simple-\(3\) cap with \(Q>0\) has
\[
 \cosh L_i=\frac{2i\gamma+1}{\sqrt{i+1-2iQ}},\quad
 \cosh L_i/q_i>R_0(i)
 :=\frac{(809/500)i+1}{\sqrt{i(i+2)}}.
\]
Both \(R_0(i)\) and \(T(\sqrt{2i})\) increase;
the first derivative has sign \((309/500)i-1>0\).
The exact intervals cover \(5\leq i\leq17\):
\[
\begin{array}{c|c|c|c}
i&q_i\text{ upper}&R_0(i)\text{ lower}&T(q_i)\text{ upper}\\ \hline
5\text{--}8&4&1.536&1.530\\
9\text{--}12&5&1.563&1.550\\
13\text{--}29&7.62&1.577&1.577.
\end{array}
\]
The squared lower margins are \(.052740,.321913,.546001\).
The upper comparisons follow by multiplying the numerator/
denominator pairs
\((27.4921,18.0316),(44.2721,28.6395)\),
\((107.2045,68.003166)\).
Every foreign far cap is strong.

Consequently every weak cap has exactly
\begin{equation}\label{app:eq:pureprefix}
 v=\gamma\omega_i+\omega_1/2,\quad
 A^2=2i\gamma^2+2\gamma-\tfrac12,\quad Q=0,\quad
 101/125<\gamma<6/5.
\end{equation}
Its own is \(5\), finite at least \(7\), or dotted.
For the upper cutoff,
\(\sigma^2/\delta\geq4\gamma-2\), with difference
numerator \([2i(\gamma-1)-1]^2\).
At \(\gamma=6/5\) this is
\(14/5>(100/61)^2\).

\subsection{Nonadjacency from one actual mixed point}
For adjacent weak owns \(a,b\), direct weighted AM--GM gives
\[
 \ip{c_i}{c_{i+1}}\geq
 \tfrac12-ab+\tfrac14(a/b+b/a)
 \geq1-ab>-11/25>-1/2.
\]
The actual product must be zero. Clearing all old roots
except \(i,i+1\) gives the two-cap Schur matrix
\[
 H=I-t\mathbf1\mathbf1^T,\qquad t=(i-1)/(2i)>0.
\]
It is positive with positive inverse. Put
\(\lambda=H^{-1}(A_i,A_{i+1})^T>0\) and
\(\mathcal D=A_i\lambda_i+A_{i+1}\lambda_{i+1}\).
The actual positive projection has displacement
\(Z=W_i\omega_i+W_{i+1}\omega_{i+1}\), with
\[
 W_i=\lambda_i a+(\lambda_i+\lambda_{i+1})/(2i),
 \quad W_{i+1}=\lambda_{i+1}b,\quad
 Z^2=\mathcal D(1+\mathcal D).
\]
Since \(Q=0\), all spectator blocks remain unchanged
positive orthogonal blocks; positivity checks every
unselected original wall.
Write \(W=W_i+W_{i+1}\), \(A_+=A_i+A_{i+1}>4\).
Then
\[
 \lambda_i+\lambda_{i+1}=iA_+,\quad
 \mathcal D=[iA_+^2+(A_i-A_{i+1})^2]/2,\quad
 W>(4/5)iA_+.
\]
It follows that
\[
 \frac{1+\mathcal D}{W^2}
 <\frac{(i+1)/2+1/16}{(16/25)i^2}
 \leq49/256,
\]
the last expression decreasing from \(i=5\).
For \(E_c=W_i\sqrt{2i}+W_{i+1}\sqrt{2i+2}\),
\[
 E_c-\norm Z=
 \frac{2W_iW_{i+1}(\sqrt{2i(2i+2)}-2i)}
      {E_c+\norm Z}<W/(4\sqrt{10}).
\]
The actual point gives
\(\ip{x}{Z}\geq1+\mathcal D-\sqrt{1+\mathcal D}\),
and \(\norm Z<1+\mathcal D\).
The weighted mean loss is therefore less than
\(7/16+1/(4\sqrt{10})<.61\), contradicting weakness
at both sites. This proves nonadjacency before using
any component-energy contradiction.

\subsection{The four actual weak-five pressures}
At \(i=5\), put \(A^2=10\gamma^2+2\gamma-1/2\) and
\(H(A,\delta)=A/(1+\sqrt{\delta/(A^2+\delta)})\).
The actual point pressures \(\ip{x}{U}\geq T\) are
\[
\begin{array}{c|l|l}
\text{own}&U&T\\ \hline
5&\gamma(\omega_4+\omega_6)/2+\omega_3/6
 &H(A,2/3-\gamma^2)\\
7,8&\gamma(\omega_4+\omega_6)/2+\omega_2/4
 &A-\sqrt{\gamma^2-3/4}\\
\text{finite }\geq9&\gamma(\omega_4+\omega_6)/2+\omega_1/2
 &H(A,1-\gamma^2)\\
\text{dotted}&(\gamma-\gamma^{-1})\omega_5+
 (\omega_4+\omega_6)/(2\gamma)+\omega_1/2
 &A-\sqrt{1-\gamma^{-2}}.
\end{array}
\]
The selected old sets are \(\{5,1,2\},\{5,1\},\{5\}\),
and the dotted pair. In the first three cases clear
those roots with their positive inverse coefficients.
For the dotted pair use \(c+n_5/\gamma\).
The selected products are zero or nonpositive and
time is positive. All unknown actual-wall products
are nonpositive by positivity; \(Q=0\) leaves pure
component displacement. Finite points are normalized
as in \eqref{eq:corepressure}, not as spacelike corrections.

Concavity gives
\(\ip{x}{\omega_4+\omega_6}\leq2\ip{x}{\omega_5}\),
\(\ip{x}{\omega_2}\leq2\ip{x}{\omega_1}\),
\(\ip{x}{\omega_3}\leq3\ip{x}{\omega_1}\).
Weak \(5\) therefore implies
\[
 x_1>\sqrt2\{T-\gamma(\sqrt{10}-.61)\}>.5965.
\]
For own \(5,7,8\), the lower bounds exceed
\(.8,.7,.65\), using
\(.809<\cos(\pi/5)<.81\),
\(.9<\cos(\pi/7)<.902\),
\(.923<\cos(\pi/8)<.9239\).
For finite owns at least \(9\), the deficit increases
from \(\gamma=.939\):
\(A'>\sqrt{10},H_A>9/10,H_\delta<0\).
At that endpoint
\(A>3.1929,\sqrt{A^2+\delta}>3.211,\sqrt\delta<.344\);
thus \(H>2.88\), giving the lower bound \(>.65\).
For \(1<\gamma<6/5\), the dotted deficit decreases:
\(A'<13/4\), the radical derivative exceeds \(1\),
and \(\sqrt{10}-.61>5/2\).
Its endpoint lower bound is
\[
 \sqrt2\{\sqrt{163/10}-\sqrt{11}/6
                -(6/5)(\sqrt{10}-.61)\}>.5965.
\]
Positive square tests
\(\sqrt{163/10}>4.03732,\sqrt{11}/6<.552771\),
\(\sqrt{10}<3.162278,\sqrt2>1.41421\)
prove it. Together with nonadjacency this supplies
both initial branches of Proposition~\ref{prop:prefix}
without assuming strong \(5\).

\section{Fresh exact-\(18\) classical layers}
\label{app:exact18layers}
This appendix supplies the short layers used in the
fresh rank-\(18\) cover, at their actual ranks rather
than by shifting a larger whole-component theorem.
All prices and energies concern the original vector.

\subsection{A local reference payment}
The long-prefix inventory already reaches \(12\)
in \(B_{16}\), and \(11\) in \(D_{16}\).
Any actual strong endpoint \(K\) supplies the
interpolation proof, in either initial-\(5\) branch.
Its endpoint sum is at most
\(f_K=\sqrt K-.61/\sqrt2\).
Its individual terminal coefficient is \(>1/7\)
for \(K=8,\ldots,12\), and \(>7/50\) for \(K=13\).
The ordinary coefficient is at least
\(\sqrt K-\sqrt{K-1}\); omitting the preceding
knot increases it. The weak-\(5\) reference
ends with slope \(1/(\sqrt K+\sqrt6)>1/7\).
The exact conservative data are
\[
\begin{array}{c|rrrrrr}
K&8&9&10&11&12&13\\ \hline
q_K>&.7683275&.797425&.823735&.8472225&.8689525&.88864\\
f_K<&2.398&2.570&2.732&2.886&3.034&3.175.
\end{array}
\]
Here \(q_K\) is the certified effective reference
energy; the entries follow by adding the positive
rational increment floors and permitted omission
losses in the prefix proof.

If an actual point/core gives \(\ip{x}{V}\geq G\),
\(V\) equals \(1\) on the first \(K\) coordinates,
and \(V^2\leq R\), pay \(\mu=1/7\) at that
individual terminal price. The residual \(V-P_K\)
has zero product with the reference and square
at most \(R-K\). One square gives
\begin{equation}\label{app:eq:localpayment}
 e_{\rm classical}>
 q_K+2(G-f_K)/7-(R-K)/49.
\end{equation}
The weak-\(5\) interpolation has higher energy
and no larger endpoint. No extra overlapping
first-coordinate gain is added to it.

\subsection{\(B_{16}\) with ordinary \(T_2\)}
At \(i=13\), deletion retains \(A_{12}+B_3\).
The right finite extensions are long \(3\) or
short \(3\), the latter giving \(F_4\).
Long \(3\) has
\(\sigma=\gamma+1/2,\delta=1/2-Q\);
own \(0,3\) is affine and larger owns are
directly strong.
Short \(3\) has
\(\sigma=\gamma+1/(2\sqrt2),\delta=1/4-Q\).
At zero own, the actual path
\(c,n_{16},n_{15},n_{14},n_{13}\) has affine
labels \(3,4,3,3\).
Positive own has ratio at least
\(1+1/\sqrt2\).
The \(F_4\) white block has no finite foreign
arm. No middle or doubled extra attachment
to \(B_3\) is finite.
The full \(B_{16}\) bridge has internal \(c\)
and is foreign-sealed; own zero is the cycle,
and positive owns have the direct strong ratio.
All left and isolated tests retain their
already proved actual matrices and full spectator.
Thus pure-helper coverage reaches \(13\).
The interpolation proof now gives
\(e(B_{16})>.8689525\), using the actual long
neighbor \(14\). Adding ordinary \(T_2>1/7\)
contradicts the norm. This is a fresh local
extension, not an out-of-domain application
of Proposition~\ref{prop:prefix}.

\subsection{\(D_{16}\) with ordinary \(T_2\)}
At \(i=12\), the right white \(D_4\) has
three simple-\(3\) leaf extensions: long leaf
\(13\), and the distinct signed forks \(15,16\).
Long \(13\) has
\(\sigma=\gamma+1/2,\delta=1/2-Q\);
its zero-own two-fork affine tree and own-\(3\)
triangle leave only directly strong cases.
A fork \(3\) has
\(\sigma=\gamma+1/4,\delta=1/2-Q\).
Own \(0,3,4,6\) gives respectively affine
\(E_8\) arms \(5,2,1\), the five-cycle,
the \(4/3/3\) star, or the \(6/3\) path;
parallel gives \(c+n_{12}\).
All remaining \(Q>0\) cases have ratio \(>2.1\);
pure finite owns at least \(7\) or dotted have
ratio \(>1.15\sqrt2>25/16\).
Only \(Q=0\), own \(5\)/fork-\(3\) remains.
Full bridges have residual \(1/24\);
zero own is a cycle and positive owns are
strong. Their internal new node cannot
take a finite foreign branch.

For own \(5\), put \(\gamma=\varphi/2\) and
\(\alpha^2=\varphi+5/8>56/25\).
The actual left \(H_4\) correction is
\(\gamma(3n_{12}/2+n_{11}+n_{10}/2)\),
with residual \(\delta_5=(5-3\varphi)/8\).
The finite future point is the positive
projection, not this spacelike correction.
Its pure scaled displacement is flat
through \(9\), with
\[
 A^2=12\alpha^2-\tfrac12,\quad
 R<12,\quad
 G>\sqrt{12-25/112}-.091>3.34.
\]
The radical bound is
\(\sqrt{\delta_5}/\alpha<.091\).
It is feasible against both actual forks,
the entire ordinary \(T_2\) and every unknown
wall by positive actual-root coefficients.
Choose a strong \(K=8\) or \(9\).
Equation~\eqref{app:eq:localpayment}, plus \(1/7\),
gives totals greater than
\[
 .7683275+2(3.34-2.398)/7-4/49+1/7>1.098,
\]
or
\[
 .797425+2(3.34-2.570)/7-3/49+1/7>1.098.
\]
Thus helper coverage reaches \(12\) before
considering the next layer.

At \(i=13\), the right \(D_3=A_3\) list is
middle \(3\), end \(3\), end \(4\), end \(5\).
Middle \(3\) has the affine four-leaf star
at zero own, the triangle at own \(3\),
and direct strength otherwise.
End \(4\) at zero own has the actual null
\[
 \sqrt2c+2n_f+2n_{14}+n_{13}+n_h,
\]
where \(h\) is the other actual fork.
Positive owns are directly strong.
End \(5\) is the sealed finite \(H_4\)
white block, of ratio \(>2.99\) even at
zero own.
The only full bridges are middle-\(3\)/full
\(D_{16}\), and end-\(3\)/full \(A_{16}\),
with residuals \(1/26\) and \(1/8+1/26\).
Zero own gives a cycle; positive owns are
strong. At rank \(16>8\), their internal
node cannot take another finite branch.

For no-left end-\(3\),
\(\sigma=\gamma+1/4,\delta=5/8-Q\).
The zero-own, own \(3,4,6\) and parallel
principals are the same \(E_8\), four-cycle,
doubled star, \(6/3\) path and null.
Positive foreign cases are strong, and
\(\gamma\geq1\) has ratio
\(\geq\sqrt{5/2}>25/16\).
The actual strength threshold obeys
\(T(\sqrt{26})<T(6)<25/16\).
Only \(Q=0\), own \(5\) or finite own
at least \(7\), remains.
Own \(5\) has its left \(H_4\) point,
now flat through \(10\), with
\[
 R<13,\qquad
 G>\sqrt{13-125/448}-.091>3.47.
\]
For finite \(.9<\gamma<1\), use both
actual negative \(A_2\) cores, corrections
\(\gamma(4n_{13}/3+2n_{\rm side}/3)\)
for sides \(12,14\), and average their
spatial pressures. With
\(\alpha=\sqrt2(\gamma+1/4)\),
\(d=4\gamma^2/3-1>0\), direct products give
\[
 A^2=13\alpha^2-5/8,\quad
 Y^2=A^2+1-14\gamma^2/9<13\alpha^2,\quad
 \ip{x}{Y}\geq A-\sqrt d.
\]
The average is flat through \(11\).
Both original cores satisfy every wall,
including both forks and all spectators.
Uniformly
\(G>\sqrt{13-125/529}-1/3>3.23,\ R<13\).
Choose strong \(K=9\) or \(10\);
even this weaker \(G\), with ordinary
\(T_2\), gives totals \(>1.047\) in
\eqref{app:eq:localpayment}.
All extra \(D_3\) roles are excluded.
Coverage reaches \(13\), so interpolation
gives \(e(D_{16})>.8689525\), again
contradicting its ordinary \(T_2\).
No typed defect is charged as ordinary.

\subsection{\(B_{17}\) with a nonshort \(A_1\)}
The generic inventory reaches \(13\).
At \(14\), the preceding complete \(B_3\)
layer reaches \(14\), with the actual
old long and short roots retained.
Choose strong \(K=12\) or \(13\).
At the right \(B_2\) layer \(15\), the
extra short \(3\) has its affine \(F_4\)
path at zero own; own \(4,6\) gives
a doubled star or \(6/3\) path;
parallel gives a feasible null.
Every own \(\gamma\geq\cos(\pi/5)\)
is strong. Positive foreign cost makes
own \(3\) strong as well.
Bridges are internal and foreign-sealed.
The sole extra weak case is therefore
\[
 Q=0,\qquad c=A p-(\omega_{15}+\omega_{17})/2.
\]
Releasing old \(16\) and retaining this
actual \(c\) gives finite \(A_{17}\):
old \(A_{15}\), \(c\), and the true
short old \(17\) form its chain.
Positive inversion gives its actual
future endpoint, feasible against released
\(16\) and every other original wall.
Its exact cosine and prefix consequence are
\[
 C_{16}=(16+15\sqrt2)/\sqrt{18}>8.7,\quad
 S_{16}^2>16(77/97)>(3.56)^2.
\]
Append \(S_{16}>3.56\) to the strong
\(K=12/13\) reference, skipping \(14,15\).
New slopes \(<.53/4,.39/3\) fit the
individual capacities \(>1/7,>7/50\).
The weak-five endpoint differs from
\(f_K\) by less than \(.00001\), so the
same strict junction bounds apply.
The two totals, including nonshort \(A_1>.1\),
exceed
\[
 .8689525+.526^2/4+.1>1.038,\qquad
 .88864+.385^2/3+.1>1.038.
\]
Thus this case is impossible.
Pure-helper coverage reaches \(15\);
interpolation gives \(e(B_{17})>.90715\),
whose addition to \(.1\) again contradicts \(1\).
Actual \(I\) has not been assigned a
nonshort debit, so \(B_{17}+I\) remains
for its separate consumer.


\section{Remaining transcription and review work}
\label{app:obligations}
The frozen source batch contains reciprocal shadows, ordinary contacts
and actual defects, the uniform \(B/D\) floor, shared actual
\(A\)-chain feet, the \(A_7\) base and uniform \(A_{\geq8},A_{\geq13}\)
costs, the combined \(N=17,18,19\) windows and paid references,
the complete fresh \(A_{12}>.65\) and \(A_{15}>.826\) refinements,
the generic-prefix inventory, the fresh exact-\(18\) classical layers,
and the complete earlier local \(A_{18}\) body used by the final
head proof. These are draft transcriptions and consolidations pending
source-bound mathematical review, not new accepted bounds or
formalization claims.

The manuscript remains incomplete while the following internal proof
bodies still have project Markdown homes:
\begin{enumerate}
\item The other exact-\(18\) row exclusions and whole \(A_{18}\).
\item All higher-dimensional covers and consumers through the
 published ceiling, with the general-only section DAG.
\item The complete whole-\(B_{18}\) terminal analysis, preserving
 the asymmetric individual-price correction.
\item Whole-\(D_{18}\) head isolation, the fresh \(.58\) inventory,
 raw-reference increment/omission tables and isolated-\(5\)
 pressure, strong \(16\), finite heads, the complete fork
 inventory and the earlier every-head pair exclusions.
\end{enumerate}
The whole-\(A_{18}\) reference and local isolated-endpoint gains
have a separate unintegrated working appendix fragment; it is
not part of this frozen batch or an additional manuscript.
The new final dotted-head estimates, true fork heights,
local adjacent-\(A_{18}\) argument and entire near/far contradiction
are written in the main text. Remaining bodies will be grouped by
technique, preserving every reviewed source version and excluding
unused historical numerical increments.

\section{External mathematical inputs}
\label{app:external}
The external inputs are spherical and affine Coxeter
classification, acute-angled compact-polytope simplicity
and wall acuteness, the Lann\'er order bound of five, and the published compact dimension
bound \(29\) for the tail beyond \(29\) \cite{Vin84,Vin85}.
The actual-wall sign and positive projection conventions are
those of \cite[p.~42, \S2.1]{Vin85}; all project-specific
component and terminal statements are proved internally.
All project-specific component, helper, forest and terminal
assertions are internal mathematics to be proved.
The elementary real-cyclotomic degree formula
\([\mathbb Q(\cos(\pi/m)):\mathbb Q]=\phi(2m)/2\) for \(m\geq3\)
is a standard algebraic input in the literal native-cosine analysis; it
imposes no bound on a polytope ground field.
Arithmeticity, field restrictions, integral dotted products,
arbitrary Coxeter-face inheritance, and abstract Gram
completions being realizable are not premises.
\begin{thebibliography}{9}
\bibitem{Vin84}
E.~B. Vinberg,
\emph{The absence of crystallographic groups of reflections in
Lobachevsky spaces of large dimension},
Trudy Moskov. Mat. Obshch. \textbf{47} (1984), 68--102;
English translation, Trans. Moscow Math. Soc. (1985), 75--112.
\href{https://www.mathnet.ru/eng/mmo444}{Primary bibliographic record}.

\bibitem{Vin85}
E.~B. Vinberg,
\emph{Hyperbolic reflection groups},
Russian Math. Surveys \textbf{40}:1 (1985), 31--75,
\href{https://doi.org/10.1070/RM1985v040n01ABEH003527}
{doi:10.1070/RM1985v040n01ABEH003527}.
\end{thebibliography}
\end{document}
